% Reformuler les idées essentielles d’un texte en vue d’un résumé — Français 1re Tech — Cours signé OnBuch+
% Programme officiel MINESEC. Compiler avec : tectonic N06-reformuler-idees-essentielles-texte-vue-resume.tex
\newcommand{\DOCMATIERE}{Français}
\newcommand{\DOCNIVEAU}{1re Tech}
\newcommand{\DOCMODULE}{Français – classe de Première technique}
\newcommand{\DOCLECON}{N06}
\newcommand{\DOCTITRE}{Reformuler les idées essentielles d’un texte en vue d’un résumé}
% OnBuch+ — préambule « cours signés » ENRICHI (compilé avec tectonic / XeLaTeX)
% Système visuel « Pop » d'OnBuch+ : fond crème (jamais blanc), bordures encre
% épaisses, ombres « dures » décalées, titres Archivo Black en capitales.
% Version MATHÉMATIQUES : graphes (pgfplots/tikz), boîtes pédagogiques
% (définition, propriété, méthode, attention, le savais-tu, exercice résolu,
% exercice, TP, bilan), page de garde et signature OnBuch+ ancrée en bas.
%
% Macros attendues dans main.tex AVANT \input{preamble} :
%   \DOCMATIERE  \DOCNIVEAU  \DOCMODULE  \DOCLECON (numéro)  \DOCTITRE
\documentclass[11pt]{article}

\usepackage{fontspec}
\usepackage[french]{babel}
\usepackage[a4paper,margin=2cm,top=3.4cm,bottom=2.8cm,headheight=1.4cm,footskip=1.3cm]{geometry}
\usepackage{amsmath,amssymb,amsfonts}
\usepackage{mathtools} % \xmapsto, \coloneqq... souvent utilisés par les modèles
\usepackage{cancel} % simplifications barrées (bilans de forces)
% \abs{z} (module, valeur absolue) et \norm{u} (norme), utilisés par les modèles
\providecommand{\abs}{}\let\abs\relax\DeclarePairedDelimiter{\abs}{\lvert}{\rvert}
\providecommand{\norm}{}\let\norm\relax\DeclarePairedDelimiter{\norm}{\lVert}{\rVert}
\usepackage[table]{xcolor} % [table] : fournit aussi \rowcolors (alternance de lignes)
\usepackage{pagecolor}
\usepackage{tikz}
\usetikzlibrary{arrows.meta,positioning,calc,shapes.geometric,shapes.misc,shapes.arrows,decorations.pathmorphing,decorations.pathreplacing,decorations.markings,patterns,shadows,mindmap,trees,fit,backgrounds,matrix,intersections,angles,quotes}
\usepackage{pgfplots}
\usepackage[version=4]{mhchem} % équations nucléaires \ce{^{235}_{92}U}
\usepackage[european,siunitx]{circuitikz} % schémas électriques
\pgfplotsset{compat=1.18}
\usepgfplotslibrary{fillbetween,groupplots} % aires entre courbes (calcul intégral)
\usepackage{enumitem}
\usepackage{fancyhdr}
% [most] charge toutes les bibliothèques tcolorbox (listings, minted, raster,
% documentation...) et gonfle inutilement la mémoire TeX sur un document long
% (~300 boîtes) : on ne charge que ce qui est réellement utilisé.
\usepackage[skins,breakable]{tcolorbox}
\usepackage{graphicx}
\usepackage{colortbl}
\usepackage{booktabs}
\usepackage{tabularx}
\usepackage{multirow}
\usepackage{diagbox} % case d'en-tête diagonale des tableaux à double entrée
% Colonnes extensibles centrée / gauche / droite (C, L, R), souvent utilisées par les modèles
\newcolumntype{C}{>{\centering\arraybackslash}X}
\newcolumntype{L}{>{\raggedright\arraybackslash}X}
\newcolumntype{R}{>{\raggedleft\arraybackslash}X}
\usepackage{array}
\usepackage{multicol}
\usepackage{siunitx}
% parse-numbers=false : accepte aussi \SI{2,5 \times 10^{-3}}{...}
\sisetup{locale=FR,output-decimal-marker={,},group-minimum-digits=4,parse-numbers=false}
\DeclareSIUnit\ohm{\ensuremath{\symup{\Omega}}} % U+2126 absent de la police
\DeclareSIUnit\division{div} % réglages d'oscilloscope (ms/div, V/div)
\DeclareSIUnit\tr{tr} % tours (vitesse de rotation, ex. \SI{1500}{\tr\per\minute})
\DeclareSIUnit\molar{M} % concentration molaire (ex. \SI{3}{\molar}, \SI{1}{\micro\molar})
\DeclareSIUnit\an{an} % année (le modèle confond parfois avec \year, un compteur TeX) — ex. \SI{5}{\kilo\meter\per\an}
\DeclareSIUnit\FCFA{FCFA} % franc CFA (ex. \SI{500}{\FCFA\per\kilo\gram})
\DeclareSIUnit\degreeBrix{\textdegree Brix} % teneur en sucre d'un sirop/jus (réfractométrie)
\DeclareSIUnit\month{mois} % le modèle utilise parfois \month, un compteur TeX, comme unité de durée
\DeclareSIUnit\microliter{\micro\liter} % le modèle écrit parfois \microliter en un seul token
\DeclareSIUnit\million{millions} % dénombrement (ex. \SI{15}{\million\per\mL} spermatozoïdes)
\usepackage{qrcode}
% \degree : commande fréquemment utilisée par les modèles pour un angle ou une
% température en dehors de \SI{}{} (non fournie par les paquets chargés).
\providecommand{\degree}{^{\circ}}
% \qed (fin de démonstration), utilisé par les modèles sans amsthm
\providecommand{\qed}{\hfill$\square$}
% \barre : double barre des valeurs interdites dans un tableau de variations (tkz-tab)
\providecommand{\barre}{\ensuremath{\|}}
% environnement proof (amsthm non chargé) utilisé par les modèles
\ifdefined\proof\else\newenvironment{proof}[1][Démonstration]{\par\noindent\textit{#1.}\ }{\qed\par}\fi
\usepackage{lastpage}
\usepackage{hyperref}
\hypersetup{hidelinks}

% ============================================================
% Palette « Pop » OnBuch+ — mêmes hex que la classe P (plus_ui.dart)
% ============================================================
\definecolor{poporange}{HTML}{F59321}
\definecolor{popdark}{HTML}{E07A0C}
\definecolor{popcream}{HTML}{FFF6E8}
\definecolor{popink}{HTML}{1C1714}
\definecolor{poppurple}{HTML}{7A5AE0}
\definecolor{popgreen}{HTML}{2E9E6A}
\definecolor{poppink}{HTML}{E0457B}
\definecolor{popblue}{HTML}{2F7DE1}
\definecolor{popgold}{HTML}{C98A12}
\definecolor{popmuted}{HTML}{8A7F76}
\definecolor{popline}{HTML}{EADFD2}
\definecolor{popgray}{HTML}{8A7F76}
\colorlet{popgrey}{popgray}
% Alias tolérants (le modèle nomme parfois les couleurs différemment)
\colorlet{popwhite}{white}
\colorlet{popblack}{popink}
\colorlet{popred}{poppink}
\colorlet{popyellow}{popgold}
% Teintes claires (fonds de tableaux/encadrés) — le modèle nomme parfois
% les couleurs avec un suffixe L (ex. popblueL) sans les avoir définies.
\colorlet{poporangeL}{poporange!20}
\colorlet{popdarkL}{popdark!20}
\colorlet{poppurpleL}{poppurple!20}
\colorlet{popgreenL}{popgreen!20}
\colorlet{poppinkL}{poppink!20}
\colorlet{popblueL}{popblue!20}
\colorlet{popgoldL}{popgold!20}
\colorlet{popredL}{popred!20}
\colorlet{popgrayL}{popgray!20}
% Teintes claires pour fonds de tableaux / zones
\colorlet{popblueL}{popblue!10!white}
\colorlet{poporangeL}{poporange!14!white}
\colorlet{popgreenL}{popgreen!12!white}
\colorlet{poppurpleL}{poppurple!12!white}
\colorlet{poppinkL}{poppink!12!white}
\colorlet{popgoldL}{popgold!14!white}

\pagecolor{popcream}

% ============================================================
% Typographie
% ============================================================
\defaultfontfeatures{Ligatures=TeX}
\setmainfont{PlusJakartaSans-Medium.ttf}[Path=../../../fonts/,BoldFont=PlusJakartaSans-Bold.ttf]
% Maths en sans-serif (Fira Math) avec chiffres et lettres droites de Plus
% Jakarta, pour que les formules se fondent dans le texte.
\usepackage[math-style=ISO,bold-style=ISO]{unicode-math}
\setmathfont{FiraMath-Regular.otf}[Path=../../../fonts/]
\setmathfont{PlusJakartaSans-Medium.ttf}[Path=../../../fonts/,range={up/num,up/{Latin,latin},\mathup}]
\usepackage{icomma} % virgule décimale sans espace en mode math
% Caractères mathématiques Unicode absents des polices de texte (Plus Jakarta,
% Archivo Black) : rendus par leur équivalent mathématique, en texte comme en
% formule — sinon ils s'affichent en carré vide (ex. « Arithmétique dans ℤ »).
\usepackage{newunicodechar}
\newunicodechar{ℝ}{\ensuremath{\mathbb{R}}}
\newunicodechar{ℤ}{\ensuremath{\mathbb{Z}}}
\newunicodechar{ℂ}{\ensuremath{\mathbb{C}}}
\newunicodechar{ℕ}{\ensuremath{\mathbb{N}}}
\newunicodechar{ℚ}{\ensuremath{\mathbb{Q}}}
\newunicodechar{𝔻}{\ensuremath{\mathbb{D}}}
\newunicodechar{α}{\ensuremath{\alpha}}
\newunicodechar{β}{\ensuremath{\beta}}
\newunicodechar{γ}{\ensuremath{\gamma}}
\newunicodechar{δ}{\ensuremath{\delta}}
\newunicodechar{ε}{\ensuremath{\varepsilon}}
\newunicodechar{θ}{\ensuremath{\theta}}
\newunicodechar{λ}{\ensuremath{\lambda}}
\newunicodechar{μ}{\ensuremath{\mu}}
\newunicodechar{σ}{\ensuremath{\sigma}}
\newunicodechar{τ}{\ensuremath{\tau}}
\newunicodechar{φ}{\ensuremath{\varphi}}
\newunicodechar{ψ}{\ensuremath{\psi}}
\newunicodechar{ω}{\ensuremath{\omega}}
\newunicodechar{Σ}{\ensuremath{\Sigma}}
% majuscules produites par \MakeUppercase dans les titres (α-aminés -> Α)
\newunicodechar{Α}{\ensuremath{\alpha}}
\newunicodechar{Β}{\ensuremath{\beta}}
\newunicodechar{Γ}{\ensuremath{\Gamma}}
\newunicodechar{Δ}{\ensuremath{\Delta}}
\newunicodechar{Ε}{\ensuremath{\varepsilon}}
\newunicodechar{Θ}{\ensuremath{\Theta}}
\newunicodechar{Λ}{\ensuremath{\Lambda}}
\newunicodechar{Μ}{\ensuremath{\mu}}
\newunicodechar{Τ}{\ensuremath{\tau}}
\newunicodechar{Φ}{\ensuremath{\Phi}}
\newunicodechar{Ψ}{\ensuremath{\Psi}}
\newunicodechar{Ω}{\ensuremath{\Omega}}
\newunicodechar{↦}{\ensuremath{\mapsto}}
\newunicodechar{∈}{\ensuremath{\in}}
\newunicodechar{∉}{\ensuremath{\notin}}
\newunicodechar{∧}{\ensuremath{\wedge}}
\newunicodechar{⇔}{\ensuremath{\Leftrightarrow}}
\newunicodechar{⟺}{\ensuremath{\Longleftrightarrow}}
\newunicodechar{⇒}{\ensuremath{\Rightarrow}}
\newunicodechar{⟹}{\ensuremath{\Longrightarrow}}
\newunicodechar{∘}{\ensuremath{\circ}}
\newunicodechar{∪}{\ensuremath{\cup}}
\newunicodechar{∩}{\ensuremath{\cap}}
\newunicodechar{≡}{\ensuremath{\equiv}}
\newunicodechar{⊂}{\ensuremath{\subset}}
\newunicodechar{⊥}{\ensuremath{\perp}}
\newunicodechar{∃}{\ensuremath{\exists}}
\newunicodechar{∀}{\ensuremath{\forall}}
\newunicodechar{∛}{\ensuremath{\sqrt[3]{\,}}}
\newunicodechar{∜}{\ensuremath{\sqrt[4]{\,}}}
\newunicodechar{⃗}{\ensuremath{{}^{\to}}}
\newunicodechar{ⁿ}{\ifmmode^{n}\else\textsuperscript{\ensuremath{n}}\fi}
\newunicodechar{ˣ}{\ifmmode^{x}\else\textsuperscript{\ensuremath{x}}\fi}
\newunicodechar{ᵇ}{\ifmmode^{b}\else\textsuperscript{\ensuremath{b}}\fi}
\newunicodechar{ʳ}{\ifmmode^{r}\else\textsuperscript{\ensuremath{r}}\fi}
\newunicodechar{ᵘ}{\ifmmode^{u}\else\textsuperscript{\ensuremath{u}}\fi}
\newunicodechar{ʸ}{\ifmmode^{y}\else\textsuperscript{\ensuremath{y}}\fi}
\newunicodechar{ᵃ}{\ifmmode^{a}\else\textsuperscript{\ensuremath{a}}\fi}
\newunicodechar{ᵉ}{\ifmmode^{e}\else\textsuperscript{\ensuremath{e}}\fi}
\newunicodechar{ᵏ}{\ifmmode^{k}\else\textsuperscript{\ensuremath{k}}\fi}
\newunicodechar{ᵖ}{\ifmmode^{p}\else\textsuperscript{\ensuremath{p}}\fi}
\newunicodechar{ᴺ}{\ifmmode^{N}\else\textsuperscript{\ensuremath{N}}\fi}
\newunicodechar{ᶜ}{\ifmmode^{c}\else\textsuperscript{\ensuremath{c}}\fi}
\newunicodechar{ʰ}{\ifmmode^{h}\else\textsuperscript{\ensuremath{h}}\fi}
\newunicodechar{ᵠ}{\ifmmode^{q}\else\textsuperscript{\ensuremath{q}}\fi}
\newunicodechar{ₙ}{\ifmmode_{n}\else\textsubscript{\ensuremath{n}}\fi}
\newunicodechar{ₐ}{\ifmmode_{a}\else\textsubscript{\ensuremath{a}}\fi}
\newunicodechar{ᵢ}{\ifmmode_{i}\else\textsubscript{\ensuremath{i}}\fi}
\newunicodechar{ᵣ}{\ifmmode_{r}\else\textsubscript{\ensuremath{r}}\fi}
\newunicodechar{ₖ}{\ifmmode_{k}\else\textsubscript{\ensuremath{k}}\fi}
\newunicodechar{ᵦ}{\ifmmode_{\beta}\else\textsubscript{\ensuremath{\beta}}\fi}
\newunicodechar{ₓ}{\ifmmode_{x}\else\textsubscript{\ensuremath{x}}\fi}
\newfontfamily\popdisplay{ArchivoBlack-Regular.ttf}[Path=../../../fonts/]
\newfontfamily\poplogo{SpaceGrotesk-Bold.ttf}[Path=../../../fonts/]
\newfontfamily\popsemi{PlusJakartaSans-SemiBold.ttf}[Path=../../../fonts/]
\newfontfamily\popxbold{PlusJakartaSans-ExtraBold.ttf}[Path=../../../fonts/]
\newfontfamily\popxboldls{PlusJakartaSans-ExtraBold.ttf}[Path=../../../fonts/,LetterSpace=3.0]

\setlist[enumerate]{leftmargin=1.7em,itemsep=5pt,topsep=4pt}
\setlist[itemize]{leftmargin=1.7em,itemsep=4pt,topsep=4pt,label={\color{poporange}\textbullet}}
\setlength{\parindent}{0pt}
\setlength{\parskip}{5pt}
\renewcommand{\arraystretch}{1.25}

% pgfplots : style Pop par défaut
\pgfplotsset{
  every axis/.append style={axis line style={popink,line width=1pt},tick style={popink},
    grid style={popline,line width=0.5pt},label style={font=\small},tick label style={font=\footnotesize}},
  popcurve/.style={poporange,line width=1.8pt,no markers,smooth},
}
% Même style disponible dans un \draw TikZ simple (hors pgfplots)
\tikzset{popcurve/.style={poporange,line width=1.8pt}}
\tikzset{popgrid/.style={popline,very thin}}
\colorlet{popcurve}{poporange} % le nom du style est parfois utilisé comme couleur

% ============================================================
% Pill / squiggle
% ============================================================
\newcommand{\poppill}[3]{%
  \tikz[baseline=-0.55ex]\node[draw=popink,line width=1.2pt,rounded corners=2.6pt,%
    fill=#2,inner xsep=6.5pt,inner ysep=3pt]{%
    \popxboldls\fontsize{7}{7}\selectfont\color{#3}\MakeUppercase{#1}};%
}
\newcommand{\popsquiggle}[1]{%
  \tikz[baseline=-0.4ex]{
    \draw[#1,line width=2.6pt,line cap=round]
      plot [smooth] coordinates {(0,0.09) (0.34,0.01) (0.68,0.21) (1.02,0.01) (1.36,0.10)};
  }%
}
% Mot-clé surligné (marqueur orange sous le texte)
\newcommand{\cle}[1]{{\popxbold\color{popdark}#1}}

% ============================================================
% En-tête / pied
% ============================================================
\pagestyle{fancy}
\fancyhf{}
\renewcommand{\headrulewidth}{0pt}
\renewcommand{\footrulewidth}{0pt}
\lhead{\raisebox{-0.15ex}{\poplogo\fontsize{13}{13}\selectfont\color{popink}OnBuch\color{poporange}+}}
\rhead{\poppill{\DOCMATIERE\ \textbullet\ \DOCNIVEAU\ \textbullet\ Leçon \DOCLECON}{white}{popink}}
\lfoot{\popxbold\fontsize{7.6}{7.6}\selectfont\color{popmuted}\MakeUppercase{Cours signé OnBuch+}\ \textbullet\ {\popsemi\DOCTITRE}}
\rfoot{\tikz[baseline=-0.6ex]\node[rounded corners=8.5pt,fill=popink,minimum height=17pt,minimum width=17pt,inner xsep=5pt,inner sep=0]{\popxbold\fontsize{8}{8}\selectfont\color{popcream}\thepage\,/\,\pageref*{LastPage}};}

% ============================================================
% Titres
% ============================================================
\newcommand{\chaptitre}[2]{%
  \begin{tcolorbox}[enhanced,colback=poporange,colframe=popink,boxrule=2.6pt,arc=11pt,
      drop shadow=popink,left=15pt,right=15pt,top=12pt,bottom=11pt,before skip=3pt,after skip=18pt]
    \poppill{#2}{popink}{popcream}\par\vspace{9pt}
    {\raggedright\hyphenpenalty=10000\exhyphenpenalty=10000\popdisplay\fontsize{22}{24}\selectfont\color{popcream}\MakeUppercase{#1}\par}
  \end{tcolorbox}
}
% Section numérotée : pastille orange avec le numéro + titre Archivo
\newcounter{popsec}
\newcommand{\coursec}[1]{%
  \stepcounter{popsec}\setcounter{popsub}{0}%
  \par\vspace{16pt}\noindent
  \begin{minipage}{\linewidth}
  \tikz[baseline=-0.7ex]\node[draw=popink,line width=1.6pt,fill=poporange,rounded corners=4pt,minimum size=22pt,inner sep=0]{\popdisplay\fontsize{12}{12}\selectfont\color{popcream}\thepopsec};%
  \hspace{8pt}{\popdisplay\fontsize{15}{17}\selectfont\color{popink}\MakeUppercase{#1}}\par
  \vspace{4pt}\noindent{\color{poporange}\rule{36pt}{3.6pt}}
  \end{minipage}\par\nopagebreak\vspace{8pt}
}
\newcounter{popsub}
\newcommand{\courssub}[1]{%
  \stepcounter{popsub}%
  \par\vspace{9pt}\noindent{\popxbold\fontsize{11.5}{13}\selectfont\color{popink}{\color{poporange}\thepopsec.\thepopsub}\hspace{6pt}#1}\par\nopagebreak\vspace{4pt}
}

% ============================================================
% Boîtes pédagogiques — même recette : fond blanc, bordure épaisse colorée,
% ombre dure encre, badge plein. Titre optionnel [#1].
% ============================================================
\tcbset{popbase/.style={breakable,enhanced,colback=white,boxrule=2.3pt,arc=9pt,
  shadow={3pt}{-3pt}{0pt}{popink},left=14pt,right=14pt,top=11pt,bottom=11pt,before skip=11pt,after skip=15pt,
  fontupper=\popsemi\fontsize{10.6}{14.5}\selectfont\color{popink},
  fontlower=\popsemi\fontsize{10.6}{14.5}\selectfont\color{popink}}}
% Coche dessinée (le glyphe ✓ n'existe pas dans les polices du document)
\newcommand{\popcheck}[1]{\tikz[baseline=-0.5ex]\draw[#1,line width=1.6pt,line cap=round,line join=round](-0.13,0.0)--(-0.04,-0.09)--(0.13,0.1);}
\providecommand{\coche}{\popcheck{popgreen}}
\providecommand{\resultat}[1]{\par\smallskip\noindent\fcolorbox{popgreen}{popgreenL}{\parbox{\dimexpr\linewidth-2\fboxsep-2\fboxrule\relax}{\textbf{Résultat.} #1}}\par\smallskip}
\newcommand{\popboxhead}[3]{\poppill{#1}{#2}{white}\ifx\relax#3\relax\else\hspace{6pt}{\popxbold\fontsize{10.6}{12}\selectfont\color{popink}#3}\fi\par\vspace{7pt}}

\newtcolorbox{aretenir}[1][]{popbase,colframe=popgreen,before upper={\popboxhead{À retenir}{popgreen}{#1}}}
\newtcolorbox{exemplebox}[1][]{popbase,colframe=poporange,before upper={\popboxhead{Exemple}{poporange}{#1}}}
\newtcolorbox{definition}[1][]{popbase,colframe=popblue,before upper={\popboxhead{Définition}{popblue}{#1}}}
\newtcolorbox{propriete}[1][]{popbase,colframe=poppurple,before upper={\popboxhead{Propriété}{poppurple}{#1}}}
\newtcolorbox{methode}[1][]{popbase,colframe=popgold,before upper={\popboxhead{Méthode}{popgold}{#1}}}
\newtcolorbox{attention}[1][]{popbase,colframe=poppink,before upper={\popboxhead{Attention}{poppink}{#1}}}
\newtcolorbox{experience}[1][]{popbase,colframe=popblue,colback=popblueL,before upper={\popboxhead{Expérience}{popblue}{#1}}}
% « Le savais-tu ? » : carte violette pleine (contexte, culture, Cameroun)
\newtcolorbox{savaistu}[1][]{popbase,colback=poppurple,colframe=popink,
  fontupper=\popsemi\fontsize{10.4}{14.2}\selectfont\color{white},
  before upper={\poppill{Le savais-tu\,?}{white}{poppurple}\ifx\relax#1\relax\else\hspace{6pt}{\popxbold\color{white}#1}\fi\par\vspace{7pt}}}
% Exercice résolu : énoncé en haut, \tcblower puis la solution sur fond crème
\newcounter{popexo}\newcounter{popexr}
\newtcolorbox{exoresolu}[1][]{popbase,colframe=poporange,colbacklower=popcream,
  segmentation style={popink,line width=1.2pt,dashed},
  before upper={\stepcounter{popexr}\popboxhead{Exercice résolu \thepopexr}{poporange}{#1}},
  before lower={\poppill{Solution}{popink}{popcream}\par\vspace{6pt}}}
% Exercice (non résolu en place) — la correction est dans la partie Corrigés
\newtcolorbox{exercice}[1][]{popbase,colframe=popink,
  before upper={\stepcounter{popexo}\popboxhead{Exercice \thepopexo}{popink}{#1}}}
% Corrigé
\newtcolorbox{corrige}[1][]{popbase,colframe=popgreen,colback=popgreenL,
  before upper={\popboxhead{Corrigé}{popgreen}{#1}}}
% Objectifs / prérequis / bilan
\newtcolorbox{objectifs}{popbase,colframe=popink,colback=poporangeL,
  before upper={\popboxhead{Objectifs de la leçon}{popink}{}}}
\newtcolorbox{prerequis}{popbase,colframe=popink,colback=popblueL,
  before upper={\popboxhead{Prérequis}{popblue}{}}}
\newtcolorbox{bilan}{popbase,colframe=popink,colback=white,boxrule=2.8pt,
  before upper={\popboxhead{Fiche bilan}{poporange}{L'essentiel à connaître pour l'examen}}}
% Figure encadrée avec légende
\newtcolorbox{popfigure}[1][]{enhanced,breakable=false,colback=white,colframe=popline,boxrule=1.4pt,arc=8pt,
  left=10pt,right=10pt,top=10pt,bottom=8pt,before skip=10pt,after skip=12pt,halign=center,
  lower separated=false,halign lower=center,
  fontlower=\popsemi\fontsize{9}{11.5}\selectfont\color{popmuted},
  #1}
\newcommand{\legende}[1]{\par\vspace{6pt}{\popsemi\fontsize{9}{11.5}\selectfont\color{popmuted}#1}}

% ============================================================
% Signature OnBuch+
% ============================================================
\newcommand{\poplogoinline}[1]{{\poplogo\fontsize{#1}{#1}\selectfont\color{popink}OnBuch\color{poporange}+}}
\newcommand{\signatureonbuch}{%
  \begin{tcolorbox}[enhanced,colback=popink,colframe=popink,boxrule=2.2pt,arc=10pt,
      drop shadow=poporange,left=14pt,right=14pt,top=10pt,bottom=10pt,before skip=0pt,after skip=0pt]
    \tikz[baseline=-0.7ex]\node[circle,fill=popgreen,draw=popcream,line width=1.3pt,minimum size=22pt,inner sep=0]{\popcheck{white}};%
    \hspace{9pt}\begin{minipage}[c]{0.62\linewidth}
      {\popxbold\fontsize{12}{13}\selectfont\color{popcream}Signé\ \textbullet\ {\poplogo OnBuch\color{poporange}+}}\par\vspace{2pt}
      {\raggedright\popsemi\fontsize{8}{10.5}\selectfont\color{popline}\MakeUppercase{Équipe pédagogique OnBuch+}\\\MakeUppercase{Conforme au programme officiel MINESEC}\par}
    \end{minipage}\hfill
    {\poplogo\fontsize{22}{22}\selectfont\color{popcream}OnBuch\color{poporange}+}
  \end{tcolorbox}%
}

% ============================================================
% Page de garde
% #1 titre · #2 matière · #3 niveau · #4 module · #5 n° leçon · #6 durée ·
% #7 description / objectifs (texte ou itemize)
% ============================================================
\newcommand{\pagegarde}[7]{%
  \thispagestyle{empty}
  \newgeometry{a4paper,margin=1.7cm,top=1.6cm,bottom=1.4cm}%
  \begin{tikzpicture}[remember picture,overlay]
    \fill[poporange!18] ([xshift=-3.2cm,yshift=-3.6cm]current page.north east) circle (4.6cm);
    \fill[poppurple!14] ([xshift=2.4cm,yshift=7.6cm]current page.south west) circle (3.8cm);
  \end{tikzpicture}%
  {\poplogo\fontsize{30}{30}\selectfont\color{popink}OnBuch\color{poporange}+}\hfill
  \raisebox{0.6ex}{\poppill{Cours signé \textbullet\ Premium}{popink}{popcream}}\par
  \vspace{1pt}\popsquiggle{poppurple}\par
  \vspace{8pt}
  \begin{tcolorbox}[enhanced,colback=poporange,colframe=popink,boxrule=2.8pt,arc=13pt,
      drop shadow=popink,left=18pt,right=18pt,top=12pt,bottom=13pt,after skip=10pt]
    \poppill{#2 \textbullet\ #3}{popink}{popcream}\hspace{5pt}\poppill{#4}{white}{popink}\par\vspace{8pt}
    {\popdisplay\fontsize{14}{14}\selectfont\color{popink}LEÇON #5}\par\vspace{4pt}
    {\raggedright\hyphenpenalty=10000\exhyphenpenalty=10000\popdisplay\fontsize{25}{27}\selectfont\color{popcream}\MakeUppercase{#1}\par}
    \vspace{8pt}
    {\popxbold\fontsize{9.5}{9.5}\selectfont\color{popink}\MakeUppercase{Durée conseillée : #6}}
  \end{tcolorbox}
  \begin{tcolorbox}[enhanced,colback=white,colframe=popink,boxrule=2.2pt,arc=10pt,
      drop shadow=popink,left=16pt,right=16pt,top=10pt,bottom=10pt,after skip=8pt]
    \poppill{Dans cette leçon}{poppurple}{white}\par\vspace{5pt}
    \popsemi\fontsize{9.3}{12.6}\selectfont\color{popink}#7
  \end{tcolorbox}
  \noindent\begin{minipage}[t]{0.487\linewidth}
    \begin{tcolorbox}[enhanced,colback=popgreen,colframe=popink,boxrule=2.2pt,arc=10pt,
        drop shadow=popink,left=11pt,right=11pt,top=10pt,bottom=10pt,height=3.1cm,valign=center]
      \tcbox[colback=white,colframe=popink,boxrule=1.3pt,arc=5pt,boxsep=0pt,nobeforeafter,
        left=3pt,right=3pt,top=3pt,bottom=3pt,valign=center]{\qrcode[height=1.9cm]{https://play.google.com/store/apps/details?id=com.luvvix.onbuch&hl=fr}}%
      \hspace{8pt}\begin{minipage}[c]{0.5\linewidth}
        {\popdisplay\fontsize{11}{12.5}\selectfont\color{white}\MakeUppercase{Télécharge OnBuch}}\par\vspace{4pt}
        {\raggedright\popsemi\fontsize{8.4}{11}\selectfont\color{white}Gratuit \textbullet\ Google Play\\Tuteur IA, annales, concours, résultats\par}
      \end{minipage}
    \end{tcolorbox}
  \end{minipage}\hfill
  \begin{minipage}[t]{0.487\linewidth}
    \begin{tcolorbox}[enhanced,colback=poppurple,colframe=popink,boxrule=2.2pt,arc=10pt,
        drop shadow=popink,left=11pt,right=11pt,top=10pt,bottom=10pt,height=3.1cm,valign=center]
      \tikz[baseline=-0.7ex]\node[circle,fill=white,draw=popink,line width=1.3pt,minimum size=1.6cm,inner sep=0]{\popdisplay\fontsize{24}{24}\selectfont\color{poppurple}+};%
      \hspace{8pt}\begin{minipage}[c]{0.6\linewidth}
        {\popdisplay\fontsize{11}{12.5}\selectfont\color{white}\MakeUppercase{Passe à OnBuch+}}\par\vspace{4pt}
        {\raggedright\popsemi\fontsize{8.4}{11}\selectfont\color{white}Tous les cours signés, Tuteur IA illimité, fiches PDF — onglet OnBuch+ de l'app\par}
      \end{minipage}
    \end{tcolorbox}
  \end{minipage}
  \vspace{8pt}
  \popcontenu
  \vfill
  \signatureonbuch
  \restoregeometry
  \newpage
}

% Rangée « Ce cours contient » (page de garde)
\newcommand{\poptile}[3]{% #1 couleur #2 gros texte #3 légende
  \begin{tcolorbox}[enhanced,colback=#1,colframe=popink,boxrule=2pt,arc=9pt,shadow={2.5pt}{-2.5pt}{0pt}{popink},
      left=6pt,right=6pt,top=9pt,bottom=9pt,halign=center,valign=center,height=1.75cm,width=\linewidth,nobeforeafter]
    {\popdisplay\fontsize{12.5}{13}\selectfont\color{white}\MakeUppercase{#2}}\par\vspace{3pt}
    {\centering\popsemi\fontsize{7.8}{9.5}\selectfont\color{white}#3\par}
  \end{tcolorbox}}
\newcommand{\popcontenu}{%
  \par\noindent{\popxboldls\fontsize{7.5}{7.5}\selectfont\color{popmuted}CE COURS SIGNÉ CONTIENT}\par\vspace{7pt}
  \noindent\begin{minipage}[t]{0.235\linewidth}\poptile{popblue}{Cours}{complet, illustré, pas à pas}\end{minipage}\hfill
  \begin{minipage}[t]{0.235\linewidth}\poptile{poporange}{Méthodes}{et exercices résolus}\end{minipage}\hfill
  \begin{minipage}[t]{0.235\linewidth}\poptile{poppink}{Exercices}{type examen + corrigés}\end{minipage}\hfill
  \begin{minipage}[t]{0.235\linewidth}\poptile{popgreen}{Fiche bilan}{l'essentiel à retenir}\end{minipage}\par}

% Sommaire
\newtcolorbox{sommaire}{enhanced,colback=white,colframe=popink,boxrule=2.2pt,arc=10pt,
  drop shadow=popink,left=18pt,right=18pt,top=13pt,bottom=13pt,before skip=6pt,after skip=20pt,
  before upper={\poppill{Sommaire}{poppurple}{white}\par\vspace{9pt}\popsemi\fontsize{10.6}{14.5}\selectfont\color{popink}}}

% Signature de fin de cours — poussée tout en bas de la dernière page
\newcommand{\findecours}{%
  \par\vfill\nopagebreak
  \begin{center}\popsquiggle{poporange}\end{center}\vspace{4pt}
  \signatureonbuch
}
\AtBeginDocument{\def\llbracket{\mathopen{[\mkern-3.5mu[}}\def\rrbracket{\mathclose{]\mkern-3.5mu]}}} % intervalles d'entiers (glyphe ⟦ absent de la police)
% Glyphes absents de Fira Math : redéfinis à partir de glyphes disponibles
\AtBeginDocument{%
  \def\setminus{\mathbin{\backslash}}\let\smallsetminus\setminus
  \def\vdots{\mathord{\vbox{\baselineskip3.2pt\lineskiplimit0pt\kern4pt\hbox{.}\hbox{.}\hbox{.}}}}%
  \def\ddots{\mathinner{\mkern1mu\raise6.5pt\hbox{.}\mkern2mu\raise3.6pt\hbox{.}\mkern2mu\raise0.7pt\hbox{.}\mkern1mu}}%
  \def\triangle{\mathord{\scalebox{0.9}{$\symup{\Delta}$}}}%
  \def\nabla{\mathord{\raisebox{\depth}{\scalebox{1}[-1]{$\symup{\Delta}$}}}}%
  \def\checkmark{\popcheck{popdark}}%
  \def\star{\mathbin{\ast}}\def\bigstar{\mathord{\ast}}%
  \def\diamond{\mathbin{\scalebox{0.6}{$\Diamond$}}}%
  \def\bigcup{\mathop{\mathchoice{\raisebox{-0.3em}{\scalebox{1.7}{$\cup$}}}{\scalebox{1.25}{$\cup$}}{\cup}{\cup}}\displaylimits}%
  \def\bigcap{\mathop{\mathchoice{\raisebox{-0.3em}{\scalebox{1.7}{$\cap$}}}{\scalebox{1.25}{$\cap$}}{\cap}{\cap}}\displaylimits}%
  \def\lesseqqgtr{\mathrel{\lessgtr}}\def\bigoplus{\mathop{\oplus}\displaylimits}%
}
\providecommand{\ssi}{si et seulement si} % abréviation fréquente
\AtBeginDocument{\providecommand{\wideparen}{\overparen}} % arc de cercle

%% FIGFIX-BEGIN v5 — correctifs globaux des figures (docs/onbuch-generation/tools/figfix)
\usepackage{adjustbox}\usepackage{etoolbox}\tcbuselibrary{raster}
% toute figure TikZ/pgfplots plus large que la ligne est réduite à la largeur disponible
\BeforeBeginEnvironment{tikzpicture}{\begin{adjustbox}{max width=\linewidth}}
\AfterEndEnvironment{tikzpicture}{\end{adjustbox}}
% légendes pgfplots lisibles par-dessus les courbes
\pgfplotsset{every axis/.append style={legend style={fill=white,fill opacity=0.92,text opacity=1,draw=black!15,inner sep=3pt}}}
% étiquettes d'arêtes (syntaxe edge["..."]) avec fond blanc
\tikzset{every edge quotes/.append style={fill=white,inner sep=1.2pt}}
%% FIGFIX-END
\begin{document}

\pagegarde{Reformuler les idées essentielles d’un texte en vue d’un résumé}{Français}{1re Tech}{Français – classe de Première technique}{N06}{14 séances (fiche de progression harmonisée)}{Cette leçon outille l'élève de Première technique pour maîtriser la contraction de texte : repérage de la structure argumentative, techniques de réduction et reformulation des idées essentielles. Elle s'appuie sur la lecture méthodique du Lion et la perle et sur les outils de la langue (registres, champs lexicaux, phrase complexe, paronymie) pour aboutir à la production de résumés conformes aux exigences du Baccalauréat technique.
\vspace{4pt}
\begin{multicols}{2}\small
\begin{itemize}
\item À la fin de cette leçon, je sais distinguer les registres de langue (soutenu, courant, familier) et les utiliser à bon escient dans une reformulation.
\item À la fin de cette leçon, je sais identifier la thèse, les arguments et les exemples dans un texte argumentatif.
\item À la fin de cette leçon, je sais dégager la structure argumentative d'un texte à contracter (plan dialectique, analytique ou thématique).
\item À la fin de cette leçon, je sais repérer et exploiter un champ lexical et un champ sémantique pour comprendre l'essentiel d'un texte.
\item À la fin de cette leçon, je sais appliquer les techniques de réduction (suppression, condensation, généralisation) sans déformer le sens.
\item À la fin de cette leçon, je sais analyser une phrase complexe pour en extraire l'idée principale et subordonner les idées secondaires.
\end{itemize}\end{multicols}}

\begin{sommaire}
\begin{multicols}{2}
\begin{enumerate}
\item Registres de langue et communication
\item La structure argumentative du texte
\item Champ lexical, champ sémantique et relations de sens
\item La phrase complexe au service de la reformulation
\item Les techniques de réduction du texte
\item Reformuler les idées essentielles et rédiger le résumé
\item Activité d'intégration
\item Méthodes et exercices résolus
\item Exercices
\item Corrigés des exercices
\item Fiche bilan
\end{enumerate}
\end{multicols}
\end{sommaire}

\begin{objectifs}
\begin{itemize}
\item À la fin de cette leçon, je sais distinguer les registres de langue (soutenu, courant, familier) et les utiliser à bon escient dans une reformulation.
\item À la fin de cette leçon, je sais identifier la thèse, les arguments et les exemples dans un texte argumentatif.
\item À la fin de cette leçon, je sais dégager la structure argumentative d'un texte à contracter (plan dialectique, analytique ou thématique).
\item À la fin de cette leçon, je sais repérer et exploiter un champ lexical et un champ sémantique pour comprendre l'essentiel d'un texte.
\item À la fin de cette leçon, je sais appliquer les techniques de réduction (suppression, condensation, généralisation) sans déformer le sens.
\item À la fin de cette leçon, je sais analyser une phrase complexe pour en extraire l'idée principale et subordonner les idées secondaires.
\item À la fin de cette leçon, je sais distinguer paronymes, hyperonymes et hyponymes pour éviter les contresens dans un résumé.
\item À la fin de cette leçon, je sais reformuler les idées essentielles d'un texte avec mes propres mots, en respectant la longueur imposée et la fidélité au texte.
\item À la fin de cette leçon, je sais rédiger et justifier ma démarche de contraction (compte-rendu) et présenter un exposé oral argumenté.
\end{itemize}
\end{objectifs}

\begin{prerequis}
\begin{itemize}
\item La phrase simple et ses expansions (classe de Seconde)
\item Les types de textes (narratif, descriptif, argumentatif)
\item Les connecteurs logiques de base (d'abord, ensuite, cependant, donc)
\item La lecture méthodique : repérage du thème et de la thèse (séances précédentes sur Le Lion et la perle)
\item Le résumé simple et le compte-rendu (classe de Seconde)
\end{itemize}
\end{prerequis}

\newpage

\coursec{Registres de langue et communication}

\courssub{Situation d'accroche : trois vies, une même compétence}

Imaginez trois scènes de la vie quotidienne camerounaise.

À Yaoundé, dans les studios de la \cle{CRTV}, un journaliste vient de recevoir un discours ministériel de vingt pages sur la crise du foncier à Douala. Il dispose de trois minutes d'antenne pour en rendre compte. Impossible de tout lire : il doit choisir, condenser, reformuler.

Quelques kilomètres plus loin, un étudiant prépare un exposé. Devant lui, un long article de \emph{Cameroon Tribune}. Sa consigne : le résumer en cent cinquante mots. S'il recopie des phrases entières, il dépasse la limite ; s'il écrit n'importe quoi, il trahit l'auteur.

À Bafoussam, une secrétaire de mairie doit transformer le compte rendu d'une réunion de deux heures en une note de synthèse d'une page pour le maire.

Dans les trois cas, une même compétence est exigée : \cle{reformuler} les idées essentielles d'un texte sans les trahir. Mais reformuler ne veut pas dire réécrire n'importe comment. Le journaliste de la CRTV ne parlera pas comme deux amis au marché Mokolo ; la note de la secrétaire ne ressemblera pas à un message WhatsApp. Chaque situation de communication impose un \cle{registre de langue} particulier.

\textbf{Problématique :} comment passer d'un texte long à un résumé fidèle, clair et personnel ? Pour y répondre, il faut d'abord comprendre ce qu'est un registre de langue, savoir le repérer dans un texte, et apprendre à choisir le bon registre pour reformuler. C'est l'objet de cette première section.

\courssub{Qu'est-ce qu'un registre de langue ?}

Commençons par une observation simple. Vous voulez annoncer la même nouvelle — votre réussite au Probatoire — à trois personnes différentes :

\begin{itemize}
\item à votre grand-mère, au village : « Grand-mère, j'ai eu mon examen ! » ;
\item à votre professeur : « Monsieur, j'ai le plaisir de vous annoncer que j'ai été admis au Probatoire. » ;
\item à votre meilleur ami : « Mec, j'ai cartonné à l'exam ! ».
\end{itemize}

L'idée est identique, mais les mots, les tournures et même le ton changent. C'est exactement ce qu'on appelle un registre de langue : on adapte sa manière de parler à la personne à qui l'on s'adresse, au lieu et aux circonstances.

\begin{definition}[Registre de langue]
Un \cle{registre de langue} est l'ensemble des marques linguistiques (choix du vocabulaire, tournures de phrases, niveau de syntaxe, ton) qu'un locuteur ou un rédacteur adopte en fonction de la \cle{situation de communication} : le destinataire, le lieu, le sujet, le degré de formalité. On distingue traditionnellement trois registres : le registre \cle{soutenu}, le registre \cle{courant} et le registre \cle{familier}.
\end{definition}

\begin{aretenir}[Les trois registres]
\begin{itemize}
\item Le registre \textbf{soutenu} : langue soignée, riche, parfois recherchée ; utilisée dans les discours officiels, les textes littéraires, les cérémonies.
\item Le registre \textbf{courant} : langue standard, correcte et naturelle ; celle de l'école, de la presse, des échanges professionnels. C'est le registre « neutre ».
\item Le registre \textbf{familier} : langue relâchée, spontanée, avec argot et abréviations ; réservée aux conversations entre proches.
\end{itemize}
\end{aretenir}

\courssub{Les marqueurs des trois registres}

Comment reconnaître concrètement un registre ? Il faut observer trois niveaux : le \cle{vocabulaire} (lexique), la \cle{syntaxe} (construction des phrases) et le \cle{ton}.

\begin{center}
\begin{tabularx}{\linewidth}{|l|X|X|X|}
\hline
\rowcolor{popblueL} \textbf{Critère} & \textbf{Soutenu} & \textbf{Courant} & \textbf{Familier} \\
\hline
Vocabulaire & riche, précis, parfois littéraire (« demeure », « allégresse », « s'entretenir ») & standard, précis mais simple (« maison », « joie », « discuter ») & argotique, abrégé (« baraque », « kif », « causer ») \\
\hline
Syntaxe & phrases longues, subordonnées, inversions (« À peine eut-il parlé que... ») & phrases de longueur moyenne, ordre sujet-verbe-complément & phrases courtes, ellipses (« Pas le temps ! »), ponctuation expressive \\
\hline
Ton & cérémonieux, distant, respectueux & neutre, poli & spontané, complice, parfois irrévérencieux \\
\hline
Situations typiques & discours officiel, hommage, texte littéraire & cours, article de presse, lettre administrative & marché, cour de récréation, réseaux sociaux \\
\hline
\end{tabularx}
\end{center}

\begin{exemplebox}[Une même réalité, trois registres]
Pour désigner le président de la République du Cameroun, on pourra entendre :
\begin{itemize}
\item registre \textbf{soutenu} : « Le locataire du bâtiment de briques » (périphrase élégante évoquant le palais présidentiel d'Etoudi) ;
\item registre \textbf{courant} : « le président », « le chef de l'État » ;
\item registre \textbf{familier} : « le vieux du palais ».
\end{itemize}
Le référent est le même ; seul le rapport au destinataire change. Dans un résumé, seul le registre courant est acceptable.
\end{exemplebox}

\begin{popfigure}
\begin{center}
\begin{tikzpicture}[font=\small]
% En-têtes
\fill[popblueL] (0,4.4) rectangle (4,5.4);
\fill[popgreenL] (4.3,4.4) rectangle (8.3,5.4);
\fill[poporangeL] (8.6,4.4) rectangle (12.6,5.4);
\node[font=\bfseries] at (2,4.9) {SOUTENU};
\node[font=\bfseries] at (6.3,4.9) {COURANT};
\node[font=\bfseries] at (10.6,4.9) {FAMILIER};
% Ligne 1
\node[draw, rounded corners, fill=white, text width=3.6cm, align=center, minimum height=1.1cm] at (2,3.5) {« Son Excellence a daigné présider la cérémonie. »};
\node[draw, rounded corners, fill=white, text width=3.6cm, align=center, minimum height=1.1cm] at (6.3,3.5) {« Le ministre a présidé la cérémonie. »};
\node[draw, rounded corners, fill=white, text width=3.6cm, align=center, minimum height=1.1cm] at (10.6,3.5) {« Le patron a géré la cérémonie. »};
% Ligne 2
\node[draw, rounded corners, fill=white, text width=3.6cm, align=center, minimum height=1.1cm] at (2,1.9) {« Les populations éprouvent d'insupportables difficultés. »};
\node[draw, rounded corners, fill=white, text width=3.6cm, align=center, minimum height=1.1cm] at (6.3,1.9) {« Les habitants ont de grosses difficultés. »};
\node[draw, rounded corners, fill=white, text width=3.6cm, align=center, minimum height=1.1cm] at (10.6,1.9) {« Les gens galèrent grave. »};
% Ligne 3
\node[draw, rounded corners, fill=white, text width=3.6cm, align=center, minimum height=1.1cm] at (2,0.3) {« Je vous saurais gré de bien vouloir... »};
\node[draw, rounded corners, fill=white, text width=3.6cm, align=center, minimum height=1.1cm] at (6.3,0.3) {« Je vous remercie de... »};
\node[draw, rounded corners, fill=white, text width=3.6cm, align=center, minimum height=1.1cm] at (10.6,0.3) {« Merci d'avance, t'es un frère ! »};
\end{tikzpicture}
\end{center}
\legende{Tableau comparatif des trois registres de langue : pour une même idée (ligne par ligne), le vocabulaire, la syntaxe et le ton varient selon la situation de communication.}
\end{popfigure}

\courssub{Le registre dépend de la situation de communication}

Pourquoi existe-t-il plusieurs registres ? Parce que toute parole s'inscrit dans une \cle{situation de communication} précise, que l'on peut schématiser ainsi :

\begin{popfigure}
\begin{center}
\begin{tikzpicture}[font=\small, >=Stealth, node distance=2.6cm]
\node[draw, rounded corners, fill=popblueL, text width=2.6cm, align=center, minimum height=1.2cm] (em) {\textbf{Émetteur}\\ celui qui parle ou écrit};
\node[draw, rounded corners, fill=popgreenL, text width=3.2cm, align=center, minimum height=1.2cm, right=of em] (msg) {\textbf{Message}\\ contenu + registre de langue};
\node[draw, rounded corners, fill=poporangeL, text width=2.6cm, align=center, minimum height=1.2cm, right=of msg] (rec) {\textbf{Récepteur}\\ celui qui écoute ou lit};
\draw[->, very thick, popblue] (em) -- (msg);
\draw[->, very thick, popblue] (msg) -- (rec);
\draw[->, dashed, popmuted] (rec.south) to[bend left=25] node[below, text width=3cm, align=center]{réaction, retour} (em.south);
\node[draw, rounded corners, fill=popcream, text width=9.5cm, align=center, below=1.6cm of msg] (ctx) {\textbf{Contexte} : lieu, moment, sujet, degré de formalité, relation entre les personnes};
\draw[->, thick, poppurple] (ctx) -- (msg.south);
\end{tikzpicture}
\end{center}
\legende{Schéma de la situation de communication : le contexte (lieu, destinataire, sujet) détermine le choix du registre dans le message.}
\end{popfigure}

Ce schéma montre une vérité essentielle : le registre n'est pas une décoration, c'est une \cle{adaptation}. Le même ministre parlera soutenu à la tribune d'une cérémonie officielle, courant dans un entretien à \emph{Cameroon Tribune}, et pourra se permettre une pointe de familier en visitant un marché de son village. Choisir le mauvais registre, c'est produire un \cle{effet de sens} non voulu : un discours familier à la radio nationale passerait pour un manque de sérieux ; un langage trop soutenu entre amis passerait pour de la prétention ou de la moquerie.

\begin{methode}[Identifier le registre d'un énoncé]
Pour déterminer le registre d'un passage, procédez en trois étapes :
\begin{enumerate}
\item \textbf{Observez le vocabulaire} : les mots sont-ils recherchés, littéraires (soutenu) ? Standards (courant) ? Argotiques, abrégés, empruntés (familier) ?
\item \textbf{Examinez la syntaxe} : phrases longues et construites avec subordonnées et inversions (soutenu) ? Phrases simples et claires (courant) ? Ellipses, exclamations, phrases nominales (familier) ?
\item \textbf{Repérez les tournures et le ton} : formules de politesse élaborées, périphrases (soutenu) ? Ton neutre et informatif (courant) ? Interpellations, familiarités, humour (familier) ?
Concluez en justifiant votre réponse par au moins deux indices précis relevés dans le texte.
\end{enumerate}
\end{methode}

\begin{exoresolu}[Repérer les registres dans trois extraits]
\textbf{Énoncé.} Indiquez le registre de chaque extrait et justifiez par deux indices.

\textbf{Extrait A} (discours officiel) : « Je voudrais, à l'entame de mon propos, rendre un vibrant hommage aux populations de la ville de Douala pour leur légendaire hospitalité. »

\textbf{Extrait B} (\emph{Cameroon Tribune}) : « Le Premier ministre a présidé hier à Yaoundé une réunion consacrée à la question du foncier urbain. Plusieurs mesures ont été annoncées. »

\textbf{Extrait C} (marché Mokolo) : « Eh ma sœur ! Ce prix-là est comment ? Tu vas me tuer avec tes tomates ! Laisse-moi ça pour deux cents, na ! »
\tcblower
\textbf{Solution détaillée.}
\begin{itemize}
\item \textbf{Extrait A : registre soutenu.} Indices : vocabulaire recherché (« à l'entame de mon propos », « vibrant hommage », « légendaire hospitalité ») ; tournure cérémonieuse et phrase construite. Situation : discours officiel.
\item \textbf{Extrait B : registre courant.} Indices : vocabulaire standard et précis (« a présidé », « réunion consacrée à », « mesures annoncées ») ; phrases claires, ton neutre et informatif. Situation : article de presse.
\item \textbf{Extrait C : registre familier.} Indices : interpellation (« Eh ma sœur ! »), tournure populaire (« ce prix-là est comment ? »), exclamations, particule empruntée aux langues locales (« na »). Situation : conversation de marché entre personnes proches.
\end{itemize}
\end{exoresolu}

\courssub{Effets de sens et adéquation à la situation}

Le choix d'un registre n'est jamais neutre : il produit un \cle{effet de sens}.

\begin{itemize}
\item Le registre \textbf{soutenu} crée de la \emph{distance} et du \emph{respect} : il solennise, il honore le destinataire. Un hommage funèbre prononcé en langage familier choquerait.
\item Le registre \textbf{courant} crée de la \emph{clarté} et de la \emph{neutralité} : il informe sans émotion excessive. C'est le registre de la transmission du savoir, donc de l'école et de l'examen.
\item Le registre \textbf{familier} crée de la \emph{proximité} et de la \emph{complicité} : il rapproche, il détend, mais il peut aussi manquer de respect s'il est employé hors de son contexte.
\end{itemize}

On parle d'\cle{adéquation registre/situation} : un énoncé est réussi lorsque son registre correspond à son destinataire, à son lieu et à son objet. Inversement, le \cle{mélange des registres} (passer du soutenu au familier dans la même phrase) produit un effet comique ou fautif.

\begin{attention}[Erreur fréquente : mélanger les registres dans un résumé]
Dans un résumé d'examen, il est interdit de laisser passer des tournures familières. Phrase fautive relevée en copie : « \emph{Le gars a dit que le foncier à Douala, c'est la galère.} » Deux fautes de registre : « le gars » (familier) et « c'est la galère » (familier). Correction au registre courant : « \emph{L'auteur affirme que la question du foncier à Douala constitue un problème majeur.} » Un seul mot familier suffit à dévaloriser tout un résumé, même si les idées sont correctes.
\end{attention}

\begin{savaistu}[Le camfranglais, une richesse... hors de l'examen]
Au Cameroun, de nombreux jeunes parlent le \cle{camfranglais}, une langue créative qui mélange français, anglais et langues locales : « \emph{Je suis venu hier soir, le gars avait déjà faya la fête} » (le garçon avait déjà gâché la fête). Ce parler témoigne d'une grande inventivité linguistique et joue un vrai rôle identitaire. Mais il relève du registre familier : il est \textbf{à proscrire absolument} dans un résumé, une dissertation ou toute copie d'examen, où seule la langue standard est attendue.
\end{savaistu}

\courssub{La règle d'or du résumé : reformuler au registre courant}

Venons-en à l'application qui nous intéresse dans cette unité : la contraction de texte. Résumer, c'est reformuler les idées essentielles d'un texte \emph{avec ses propres mots}. Mais avec quels mots ?

\begin{aretenir}[Règle d'or du résumé]
Un résumé se rédige toujours au registre \cle{courant} (ou soutenu si le texte de départ l'exige), \textbf{jamais au registre familier}. Reformuler, c'est \cle{neutraliser le style sans trahir le sens} : on conserve intégralement les idées de l'auteur, mais on les exprime dans une langue standard, claire et impersonnelle.
\end{aretenir}

Cette règle a deux conséquences pratiques :

\begin{enumerate}
\item Si le texte de départ contient des passages familiers (dialogues, citations de témoins, proverbes populaires), il faut les \emph{transposer} au registre courant dans le résumé.
\item Si le texte de départ est très soutenu ou très littéraire, on peut le \emph{simplifier} vers le registre courant, à condition de ne perdre aucune idée essentielle.
\end{enumerate}

\begin{exoresolu}[Transposition guidée : du familier au soutenu]
\textbf{Énoncé.} Voici le témoignage d'un habitant de Douala, cité dans un article sur la crise du foncier : « \emph{Moi, je te dis, ces histoires de terrain, ça me dépasse ! Le gars de la mairie m'a baladé pendant des mois, et à la fin, mon dossier a carrément disparu. C'est du vol, point !} »

Transposez ce témoignage au registre courant/soutenu pour l'intégrer dans un résumé.
\tcblower
\textbf{Solution détaillée.}

Étape 1 : repérer les marques du familier — « je te dis », « ça me dépasse », « le gars », « m'a baladé », « carrément », « c'est du vol, point ».

Étape 2 : identifier l'idée essentielle derrière chaque marque — l'habitant se déclare dépassé par les procédures ; un employé municipal lui a donné de faux espoirs pendant des mois ; son dossier a disparu ; il dénonce une injustice.

Étape 3 : reformuler au registre courant, à la troisième personne, en phrases construites.

\textbf{Transposition possible :} « Ce témoin se déclare dépassé par les procédures foncières. Il affirme qu'un employé de la mairie a fait durer sa démarche pendant plusieurs mois et que son dossier a finalement disparu, ce qu'il considère comme une injustice flagrante. »

On constate que toutes les idées sont conservées, que le ton est neutre et qu'aucune trace du familier ne subsiste : le sens est fidèle, le style est neutralisé.
\end{exoresolu}

\begin{methode}[Transposer un passage familier au registre courant]
\begin{enumerate}
\item \textbf{Repérez} les mots et tournures familiers (argot, ellipses, interpellations, exclamations).
\item \textbf{Dégagez} l'idée que chacun exprime réellement (demandez-vous : « que veut dire l'auteur ? »).
\item \textbf{Remplacez} chaque marque par un équivalent standard : « le gars » devient « l'homme », « l'employé », « l'auteur » ; « c'est la galère » devient « c'est très difficile ».
\item \textbf{Reconstruisez} des phrases complètes, à la troisième personne, avec un verbe de parole neutre : « affirme », « déclare », « explique », « dénonce ».
\item \textbf{Vérifiez} qu'aucune idée n'a été perdue et qu'aucun mot familier ne subsiste.
\end{enumerate}
\end{methode}

\begin{exercice}[À vous de jouer]
Relevez le registre des énoncés suivants, puis transposez au registre courant ceux qui sont familiers :
\begin{enumerate}
\item « Le chef a dit que l'école, c'est trop important, faut pas déconner avec ça. »
\item « Nul n'est censé ignorer que l'éducation constitue le socle du développement d'une nation. »
\item « L'auteur soutient que l'école joue un rôle essentiel dans le développement. »
\end{enumerate}
\end{exercice}

\begin{corrige}[Exercice : registres et transposition]
\begin{enumerate}
\item Registre familier (« le chef », « c'est trop important », « déconner »). Transposition : « Le responsable affirme que l'école revêt une importance capitale et qu'il ne faut pas la négliger. »
\item Registre soutenu (formule juridique « nul n'est censé ignorer », vocabulaire abstrait « le socle du développement »). Aucune transposition nécessaire ; on peut éventuellement simplifier vers le courant dans un résumé : « L'éducation est le fondement du développement d'une nation. »
\item Registre courant : c'est déjà la formulation idéale pour un résumé.
\end{enumerate}
\end{corrige}

\begin{aretenir}[Bilan de la section]
\begin{itemize}
\item Un \cle{registre de langue} est l'ensemble des marques linguistiques adaptées à une situation de communication ; on distingue les registres \cle{soutenu}, \cle{courant} et \cle{familier}.
\item On identifie un registre en observant le vocabulaire, la syntaxe et le ton.
\item Le choix du registre produit des effets de sens : distance et respect (soutenu), clarté et neutralité (courant), proximité et complicité (familier).
\item Règle d'or du résumé : reformuler au registre courant ou soutenu, \textbf{jamais au familier} ; neutraliser le style sans trahir le sens.
\end{itemize}
Cette maîtrise des registres est la première étape de la reformulation. La section suivante vous apprendra à repérer la structure argumentative d'un texte, c'est-à-dire le squelette d'idées que votre résumé devra restituer.
\end{aretenir}

\coursec{La structure argumentative du texte}

Dans la section précédente, nous avons vu que tout énoncé s'inscrit dans un \cle{registre de langue} (soutenu, courant, familier) choisi en fonction de la situation de communication. Mais un texte n'est pas seulement une affaire de mots : c'est aussi une affaire d'\emph{organisation des idées}. Lorsqu'un journaliste, un orateur ou un élève veut convaincre, il ne jette pas ses idées en désordre : il les \textbf{structure}. C'est précisément cette architecture que nous allons apprendre à démonter, car c'est la clé de la \cle{contraction de texte} : pour résumer un texte, il faut d'abord comprendre comment il est construit.

\courssub{Qu'est-ce qu'une structure argumentative ?}

\begin{definition}[Structure argumentative]
La \cle{structure argumentative} d'un texte est l'\textbf{organisation logique des idées} qui soutiennent une \cle{these} (l'idée défendue par l'auteur). Elle repose sur trois étages hiérarchisés :
\begin{itemize}
\item la \textbf{thèse} : l'idée principale que l'auteur veut faire admettre ;
\item les \textbf{arguments} : les raisons avancées pour justifier la thèse ;
\item les \textbf{exemples} : les faits, chiffres ou anecdotes qui illustrent concrètement chaque argument.
\end{itemize}
\end{definition}

Prenons une image simple. Un avocat au tribunal ne dit pas n'importe quoi dans n'importe quel ordre : il annonce d'abord sa conviction (« mon client est innocent »), puis il donne ses raisons (« il était ailleurs au moment des faits », « aucun témoin ne l'a vu »), puis il appuie chaque raison par une preuve (un billet de train, un témoignage). Un texte argumentatif fonctionne exactement de la même manière : c'est un \textbf{édifice} dont la thèse est le toit, les arguments les piliers et les exemples les briques.

\begin{exemplebox}[Les trois étages en miniature]
\emph{Thèse} : « Le Cameroun doit réduire ses importations de riz. »\\
\emph{Argument 1} : « Ces importations coûtent des milliards de francs CFA chaque année. » \emph{Exemple} : « En 2022, la facture dépassait 300 milliards. »\\
\emph{Argument 2} : « Le pays dispose de terres rizicoles sous-exploitées. » \emph{Exemple} : « Les plaines de Yagoua et de Ndom pourraient augmenter fortement leur rendement. »
\end{exemplebox}

\begin{aretenir}[La hiérarchie à retenir]
\textbf{Thèse} (1 idée) $\rightarrow$ \textbf{arguments} (2 à 4 raisons) $\rightarrow$ \textbf{exemples} (illustrations). Dans un résumé, on conserve la thèse et les arguments ; on supprime presque toujours les exemples.
\end{aretenir}

\courssub{Les trois grands plans argumentatifs}

Tous les textes argumentatifs ne sont pas bâtis sur le même plan. En classe de Première technique, il faut en connaître trois.

\begin{definition}[Plan analytique]
Le \cle{plan analytique} procède en trois temps : \textbf{constat} (on décrit un problème) $\rightarrow$ \textbf{causes} (on explique d'où vient le problème) $\rightarrow$ \textbf{solutions} (on propose des remèdes). C'est le plan typique de l'éditorial de presse et du discours politique.
\end{definition}

\begin{definition}[Plan dialectique]
Le \cle{plan dialectique} met en débat deux points de vue opposés : \textbf{thèse} (une première opinion) $\rightarrow$ \textbf{antithèse} (l'opinion contraire) $\rightarrow$ \textbf{synthèse} (la position de l'auteur, qui dépasse ou tranche l'opposition). C'est le plan de la dissertation au Baccalauréat technique.
\end{definition}

\begin{definition}[Plan thématique]
Le \cle{plan thematique} examine successivement les différents \textbf{aspects} d'un même sujet, sans opposition ni problème à résoudre : par exemple, les aspects économiques, sociaux et culturels d'une question. Les arguments s'y succèdent comme les chapitres d'un exposé.
\end{definition}

\begin{popfigure}
\begin{center}
\begin{tikzpicture}[
  bloc/.style={draw=popblue, fill=popblueL, rounded corners=2pt, text width=2.6cm, align=center, minimum height=0.9cm, font=\small},
  fleche/.style={-{Stealth[length=2.5mm]}, thick, popdark},
  titre/.style={font=\bfseries\small, popdark}
]
% Plan analytique
\node[titre] at (0,4.4) {Plan analytique};
\node[bloc, align=center] (c1) at (0,3.4){Constat\\ (le problème)};
\node[bloc, align=center] (c2) at (0,2.1){Causes\\ (explications)};
\node[bloc, align=center] (c3) at (0,0.8){Solutions\\ (propositions)};
\draw[fleche] (c1) -- (c2);
\draw[fleche] (c2) -- (c3);
% Plan dialectique
\node[titre] at (5.2,4.4) {Plan dialectique};
\node[bloc, draw=poporange, fill=poporangeL, align=center] (d1) at (5.2,3.4){Thèse\\ (1\textsuperscript{re} opinion)};
\node[bloc, draw=poporange, fill=poporangeL, align=center] (d2) at (5.2,2.1){Antithèse\\ (opinion contraire)};
\node[bloc, draw=poporange, fill=poporangeL, align=center] (d3) at (5.2,0.8){Synthèse\\ (position de l'auteur)};
\draw[fleche] (d1) -- (d2);
\draw[fleche] (d2) -- (d3);
% Plan thématique
\node[titre] at (10.4,4.4) {Plan thématique};
\node[bloc, draw=popgreen, fill=popgreenL, align=center] (t1) at (10.4,3.4){Aspect 1\\ (économique)};
\node[bloc, draw=popgreen, fill=popgreenL, align=center] (t2) at (10.4,2.1){Aspect 2\\ (social)};
\node[bloc, draw=popgreen, fill=popgreenL, align=center] (t3) at (10.4,0.8){Aspect 3\\ (culturel)};
\draw[fleche] (t1) -- (t2);
\draw[fleche] (t2) -- (t3);
\end{tikzpicture}
\end{center}
\legende{Schéma comparatif des trois plans argumentatifs. Dans le plan thématique, l'ordre des aspects est plus libre ; dans les plans analytique et dialectique, l'ordre est logique et ne peut pas être inversé.}
\end{popfigure}

\begin{attention}[Ne pas confondre les plans]
Le plan \textbf{dialectique} suppose une \emph{opposition} réelle entre deux thèses (pour / contre). Si le texte expose simplement plusieurs facettes d'un sujet sans les opposer, c'est un plan \textbf{thématique}. Et si le texte part d'un problème pour aboutir à des solutions, c'est un plan \textbf{analytique}. À l'examen, justifie toujours ton choix par des indices précis du texte.
\end{attention}

\courssub{Les connecteurs logiques : les panneaux de signalisation du texte}

Comment repérer concrètement la structure argumentative dans un texte ? En suivant les \cle{connecteurs logiques}, ces petits mots qui fonctionnent comme des panneaux de signalisation pour le lecteur.

\begin{center}
\begin{tabularx}{\linewidth}{|l|X|X|}
\hline
\rowcolor{popblueL} \textbf{Relation logique} & \textbf{Connecteurs fréquents} & \textbf{Fonction dans l'argumentation} \\
\hline
Cause, justification & car, parce que, en effet, puisque & introduisent un \textbf{argument} \\
\hline
Conséquence & donc, par conséquent, c'est pourquoi, ainsi & annoncent une \textbf{conclusion} ou la thèse \\
\hline
Opposition & mais, cependant, pourtant, en revanche, néanmoins & signalent l'\textbf{antithèse} ou une objection \\
\hline
Illustration & par exemple, ainsi, c'est le cas de & introduisent un \textbf{exemple} (idée secondaire) \\
\hline
Addition & d'abord, ensuite, de plus, enfin, en outre & numérotent les \textbf{arguments} \\
\hline
\end{tabularx}
\end{center}

\begin{methode}[Dégager la structure argumentative d'un texte]
\begin{enumerate}
\item \textbf{Lire} le texte deux fois, puis \textbf{souligner la thèse} : elle se trouve souvent au début ou à la fin, et elle est fréquemment introduite par « il faut », « nous devons », « c'est pourquoi ».
\item \textbf{Numéroter les arguments} dans la marge (A1, A2, A3...) en repérant les connecteurs de cause (« car », « en effet ») et d'addition (« d'abord », « ensuite », « enfin »).
\item \textbf{Barrer les exemples, les répétitions et les digressions} : tout ce qui est introduit par « par exemple » ou qui reformule une idée déjà exprimée est secondaire.
\item \textbf{Identifier le plan} : analytique, dialectique ou thématique, en observant l'enchaînement des grandes parties.
\item \textbf{Construire le schéma argumentatif} : la thèse au sommet, les arguments au niveau intermédiaire, les exemples en bas.
\end{enumerate}
\end{methode}

\courssub{Application : dégager la structure d'un éditorial de presse camerounaise}

Mettons la méthode en pratique sur un texte d'actualité, de type éditorial publié dans un quotidien camerounais.

\begin{exemplebox}[Texte à analyser (éditorial, environ 180 mots)]
« \textbf{(1)} Le Cameroun doit urgemment réduire sa dépendance au riz importé. \textbf{(2)} En effet, la facture des importations alimentaires pèse lourdement sur la balance commerciale du pays : chaque année, des centaines de milliards de francs CFA partent à l'étranger pour acheter un produit que nos sols peuvent fournir. \textbf{(3)} Par exemple, la seule campagne 2022 a coûté plus de 300 milliards de francs CFA aux importateurs. \textbf{(4)} De plus, nos plaines rizicoles restent tragiquement sous-exploitées : à Yagoua, dans l'Extrême-Nord, ou à Ndom, dans le Littoral, des milliers d'hectares attendent l'eau et les intrants. \textbf{(5)} Certes, les partisans du libre-échange objectent que le riz importé est moins cher pour le consommateur. \textbf{(6)} Cependant, cet argument oublie que chaque franc dépensé à l'importation est un emploi perdu pour nos jeunes ruraux. \textbf{(7)} C'est pourquoi l'État doit investir dans l'irrigation, encadrer le crédit agricole et protéger le riz local. \textbf{(8)} L'autosuffisance rizicole n'est pas un rêve : c'est une décision politique. »
\end{exemplebox}

Appliquons la méthode pas à pas. La \textbf{thèse} est énoncée dès la phrase (1) et reformulée en conclusion (8) : \emph{le Cameroun doit réduire sa dépendance au riz importé}. Le connecteur « \emph{En effet} » (2) annonce le premier argument (le coût des importations), illustré par un exemple chiffré (3) introduit par « \emph{par exemple} ». « \emph{De plus} » (4) signale le deuxième argument (les terres sous-exploitées), lui-même illustré par les cas de Yagoua et Ndom. « \emph{Certes... Cependant} » (5--6) marque l'objection et sa réfutation : c'est la trace d'un raisonnement qui intègre l'antithèse. Enfin, « \emph{C'est pourquoi} » (7) introduit les solutions, c'est-à-dire la conclusion. Le plan est donc \textbf{analytique} (constat $\rightarrow$ causes $\rightarrow$ solutions), avec une courte séquence dialectique interne (5--6).

\begin{popfigure}
\begin{center}
\begin{tikzpicture}[
  these/.style={draw=popdark, fill=poppurpleL, rounded corners=3pt, text width=6.5cm, align=center, font=\small\bfseries},
  arg/.style={draw=popblue, fill=popblueL, rounded corners=2pt, text width=4.2cm, align=center, font=\small},
  ex/.style={draw=popmuted, fill=popcream, rounded corners=2pt, text width=4.0cm, align=center, font=\footnotesize},
  lien/.style={-{Stealth[length=2mm]}, popdark, thick}
]
\node[these] (T) at (0,0) {THÈSE : Le Cameroun doit réduire sa dépendance au riz importé};
\node[arg] (A1) at (-4.2,-2.2) {ARG. 1 : Les importations coûtent des milliards au pays};
\node[arg] (A2) at (4.2,-2.2) {ARG. 2 : Les plaines rizicoles locales sont sous-exploitées};
\node[ex] (E1) at (-4.2,-4.2) {Ex. : 300 milliards de FCFA en 2022};
\node[ex] (E2) at (4.2,-4.2) {Ex. : Yagoua, Ndom (milliers d'hectares)};
\node[arg, draw=poporange, fill=poporangeL] (A3) at (0,-6.2) {RÉFUTATION : « le riz importé est moins cher » $\rightarrow$ mais chaque franc importé = un emploi rural perdu};
\draw[lien] (T) -- (A1);
\draw[lien] (T) -- (A2);
\draw[lien] (A1) -- (E1);
\draw[lien] (A2) -- (E2);
\draw[lien] (T) -- (A3);
\end{tikzpicture}
\end{center}
\legende{Schéma argumentatif (arbre) de l'éditorial : la thèse au sommet, les arguments au niveau intermédiaire, les exemples en feuilles. Pour le résumé, on conserve les cadres violets et bleus ; on élimine les cadres crème (exemples).}
\end{popfigure}

\courssub{Idées essentielles et idées secondaires : le crible du résumé}

Tout le travail précédent aboutit à une distinction capitale pour la contraction de texte.

\begin{definition}[Idées essentielles et idées secondaires]
Les \cle{idees essentielles} sont la \textbf{thèse} et les \textbf{arguments} : sans elles, le raisonnement s'effondre. Les \cle{idees secondaires} sont les \textbf{exemples}, les \textbf{répétitions} (une même idée reformulée), les \textbf{digressions} (anecdotes, précisions anecdotiques) et les détails chiffrés ou descriptifs : on peut les supprimer sans détruire le sens global.
\end{definition}

\begin{exemplebox}[Un exemple chiffré]
Dans un texte argumentatif de \num{500} mots, la thèse et les arguments représentent souvent \textbf{moins de 200 mots}. Le reste — plus de la moitié du texte ! — est constitué d'exemples, de reformulations et d'illustrations que le résumé doit éliminer. C'est pourquoi savoir distinguer l'essentiel du secondaire est la compétence décisive de la contraction.
\end{exemplebox}

\begin{attention}[L'erreur classique : confondre exemple et argument]
Beaucoup d'élèves recopient dans leur résumé la phrase « la campagne 2022 a coûté 300 milliards » en la prenant pour une idée essentielle. Or c'est un \textbf{exemple} : il illustre l'argument « les importations coûtent cher », mais il ne le remplace pas. \textbf{Test pratique} : supprime la phrase. Si le raisonnement tient encore debout, c'était une idée secondaire ; s'il s'effondre, c'était un argument.
\end{attention}

\begin{savaistu}[Le résumé au Baccalauréat technique]
À l'épreuve de Français du Probatoire et du Baccalauréat technique, le résumé exigé représente généralement le \textbf{quart du texte initial}, avec une marge de tolérance de $\pm 10\,\%$. Pour un texte de 400 mots, il faut donc produire environ 100 mots (entre 90 et 110). Compter ses mots et respecter cette contrainte fait partie de la note !
\end{savaistu}

\begin{exoresolu}[Identifier thèse, arguments et exemples]
\textbf{Énoncé.} Dans le texte suivant, indiquez la thèse, les deux arguments et les deux exemples : « Le téléphone portable en classe doit être strictement interdit. D'abord, il détourne l'attention des élèves : une étude menée dans un lycée de Douala montre qu'un élève consulte son écran en moyenne toutes les six minutes. Ensuite, il favorise la triche pendant les évaluations, comme l'a montré la fraude massive déjouée au dernier Baccalauréat dans la région du Centre. L'école doit donc rester un sanctuaire de la concentration. »
\tcblower
\textbf{Solution détaillée.}
\begin{itemize}
\item \textbf{Thèse} : « Le téléphone portable en classe doit être strictement interdit » (reformulée en conclusion : « L'école doit rester un sanctuaire de la concentration »).
\item \textbf{Argument 1} (introduit par « D'abord ») : le téléphone détourne l'attention des élèves. \textbf{Exemple associé} : l'étude du lycée de Douala (une consultation toutes les six minutes).
\item \textbf{Argument 2} (introduit par « Ensuite ») : le téléphone favorise la triche. \textbf{Exemple associé} : la fraude déjouée au Baccalauréat dans la région du Centre.
\end{itemize}
Pour un résumé, on conserverait : thèse + argument 1 + argument 2 + conclusion, soit environ 30 mots sur un texte d'environ 90 mots.
\end{exoresolu}

\begin{exercice}[À vous de jouer]
Relevez dans un journal camerounais (ou dans le manuel) un éditorial d'environ 300 mots. 1) Soulignez la thèse. 2) Numérotez les arguments en repérant les connecteurs logiques. 3) Barrez les exemples et les répétitions. 4) Identifiez le plan (analytique, dialectique ou thématique) et justifiez votre réponse. 5) Dessinez le schéma argumentatif du texte.
\end{exercice}

\begin{corrige}[Exercice — éléments attendus]
Le corrigé dépend du texte choisi, mais la copie doit présenter : une thèse citée entre guillemets ; des arguments numérotés et formulés en une phrase chacun ; la liste des connecteurs repérés avec leur fonction (cause, opposition, illustration...) ; le nom du plan justifié par l'enchaînement des parties ; un schéma en arbre conforme au modèle du cours (thèse au sommet, arguments au niveau 2, exemples en feuilles). Toute confusion exemple/argument est sanctionnée : appliquer le test de suppression.
\end{corrige}

Maintenant que nous savons \textbf{dégager} la structure argumentative d'un texte, il reste à apprendre les outils linguistiques qui permettront de \textbf{reformuler} ces idées essentielles avec nos propres mots : le champ lexical, le champ sémantique et les relations de sens feront l'objet de la section suivante.

\coursec{Champ lexical, champ sémantique et relations de sens}

Dans la section précédente, nous avons appris à repérer la \cle{structure argumentative} d'un texte : la thèse, les arguments, les exemples. Mais pour résumer un texte, il ne suffit pas de savoir \emph{où} sont les idées ; il faut encore pouvoir les \emph{reformuler} sans les trahir. Or reformuler, c'est jouer avec les mots : les regrouper, les remplacer, les généraliser. C'est précisément ce que permettent les notions de \cle{champ lexical}, de \cle{champ sémantique} et de \cle{relations de sens} (hyperonymie, hyponymie, paronymie) que nous allons étudier maintenant. Ces outils de la langue deviendront, dans les sections suivantes, de véritables techniques de réduction du texte.

\courssub{Le champ lexical : des mots qui tournent autour d'une même idée}

\textbf{Intuition.} Lisez cette phrase inspirée de l'univers de \emph{Le Lion et la perle} : « Le soir, les anciens s'asseyaient sous le baobab, les griots chantaient les ancêtres, et les jeunes filles exécutaient la danse du mariage devant le chef du village. » Même sans qu'on vous dise le sujet du passage, vous devinez immédiatement qu'il parle de \emph{la tradition}. Pourquoi ? Parce que plusieurs mots (\emph{anciens, griots, ancêtres, danse, mariage, chef du village}) renvoient tous à cette même notion. Ce faisceau de mots, c'est un champ lexical.

\begin{definition}[Champ lexical]
Le \cle{champ lexical} est l'ensemble des mots et expressions qui, dans un texte, se rapportent à une même notion, un même thème ou une même idée. Ces mots peuvent être de natures différentes : noms (\emph{le griot}), verbes (\emph{chanter}), adjectifs (\emph{ancestral}), adverbes (\emph{traditionnellement}).
\end{definition}

\begin{methode}[Repérer un champ lexical dans un texte]
\begin{enumerate}
\item Lire le texte une première fois et dégager le thème général.
\item Souligner tous les mots et expressions qui renvoient à ce thème, quelle que soit leur nature grammaticale.
\item Regrouper ces mots en une liste.
\item Vérifier que chaque mot de la liste évoque bien la même notion ; éliminer les mots qui appartiennent à un autre thème.
\item Nommer le champ lexical par un terme général (ex. : « le champ lexical de la tradition »).
\end{enumerate}
\end{methode}

\begin{exemplebox}[Champ lexical de la tradition dans \emph{Le Lion et la perle}]
Dans la pièce de Wole Soyinka, l'univers du village d'Ilujinle est traversé par un riche champ lexical de la tradition : \emph{les anciens, la coutume, le mariage coutumier, la dot, la danse, le tam-tam, les masques, les ancêtres, le chef (le Bale), la sagesse}. À l'opposé, un autre champ lexical, celui de la \emph{modernité}, se construit autour de mots comme \emph{l'école, le progrès, la photographie, la ville, le train, les manières occidentales}. Le conflit entre Baroka (le Lion) et Lakunle (l'instituteur) est donc aussi un conflit entre deux champs lexicaux.
\end{exemplebox}

\begin{popfigure}
\begin{tikzpicture}[mindmap, grow cyclic, every node/.style={concept, font=\small},
root concept/.append style={concept color=popblue, font=\bfseries},
level 1/.append style={level distance=3.2cm, sibling angle=90, concept color=popblueL},
level 2/.append style={level distance=2.2cm, sibling angle=50, concept color=poporangeL}]
\node[root concept]{La tradition}
  child { node {Personnes}
    child { node {les anciens} }
    child { node {le griot} }
    child { node {le Bale (chef)} }
  }
  child { node {Pratiques}
    child { node {la dot} }
    child { node {la danse} }
    child { node {le tam-tam} }
  }
  child { node {Croyances}
    child { node {les ancêtres} }
    child { node {les masques} }
  }
  child { node {Valeurs}
    child { node {la coutume} }
    child { node {la sagesse} }
  };
\end{tikzpicture}
\legende{Carte sémantique du champ lexical de « la tradition » dans \emph{Le Lion et la perle} : les mots du champ se regroupent en sous-thèmes (personnes, pratiques, croyances, valeurs).}
\end{popfigure}

\begin{attention}[Un champ lexical n'est pas une liste de synonymes]
Les mots d'un champ lexical ne se ressemblent pas forcément et ne veulent pas dire la même chose : \emph{griot}, \emph{dot} et \emph{tam-tam} ne sont pas synonymes, mais ils renvoient tous à la tradition. Ne confonds pas \emph{champ lexical} (mots liés à un même thème) et \emph{synonymie} (mots de sens équivalent).
\end{attention}

\courssub{Le champ sémantique : le sens partagé au-delà des mots}

\textbf{Intuition.} Prenons trois mots : \emph{vieillard} (nom), \emph{se transmettre} (verbe), \emph{ancestral} (adjectif). Ils n'ont ni la même nature ni la même forme. Pourtant, ils partagent quelque chose : l'idée de \emph{ce qui vient du passé, de l'ancien temps}. Ce « quelque chose » de commun, c'est un sème, un trait de sens.

\begin{definition}[Champ sémantique]
Le \cle{champ sémantique} est l'ensemble des traits de sens (appelés \cle{sèmes}) communs à plusieurs mots de natures grammaticales différentes. Là où le champ lexical s'intéresse aux \emph{mots} réunis autour d'un thème, le champ sémantique s'intéresse au \emph{sens} que ces mots partagent.
\end{definition}

\begin{exemplebox}[Champ lexical et champ sémantique : la différence]
\begin{itemize}
\item Champ lexical de la tradition : \emph{griot, dot, ancêtres, danse, coutume} (des mots autour d'un thème).
\item Champ sémantique de l'ancienneté : \emph{vieux, ancien, ancestral, vieillir, archaïque, antérieur} partagent le sème « + ancien » ; \emph{transmettre, hériter, léguer} partagent le sème « + passage d'une génération à l'autre ».
\end{itemize}
\end{exemplebox}

\begin{aretenir}[Lien entre les deux notions]
Dans un texte, le champ lexical d'un thème est souvent construit à partir de mots dont les champs sémantiques se recoupent. Pour résumer, repérer le sens commun (le champ sémantique) permet de trouver le mot général capable de remplacer plusieurs mots précis.
\end{aretenir}

\courssub{Hyperonymie et hyponymie : du général au particulier}

\textbf{Intuition.} Si je dis : « Le marché de Mokolo vend des mangues, des bananes, des safous et des avocats », je peux dire plus court : « Le marché de Mokolo vend des fruits. » Le mot \emph{fruit} contient en lui tous les autres. Cette relation du général au particulier est fondamentale pour la contraction de texte.

\begin{definition}[Hyperonymie et hyponymie]
\begin{itemize}
\item Un \cle{hyperonyme} est un mot au sens \emph{général} qui englobe le sens d'autres mots. Exemple : \emph{fruit} est l'hyperonyme de \emph{mangue}.
\item Un \cle{hyponyme} est un mot au sens \emph{particulier}, inclus dans le sens de l'hyperonyme. Exemples : \emph{mangue, banane, safou} sont des hyponymes (ou co-hyponymes) de \emph{fruit}.
\end{itemize}
On parle d'\cle{hyperonymie} (relation du particulier vers le général) et d'\cle{hyponymie} (relation du général vers le particulier).
\end{definition}

\begin{popfigure}
\begin{tikzpicture}[
every node/.style={font=\small},
hyp/.style={draw=popdark, fill=popgreenL, rounded corners=2pt, minimum width=2.1cm, minimum height=0.7cm},
hyper/.style={draw=popdark, fill=popblue, text=white, font=\bfseries\small, rounded corners=2pt, minimum width=4.2cm, minimum height=0.8cm},
lien/.style={-{Stealth[length=2mm]}, thick, popdark}]
\node[hyper] (h) {Produits agricoles du Cameroun};
\node[hyp, below left=1.4cm and 2.6cm of h] (a) {cacao};
\node[hyp, below left=1.4cm and 0.2cm of h] (b) {café};
\node[hyp, below right=1.4cm and 0.2cm of h] (c) {banane};
\node[hyp, below right=1.4cm and 2.6cm of h] (d) {coton};
\draw[lien] (h.south) -- (a.north);
\draw[lien] (h.south) -- (b.north);
\draw[lien] (h.south) -- (c.north);
\draw[lien] (h.south) -- (d.north);
\node[font=\footnotesize\itshape, popmuted, above=0.15cm of h] {Hyperonyme (terme générique)};
\node[font=\footnotesize\itshape, popmuted, below=0.15cm of b.south west, anchor=north west, xshift=-2.6cm] {Hyponymes (termes spécifiques)};
\end{tikzpicture}
\legende{Arbre hiérarchique : l'hyperonyme « produits agricoles du Cameroun » domine ses hyponymes. Pour résumer, on remonte la flèche : on remplace les termes spécifiques par le terme générique.}
\end{popfigure}

\begin{methode}[Condenser par généralisation]
Pour réduire un texte sans perdre l'essentiel :
\begin{enumerate}
\item Repérer dans la phrase une \emph{énumération} de termes précis (liste de noms, souvent séparés par des virgules ou par « et »).
\item Vérifier que ces termes sont des \emph{co-hyponymes}, c'est-à-dire qu'ils relèvent d'un même hyperonyme.
\item Trouver l'hyperonyme qui les englobe tous.
\item Remplacer toute l'énumération par cet hyperonyme, éventuellement accompagné d'un adjectif précisant (\emph{des arbres fruitiers}, \emph{des produits agricoles}).
\item Relire pour vérifier qu'aucune idée essentielle n'est perdue : si l'un des termes de la liste a une importance particulière dans l'argumentation, il faut le conserver.
\end{enumerate}
\end{methode}

\begin{exemplebox}[La généralisation en action]
Phrase du texte : « Manguiers, palmiers, bananiers et goyaviers bordaient la cour. »\\
Résumé : « Des \cle{arbres fruitiers} bordaient la cour. »\\
Quatre hyponymes (\emph{manguiers, palmiers, bananiers, goyaviers}) ont été remplacés par un seul hyperonyme (\emph{arbres fruitiers}) : la phrase passe de 9 à 6 mots, et l'idée essentielle est intacte.
\end{exemplebox}

\begin{attention}[Ne pas généraliser à l'excès]
L'hyperonyme choisi ne doit pas être \emph{trop} général. Remplacer « mangues, bananes, safous » par « des choses » ou « des produits » fait perdre une information essentielle. Choisis toujours l'hyperonyme le plus \emph{proche} de la liste : \emph{fruits} convient, \emph{végétaux} est déjà trop large. De même, si le texte insiste sur un élément précis (par exemple le \emph{cacao} dans un passage sur l'économie camerounaise), conserve ce terme dans ton résumé.
\end{attention}

\courssub{La paronymie : attention aux faux amis de la langue}

\textbf{Intuition.} « Le différend entre les deux villages » et « le différent entre les deux villages » : une seule lettre change, mais la seconde phrase est fautive. Certains mots se ressemblent tellement par la forme qu'on les confond, alors que leur sens est totalement différent. Dans une reformulation, cette confusion peut produire un véritable \emph{contresens}.

\begin{definition}[Paronymie]
La \cle{paronymie} est la relation entre des mots de \emph{formes voisines} (même racine, prononciation ou orthographe proches) mais de \emph{sens différents}. Ces mots sont appelés \cle{paronymes}. Exemples : \emph{éminence} (ce qui est élevé, distingué) / \emph{imminence} (ce qui va arriver très vite) ; \emph{collision} (choc brutal) / \emph{collusion} (entente secrète et malhonnête) ; \emph{différend} (désaccord, conflit) / \emph{différent} (qui n'est pas pareil).
\end{definition}

\begin{exemplebox}[Quelques paronymes à connaître]
\begin{center}
\begin{tabularx}{\linewidth}{|l|X|l|X|}\hline
\rowcolor{popblueL}
\textbf{Paronyme 1} & \textbf{Sens} & \textbf{Paronyme 2} & \textbf{Sens} \\ \hline
éminence & hauteur, distinction & imminence & arrivée proche \\ \hline
collision & choc entre deux corps & collusion & entente secrète \\ \hline
différend & conflit, désaccord & différent & pas semblable \\ \hline
invention & création nouvelle & intention & volonté de faire \\ \hline
abjurer & renier solennellement & adjurer & supplier instamment \\ \hline
\end{tabularx}
\end{center}
\end{exemplebox}

\begin{attention}[Erreur fréquente : confondre paronymes et synonymes]
Un synonyme a le \emph{même sens} qu'un autre mot (\emph{conflit} $\approx$ \emph{différend}) ; un paronyme a seulement une \emph{forme proche} (\emph{différend} $\neq$ \emph{différent}). Dans un résumé, remplacer un mot du texte par son paronyme est une faute grave : cela change le sens de l'idée de l'auteur et produit un contresens. Avant d'employer un mot de remplacement, demande-toi toujours : « Ce mot signifie-t-il vraiment la même chose, ou lui ressemble-t-il seulement ? » En cas de doute, vérifie dans le dictionnaire.
\end{attention}

\courssub{Utilité de ces notions pour la contraction de texte}

Toutes ces notions convergent vers un même objectif : t'aider à \emph{reformuler les idées essentielles} d'un texte en moins de mots, sans les déformer.

\begin{aretenir}[Les trois usages en résumé]
\begin{enumerate}
\item Le \cle{champ lexical} permet d'identifier le thème d'un passage et de vérifier que ton résumé reste centré sur ce thème.
\item L'\cle{hyperonymie} permet de \emph{généraliser} : une énumération d'hyponymes est remplacée par l'hyperonyme (technique de réduction par condensation).
\item La \cle{paronymie} impose la \emph{vigilance} : dans la reformulation, tout mot de remplacement doit être contrôlé pour éviter le contresens.
\end{enumerate}
\end{aretenir}

\begin{exoresolu}[Classer les mots d'un texte et condenser]
\textbf{Énoncé.} Voici un extrait adapté de l'univers d'Ilujinle : « Autour de la place du village, les femmes vendaient des mangues, des bananes, des safous et des papayes, tandis que résonnaient les tam-tams et que les danseurs, coiffés de masques, célébraient les ancêtres devant les anciens rassemblés. »\\
1) Relevez les mots appartenant au champ lexical de la tradition.\\
2) Relevez l'énumération et remplacez-la par un hyperonyme pour condenser la phrase.\\
3) Le mot \emph{célébraient} peut-il être remplacé par \emph{accéléraient} ? Justifiez.
\tcblower
\textbf{Solution détaillée.}\\
1) Champ lexical de la tradition : \emph{les tam-tams, les danseurs, les masques, les ancêtres, les anciens} (on peut ajouter \emph{la place du village} comme cadre coutumier). Ces mots, de natures différentes (noms communs concrets ou abstraits), renvoient tous à la même notion.\\
2) L'énumération « des mangues, des bananes, des safous et des papayes » est une liste de quatre co-hyponymes. Leur hyperonyme est \emph{fruits}. Phrase condensée : « Autour de la place du village, les femmes vendaient des fruits, tandis que les danseurs masqués célébraient les ancêtres devant les anciens, au son des tam-tams. » L'idée essentielle est conservée en moins de mots.\\
3) Non. \emph{Célébrer} (honorer par une fête, un rite) et \emph{accélérer} (rendre plus rapide) sont des paronymes : leurs formes sont proches mais leurs sens sont différents. Le remplacement produirait un contresens total.
\end{exoresolu}

\begin{exercice}[À toi de jouer]
1) Dans la phrase suivante, relève les mots du champ lexical de l'école et de la modernité : « Lakunle, l'instituteur, voulait apporter le progrès au village : l'école, les livres, les manières de la ville et le refus de la dot. »\\
2) Construis un arbre hyperonyme $\rightarrow$ hyponymes à partir du mot \emph{bétail} (au moins trois hyponymes).\\
3) Condense : « Dans les champs, on cultivait le maïs, le manioc, l'igname et l'arachide. »\\
4) Choisis le bon paronyme : « La (collision / collusion) entre les deux motos a bloqué la route. » — « Les deux commerçants étaient en (collision / collusion) pour tromper les clients. »
\end{exercice}

\begin{corrige}[Exercice : champ lexical, hyperonymes et paronymes]
1) Champ lexical de l'école et de la modernité : \emph{instituteur, progrès, école, livres, manières de la ville}. Le \emph{refus de la dot} appartient au thème du conflit entre modernité et tradition.\\
2) Hyperonyme : \emph{bétail} ; hyponymes possibles : \emph{bœufs, chèvres, moutons, porcs, volailles}.\\
3) « Dans les champs, on cultivait des \cle{produits vivriers} (ou des \emph{cultures vivrières}). » Les quatre hyponymes sont remplacés par l'hyperonyme.\\
4) « La \emph{collision} entre les deux motos a bloqué la route » (choc). « Les deux commerçants étaient en \emph{collusion} pour tromper les clients » (entente secrète).
\end{corrige}

\begin{savaistu}[Le vocabulaire, richesse des langues du Cameroun]
Le Cameroun compte plus de 250 langues nationales. Dans chacune, les champs lexicaux reflètent la culture : les langues des peuples forestiers possèdent de nombreux hyponymes pour désigner les variétés de palmiers ou de bananes plantains, là où le français n'a qu'un seul mot. Comprendre les relations de sens, c'est aussi comprendre comment une communauté organise son monde.
\end{savaistu}

En résumé, le champ lexical et le champ sémantique t'aident à \emph{voir} les idées derrière les mots ; l'hyperonymie te donne un outil puissant de \emph{condensation} ; la paronymie t'impose une \emph{discipline de précision}. Nous pouvons maintenant passer à l'outil grammatical de la reformulation : la phrase complexe.

\coursec{La phrase complexe au service de la reformulation}

Dans la section précédente, nous avons appris à repérer les champs lexicaux et les relations de sens (paronymie, hyperonymie, hyponymie) pour comprendre comment les mots s'organisent autour d'une idée. Mais dans un texte argumentatif comme \emph{Le Lion et la perle}, les idées ne sont pas seulement portées par des mots : elles sont construites dans des \cle{phrases}. Or, pour contracter un texte, il ne suffit pas de remplacer des mots par des synonymes ; il faut savoir \textbf{découper la phrase}, repérer ce qui est essentiel et ce qui est secondaire. C'est précisément le rôle de l'analyse de la \cle{phrase complexe} : elle est l'outil grammatical principal de la reformulation. Cette section vous donne la méthode complète, pas à pas.

\courssub{Qu'est-ce qu'une phrase complexe ?}

Commençons par une intuition simple. Comparez ces deux énoncés :

\begin{itemize}
\item « Le chef refusa la proposition. » (un seul verbe conjugué : \emph{refusa})
\item « Le chef, qui était très âgé, refusa la proposition que lui avait faite l'étranger. » (trois verbes conjugués : \emph{était}, \emph{refusa}, \emph{avait faite})
\end{itemize}

Le premier énoncé est une \cle{phrase simple} : un seul verbe conjugué, une seule proposition. Le second est une \cle{phrase complexe} : plusieurs verbes conjugués, donc plusieurs propositions réunies dans une même phrase.

\begin{definition}[Phrase complexe]
Une \cle{phrase complexe} est une phrase qui comporte \textbf{au moins deux verbes conjugués} répartis en plusieurs \cle{propositions}. Une proposition est un groupe de mots organisé autour d'un verbe conjugué, pouvant fonctionner comme une phrase autonome sur le plan grammatical.
\end{definition}

\begin{attention}[Ne pas confondre]
Attention à deux pièges fréquents :
\begin{itemize}
\item Un verbe à l'\textbf{infinitif} ou au \textbf{participe} ne suffit pas à créer une proposition conjuguée. « Voulant plaire, le chef refusa » contient un participe présent (\emph{voulant}) et un verbe conjugué (\emph{refusa}) : on parle de proposition \emph{participiale}, rattachée à la phrase complexe.
\item Le nombre de \textbf{mots} ne compte pas : une phrase de trente mots avec un seul verbe conjugué reste une phrase simple.
\end{itemize}
\end{attention}

\courssub{Les types de propositions : rappels}

Pour analyser une phrase complexe, il faut d'abord distinguer les différentes propositions qui la composent.

\textbf{1. La proposition principale.} C'est la proposition qui ne dépend d'aucune autre. Elle porte le verbe « pivot » de la phrase. Exemple : « Le chef \textbf{refusa} la proposition. »

\textbf{2. Les propositions subordonnées.} Elles dépendent d'une autre proposition et ne peuvent pas être employées seules. On en distingue trois grandes familles au programme de Première technique :

\begin{center}
\begin{tabularx}{\linewidth}{|l|X|X|}
\hline
\rowcolor{popblueL} \textbf{Type} & \textbf{Marque de reconnaissance} & \textbf{Exemple} \\
\hline
Subordonnée \cle{relative} & Introduite par un pronom relatif : \emph{qui, que, où, dont, lequel} & « Le chef \textbf{qui était âgé} refusa. » \\
\hline
Subordonnée \cle{conjonctive} & Introduite par une conjonction de subordination : \emph{que, parce que, bien que, si, lorsque, puisque} & « Le chef refusa \textbf{parce qu'il se méfiait de l'étranger}. » \\
\hline
Subordonnée \cle{participiale} & Construite autour d'un participe présent ou passé & « \textbf{Connaissant les traditions}, le chef refusa. » \\
\hline
\end{tabularx}
\end{center}

\textbf{3. La coordination et la juxtaposition.} Toutes les propositions d'une phrase complexe ne sont pas subordonnées. Deux propositions peuvent être :

\begin{itemize}
\item \textbf{coordonnées} : reliées par une conjonction de coordination (\emph{mais, ou, et, donc, or, ni, car}) : « Le chef était âgé, \textbf{mais} il restait lucide. »
\item \textbf{juxtaposées} : placées côte à côte, séparées par une virgule ou un point-virgule : « Le chef écouta, il réfléchit, il refusa. »
\end{itemize}

\begin{aretenir}[Le test du verbe conjugué]
Pour compter les propositions d'une phrase complexe, soulignez tous les \textbf{verbes conjugués} : il y a autant de propositions que de verbes conjugués. Ensuite, demandez-vous pour chacun : « Cette proposition peut-elle exister seule ? » Si oui, c'est une principale ou une coordonnée ; si non, c'est une subordonnée.
\end{aretenir}

\courssub{Hiérarchiser les idées : la principale porte l'essentiel}

Voici le principe qui fonde toute la technique de la contraction de texte. Faisons une petite expérience de pensée : prenez la phrase « Baroka, qui était le chef du village et qui redoutait la modernité, s'opposa à l'ouverture de l'école. » Supprimons les subordonnées : il reste « Baroka s'opposa à l'ouverture de l'école. » Le sens fondamental survit. En revanche, supprimons la principale : il reste « qui était le chef du village et qui redoutait la modernité » : l'information essentielle (que s'est-il passé ?) a disparu.

\begin{propriete}[Principe de hiérarchisation]
Dans une phrase complexe, la \cle{proposition principale} porte généralement \textbf{l'idée essentielle}. Les propositions subordonnées apportent des précisions secondaires (description, circonstance, explication) qui peuvent, dans un résumé, être \textbf{supprimées} ou \textbf{condensées}. Attention toutefois : certaines subordonnées de cause, de condition ou d'opposition peuvent contenir une nuance indispensable au raisonnement de l'auteur.
\end{propriete}

\begin{exemplebox}[Réduction d'une phrase complexe]
Phrase de départ : « Le chef, qui était très âgé et qui connaissait toutes les traditions, refusa catégoriquement la proposition. »

Résumé : « Le chef refusa la proposition. »

Les deux relatives (\emph{qui était très âgé}, \emph{qui connaissait toutes les traditions}) sont descriptives : elles caractérisent le chef mais ne font pas avancer l'idée. L'adverbe \emph{catégoriquement} est une nuance d'intensité, supprimable dans un résumé très serré.
\end{exemplebox}

\begin{popfigure}
\begin{center}
\begin{tikzpicture}[
bloc/.style={draw=popdark, rounded corners=2pt, align=left, inner sep=5pt, font=\small},
princ/.style={bloc, fill=popblueL, draw=popblue, thick},
sub/.style={bloc, fill=poporangeL, draw=poporange},
lien/.style={-{Stealth[length=2.5mm]}, thick, popmuted}
]
\node[princ, text width=9.5cm] (P) at (0,0) {\textbf{Proposition principale (idée essentielle)}\\ « Baroka \textbf{s'opposa} à l'ouverture de l'école. »};
\node[sub, text width=4.6cm] (S1) at (-3.2,-2.3) {\textbf{Subordonnée relative}\\ « qui était le chef du village »\\ \emph{(précision : identité)}};
\node[sub, text width=4.6cm] (S2) at (3.2,-2.3) {\textbf{Subordonnée relative}\\ « qui redoutait la modernité »\\ \emph{(précision : motif psychologique)}};
\draw[lien] (S1.north) -- (P.south west);
\draw[lien] (S2.north) -- (P.south east);
\node[font=\small\itshape, popmuted] at (0,-3.8) {Résumé : on garde le bloc bleu, on supprime ou condense les blocs oranges.};
\end{tikzpicture}
\end{center}
\legende{Schéma d'analyse d'une phrase complexe : la proposition principale (en bleu) encadre l'idée essentielle ; les subordonnées (en orange, en retrait) apportent des précisions supprimables ou condensables.}
\end{popfigure}

\courssub{La technique du découpage : isoler les propositions}

Avant de réduire, il faut \textbf{voir} la structure de la phrase. La technique du découpage consiste à réécrire la phrase en plaçant chaque proposition sur une ligne, avec un retrait pour les subordonnées.

\begin{methode}[Découper une phrase complexe]
\begin{enumerate}
\item \textbf{Souligner} tous les verbes conjugués (et les participes).
\item \textbf{Encadrer} la proposition principale : c'est celle qui ne dépend de rien.
\item \textbf{Indenter} chaque subordonnée sous le mot qu'elle complète, en indiquant son type (relative, conjonctive, participiale).
\item \textbf{Évaluer} l'apport de chaque subordonnée : décrit-elle ? explique-t-elle ? pose-t-elle une condition ?
\item \textbf{Décider} : supprimer, condenser ou conserver.
\end{enumerate}
\end{methode}

\begin{exemplebox}[Découpage guidé]
Phrase : « Lakunle, qui voulait moderniser le village, proposa à Sidi de l'épouser sans payer le prix de la dot, parce qu'il considérait cette coutume comme arriérée. »

Découpage :
\begin{itemize}
\item [P] Lakunle \textbf{proposa} à Sidi de l'épouser sans payer le prix de la dot \hfill (principale)
\item \hspace{0.6cm} qui \textbf{voulait} moderniser le village \hfill (relative, descriptive)
\item \hspace{0.6cm} parce qu'il \textbf{considérait} cette coutume comme arriérée \hfill (conjonctive de cause)
\end{itemize}

Analyse : la relative est supprimable. La subordonnée de cause, en revanche, explique \emph{pourquoi} Lakunle refuse la dot : c'est une nuance argumentative. On la condense plutôt qu'on ne la supprime : « Lakunle proposa à Sidi un mariage sans dot, jugeant cette coutume arriérée. »
\end{exemplebox}

\courssub{Transformer les subordonnées pour alléger la phrase}

Réduire ne signifie pas toujours supprimer : on peut \textbf{condenser}, c'est-à-dire dire la même chose en moins de mots. L'outil majeur est la transformation d'une subordonnée en \cle{expansion du nom} (adjectif, apposition, complément du nom), ou l'inverse.

\begin{center}
\begin{tabularx}{\linewidth}{|l|X|X|}
\hline
\rowcolor{popblueL} \textbf{Transformation} & \textbf{Avant (lourd)} & \textbf{Après (allégé)} \\
\hline
Relative $\rightarrow$ adjectif & « un chef qui est âgé » & « un chef âgé » \\
\hline
Relative $\rightarrow$ apposition & « Baroka, qui est le chef du village, » & « Baroka, chef du village, » \\
\hline
Relative $\rightarrow$ complément du nom & « la coutume qui concerne la dot » & « la coutume de la dot » \\
\hline
Conjonctive $\rightarrow$ participiale & « parce qu'il jugeait la coutume arriérée » & « jugeant la coutume arriérée » \\
\hline
Conjonctive $\rightarrow$ nom & « parce qu'il était pauvre » & « à cause de sa pauvreté » / « par pauvreté » \\
\hline
\end{tabularx}
\end{center}

\begin{attention}[La nuance essentielle]
L'erreur la plus grave en contraction est de \textbf{supprimer une subordonnée qui contient une nuance essentielle} : une cause (\emph{parce que}), une condition (\emph{si}), une concession (\emph{bien que}), une conséquence (\emph{si bien que}). Dans un texte argumentatif, ces relations logiques \emph{sont} le raisonnement de l'auteur. Règle d'or : on supprime les \textbf{descriptions}, on condense les \textbf{explications}, on conserve les \textbf{articulations logiques}.
\end{attention}

\courssub{Application : réduire une phrase du \emph{Lion et la perle}}

Mettons toute la méthode en œuvre sur une phrase longue inspirée de la pièce de Wole Soyinka.

\begin{exoresolu}[De 40 mots à 15 mots]
\textbf{Énoncé.} Réduisez la phrase suivante à environ quinze mots, sans en trahir le sens : « Baroka, qui était le chef traditionnel du village d'Ilujinle et qui redoutait profondément l'influence de la modernité, s'opposa fermement à l'ouverture de l'école, parce qu'il craignait que l'instruction ne détruise les coutumes ancestrales. » (40 mots)

\tcblower
\textbf{Solution détaillée.}
\begin{enumerate}
\item \textbf{Verbes conjugués} : \emph{était}, \emph{redoutait}, \emph{s'opposa}, \emph{craignait}, \emph{détruise} (subjonctif dans la subordonnée). Cinq propositions.
\item \textbf{Principale} : « Baroka s'opposa fermement à l'ouverture de l'école. » C'est l'idée essentielle.
\item \textbf{Évaluation} : les deux relatives (\emph{qui était le chef...}, \emph{qui redoutait...}) sont descriptives $\rightarrow$ supprimables ou condensables en apposition. La conjonctive de cause (\emph{parce qu'il craignait...}) est argumentative $\rightarrow$ à condenser, non à supprimer.
\item \textbf{Condensation} : « parce qu'il craignait que l'instruction ne détruise les coutumes ancestrales » $\rightarrow$ « craignant que l'instruction ne détruise les coutumes » (participiale) ou « par crainte pour les coutumes ».
\end{enumerate}
\textbf{Résumé (15 mots)} : « Baroka, chef d'Ilujinle, s'opposa à l'école, craignant que l'instruction ne détruise les coutumes. »
\end{exoresolu}

\begin{popfigure}
\begin{center}
\begin{tikzpicture}[
avant/.style={draw=poporange, fill=poporangeL, rounded corners=3pt, text width=10.5cm, align=left, inner sep=6pt, font=\small},
apres/.style={draw=popgreen, fill=popgreenL, rounded corners=3pt, text width=10.5cm, align=left, inner sep=6pt, font=\small},
fleche/.style={-{Stealth[length=3.5mm]}, very thick, poppurple}
]
\node[avant, align=center] (A) at (0,0){\textbf{AVANT — 40 mots}\\ « Baroka, \textcolor{poporange}{qui était le chef traditionnel du village d'Ilujinle} et \textcolor{poporange}{qui redoutait profondément l'influence de la modernité}, s'opposa \textcolor{poporange}{fermement} à l'ouverture de l'école, parce qu'il craignait que l'instruction ne détruise les coutumes \textcolor{poporange}{ancestrales}. »};
\node[apres, align=center] (B) at (0,-3.8){\textbf{APRÈS — 15 mots}\\ « Baroka, chef d'Ilujinle, s'opposa à l'école, craignant que l'instruction ne détruise les coutumes. »};
\draw[fleche] (A.south) -- node[right, font=\small\itshape, align=left, popdark]{1. repérer la principale\\2. évaluer chaque subordonnée\\3. supprimer / condenser} (B.north);
\end{tikzpicture}
\end{center}
\legende{Diagramme avant/après : les passages en orange sont supprimés (descriptions, intensifieurs) ; la subordonnée de cause est condensée en proposition participiale. La phrase passe de 40 à 15 mots sans perdre l'idée essentielle ni le raisonnement.}
\end{popfigure}

\courssub{Entraînement}

\begin{exercice}[Réduire des phrases complexes]
Réduisez chaque phrase à douze mots environ, en justifiant vos choix (suppression ou condensation) :
\begin{enumerate}
\item « Sidi, qui était la plus belle jeune fille du village et qui en était parfaitement consciente, refusa d'épouser Lakunle, parce qu'il ne voulait pas payer la dot. »
\item « Lorsque l'étranger arriva dans le village avec son appareil photographique, les habitants, qui n'avaient jamais vu un tel objet, furent saisis d'étonnement. »
\item « Bien que Baroka fût âgé, il restait redoutable, car sa ruse était connue de tous les habitants d'Ilujinle. »
\end{enumerate}
\end{exercice}

\begin{corrige}[Exercice : Réduire des phrases complexes]
\begin{enumerate}
\item « Sidi refusa d'épouser Lakunle, qui ne voulait pas payer la dot. » (12 mots). Les relatives descriptives (\emph{la plus belle...}, \emph{consciente...}) sont supprimées ; la cause, essentielle au conflit, est conservée sous forme condensée.
\item « L'arrivée de l'étranger et de son appareil étonna les villageois. » (10 mots). La subordonnée de temps est nominalisée (\emph{lorsque l'étranger arriva} $\rightarrow$ \emph{l'arrivée de l'étranger}) ; la relative descriptive est supprimée.
\item « Malgré son âge, Baroka restait redoutable par sa ruse. » (9 mots). La concession (\emph{bien que}) devient \emph{malgré} + nom ; la cause (\emph{car}) est condensée : ces deux articulations logiques devaient être conservées.
\end{enumerate}
\end{corrige}

\begin{aretenir}[L'essentiel de la section]
\begin{itemize}
\item Une \cle{phrase complexe} contient au moins deux verbes conjugués répartis en propositions.
\item La \cle{proposition principale} porte l'idée essentielle ; les subordonnées (relatives, conjonctives, participiales) apportent des précisions.
\item Pour réduire : 1) repérer la principale, 2) évaluer l'apport de chaque subordonnée, 3) supprimer les descriptions, condenser les explications, conserver les articulations logiques (cause, condition, concession).
\item Les transformations (relative $\rightarrow$ adjectif/apposition/complément du nom, conjonctive $\rightarrow$ participiale ou nom) permettent de dire la même chose en moins de mots.
\end{itemize}
\end{aretenir}

Vous savez désormais découper et alléger une phrase complexe. La section suivante élargira cette compétence à l'échelle du texte entier : il s'agira d'appliquer les \textbf{techniques de réduction} (suppression, condensation, généralisation) non plus à une phrase, mais à des paragraphes complets du \emph{Lion et la perle}.

\coursec{Les techniques de réduction du texte}

Dans la section précédente, nous avons vu comment la phrase complexe permet de resserrer les liens logiques et de dire plus avec moins de mots. Cette compétence trouve maintenant son application naturelle : la \cle{contraction de texte}. En effet, reformuler une phrase complexe de manière économique suppose que l'on sache d'abord \emph{quoi} enlever, \emph{quoi} fusionner, \emph{quoi} généraliser. C'est précisément l'objet des techniques de réduction que nous allons étudier : ce sont les quatre gestes fondamentaux du résumeur.

\begin{definition}[La contraction de texte]
La \cle{contraction de texte} (ou résumé) est la réduction d'un texte à une \textbf{fraction imposée} de sa longueur (généralement le quart au Probatoire et au Baccalauréat technique), en en conservant \textbf{toutes les idées essentielles} et en respectant la pensée de l'auteur. Le résumé est rédigé avec les \textbf{propres mots} du candidat, dans un français correct, et son nombre de mots doit être indiqué à la fin.
\end{definition}

Pour passer d'un texte long à un texte court sans trahir son sens, le résumeur dispose de quatre outils complémentaires : la \cle{suppression}, la \cle{condensation}, la \cle{généralisation} et la \cle{transformation}. Étudions-les une à une.

\courssub{La suppression : éliminer ce qui n'est pas indispensable}

\courssub{Intuition et principe}

Imaginez un sac de riz que l'on vanne : on souffle sur les grains pour faire s'envoler la bale (l'enveloppe vide) et ne garder que le grain nourricier. La suppression fait exactement cela avec le texte : elle élimine tout ce qui ne porte pas l'idée essentielle.

\begin{definition}[La suppression]
La \cle{suppression} consiste à éliminer purement et simplement les éléments \textbf{secondaires} du texte : répétitions, digressions, exemples, énumérations, détails descriptifs, figures de style ornementales, citations illustratives. Ce qui reste après suppression doit former un texte encore cohérent et fidèle à la pensée de l'auteur.
\end{definition}

\courssub{Que faut-il supprimer ?}

\begin{itemize}
\item \textbf{Les répétitions} : l'auteur reprend parfois la même idée sous d'autres mots pour insister. On ne garde qu'une seule formulation.
\item \textbf{Les digressions} : anecdotes, souvenirs, parenthèses qui s'écartent du sujet.
\item \textbf{Les exemples} : ils illustrent une idée générale ; c'est l'idée générale qu'il faut garder, pas l'exemple.
\item \textbf{Les énumérations} : une liste de termes proches peut être remplacée par un mot générique (voir la généralisation).
\item \textbf{Les figures de style ornementales} : comparaisons, métaphores, hyperboles qui décorent sans ajouter d'idée.
\item \textbf{Les citations et références} : elles appuient l'argumentation, elles ne la constituent pas.
\end{itemize}

\begin{exemplebox}[Suppression en action]
\textbf{Texte original (47 mots)} : « La corruption, ce véritable cancer qui ronge notre pays comme une termite attaque le bois, freine le développement. Par exemple, à Douala, un commerçant a dû payer des pots-de-vin à cinq services différents, à savoir la mairie, les impôts, la douane, la police et le ministère du Commerce, avant d'ouvrir sa boutique. »

\textbf{Après suppression (12 mots)} : « La corruption freine le développement du pays. »

On a supprimé la comparaison (« cancer », « termite »), l'exemple du commerçant et l'énumération des cinq services. L'idée essentielle demeure intacte.
\end{exemplebox}

\begin{attention}[Piège : supprimer l'essentiel]
Ne confondez jamais \textbf{exemplebox} et \textbf{argument}. L'exemple illustre, l'argument prouve. Si une phrase apporte une raison nouvelle, elle est essentielle et doit être conservée, même si elle est courte. Demandez-vous toujours : « Si j'enlève cette phrase, une idée de l'auteur disparaît-elle ? » Si oui, gardez-la.
\end{attention}

\courssub{La condensation : dire plus avec moins de mots}

\courssub{Intuition et principe}

La suppression enlève des blocs entiers ; la \cle{condensation}, elle, travaille à l'intérieur de ce qui reste. C'est l'art de \textbf{fusionner} plusieurs phrases ou plusieurs mots en une formulation plus brève, sans perdre aucune information essentielle. C'est ici que la phrase complexe, étudiée dans la section précédente, devient un outil précieux : une subordonnée peut remplacer toute une phrase.

\begin{definition}[La condensation]
La \cle{condensation} consiste à fusionner plusieurs énoncés en un seul, ou à remplacer un groupe de mots par un mot unique de sens équivalent. Elle s'appuie sur trois procédés principaux : la \textbf{nominalisation}, la \textbf{subordination} et le \textbf{remplacement lexical} (un mot pour plusieurs).
\end{definition}

\courssub{Les trois procédés de condensation}

\textbf{1. La nominalisation.} On transforme un verbe (ou une phrase entière) en nom : « le gouvernement décida de construire des écoles » devient « la décision gouvernementale de construction d'écoles ». Une phrase complète se réduit à un groupe nominal.

\textbf{2. La subordination.} Deux phrases juxtaposées deviennent une phrase complexe : « Les pluies furent abondantes. Les récoltes furent bonnes. » devient « Les pluies abondantes ont donné de bonnes récoltes. »

\textbf{3. Le remplacement lexical.} Un mot unique remplace une périphrase : « les personnes qui ne savent ni lire ni écrire » devient « les analphabètes ».

\begin{savaistu}[La nominalisation, reine de l'administration]
La nominalisation (« il décida de construire » $\rightarrow$ « la décision de construire ») est l'outil de condensation préféré des rédacteurs administratifs camerounais. Lisez un arrêté ou un communiqué du ministère : vous y trouverez partout « la mise en œuvre », « l'opérationnalisation », « la sensibilisation des populations ». Ce style nominal permet de condenser des actions entières en quelques mots. Attention toutefois : au résumé, une nominalisation excessive rend le texte lourd et abstrait. Dosez-la !
\end{savaistu}

\begin{exemplebox}[Condensation en action]
\textbf{Avant (28 mots)} : « Les jeunes quittent les villages. Ils partent vers les villes parce qu'ils espèrent y trouver du travail et une vie meilleure. »

\textbf{Après (11 mots)} : « L'exode rural est motivé par la recherche d'emploi et d'une vie meilleure. »

Deux phrases condensées en une seule grâce à la nominalisation (« quittent les villages » $\rightarrow$ « l'exode rural ») et à la subordination implicite de cause.
\end{exemplebox}

\courssub{La généralisation : du particulier au général}

\courssub{Intuition et principe}

Vous avez étudié dans cette unité les relations d'\cle{hypéronymie} et d'\cle{hyponymie} : « fruit » est l'hyperonyme de « mangue », « banane », « ananas ». La généralisation exploite exactement cette relation : au lieu de citer tous les hyponymes, on remonte à l'hyperonyme.

\begin{definition}[La généralisation]
La \cle{généralisation} consiste à remplacer un terme particulier, un exemple ou une énumération par le terme général (l'\textbf{hyperonyme}) qui les englobe. Le cas particulier cède la place à la catégorie ; l'exemple cède la place à l'idée qu'il illustre.
\end{definition}

\begin{exemplebox}[Généralisation en action]
\textbf{Avant (19 mots)} : « Les élèves pratiquent le football, le handball, le basketball, l'athlétisme et le volley-ball dans les lycées camerounais. »

\textbf{Après (8 mots)} : « Les élèves pratiquent divers sports dans les lycées camerounais. »

L'énumération des cinq sports (hyponymes) est remplacée par l'hyperonyme « divers sports ».
\end{exemplebox}

\begin{attention}[Piège : la généralisation excessive]
Généraliser ne signifie pas \textbf{vaguer}. Si l'auteur affirme que « le paludisme tue des milliers d'enfants chaque année au Cameroun », ne résumez pas par « les maladies posent des problèmes » : vous perdez une information essentielle (laquelle ? quelle ampleur ? où ?). La bonne généralisation conserve la précision de la pensée : « le paludisme reste meurtrier au Cameroun ».
\end{attention}

\courssub{La transformation : changer la forme sans changer le sens}

\courssub{Intuition et principe}

La \cle{transformation} est la technique la plus délicate : il s'agit de réécrire un énoncé sous une autre forme grammaticale tout en conservant exactement son sens. Elle prouve à l'examinateur que vous avez réellement compris le texte, puisque vous êtes capable de le redire autrement.

\begin{definition}[La transformation]
La \cle{transformation} consiste à modifier la structure grammaticale d'un énoncé sans en altérer le sens. Les principales transformations sont : la \textbf{nominalisation} (verbe $\rightarrow$ nom), le changement de \textbf{voix} (active $\leftrightarrow$ passive), le changement de \textbf{type de phrase} (négation $\rightarrow$ affirmation équivalente, question rhétorique $\rightarrow$ affirmation), et le remplacement d'une \textbf{citation au discours direct} par une proposition au discours indirect.
\end{definition}

\begin{center}
\begin{tabularx}{\linewidth}{|l|X|X|X|}
\hline
\rowcolor{popblueL} \textbf{Technique} & \textbf{Définition} & \textbf{Avant} & \textbf{Après} \\
\hline
Suppression & Éliminer le secondaire (exemples, répétitions, figures) & « Ce fléau, ce cancer, cette plaie détruit la jeunesse. » & « Ce fléau détruit la jeunesse. » \\
\hline
Condensation & Fusionner plusieurs mots ou phrases en une formulation brève & « Il lutte contre l'ignorance et il lutte contre la pauvreté. » & « Il lutte contre l'ignorance et la pauvreté. » \\
\hline
Généralisation & Remplacer le particulier par le général (hyperonyme) & « mangues, bananes, ananas, papayes » & « fruits » \\
\hline
Transformation & Changer la forme sans changer le sens & « Le maire a lancé les travaux. » & « Le lancement des travaux par le maire... » \\
\hline
\end{tabularx}
\end{center}

\begin{exemplebox}[Transformations combinées]
\textbf{Avant (22 mots)} : « N'est-il pas vrai que les populations ont été sensibilisées par les autorités sanitaires sur les dangers du choléra ? »

\textbf{Après (11 mots)} : « Les autorités sanitaires ont sensibilisé les populations aux dangers du choléra. »

La question rhétorique devient une affirmation (changement de type de phrase) et la voix passive devient active : le sens est strictement identique.
\end{exemplebox}

\courssub{Ce qu'il ne faut jamais faire}

\begin{attention}[Les trois interdits du résumé]
\begin{enumerate}
\item \textbf{Ne jamais copier-coller des phrases entières du texte.} Un résumé qui juxtapose des phrases recopiées n'est pas un résumé : c'est un \emph{faux résumé}, lourdement pénalisé à l'examen. Vous devez reformuler avec vos propres mots.
\item \textbf{Ne jamais ajouter ses opinions personnelles.} Le résumé restitue la pensée de l'auteur, pas la vôtre. Les formules comme « je pense que », « à mon avis », « l'auteur a raison » sont interdites.
\item \textbf{Ne jamais déformer la pensée de l'auteur.} Ni exagération, ni atténuation, ni inversion de la thèse. Si l'auteur est nuancé, votre résumé doit être nuancé.
\end{enumerate}
\end{attention}

\courssub{La méthode complète en cinq étapes}

\begin{methode}[Réduire un texte en cinq étapes]
\begin{enumerate}
\item \textbf{Lire} le texte deux ou trois fois pour bien le comprendre, en repérant le thème et la thèse de l'auteur.
\item \textbf{Dégager la structure argumentative} : repérer les grandes parties du raisonnement (thèse, arguments, illustrations, conclusion).
\item \textbf{Souligner l'essentiel} : marquer les phrases-clés (idées principales, connecteurs logiques, thèse) et écarter mentalement exemples, répétitions et digressions.
\item \textbf{Réduire} en appliquant les quatre techniques : suppression, condensation, généralisation, transformation.
\item \textbf{Reformuler et compter} : rédiger le résumé dans un français correct et cohérent, compter les mots, vérifier la marge autorisée et inscrire le nombre de mots à la fin.
\end{enumerate}
\end{methode}

\courssub{Le comptage des mots et la marge de l'examen}

À l'examen, la longueur du résumé est imposée : en général le \textbf{quart} du texte, avec une marge de tolérance de $\pm 10\,\%$. Tout dépassement est sanctionné, car résumer, c'est aussi respecter une contrainte.

\begin{exemplebox}[Un exemple chiffré]
Pour un texte de \textbf{480 mots}, le résumé attendu fait le quart, soit \textbf{120 mots}. La fourchette acceptée est de $120 \pm 12$, c'est-à-dire de \textbf{108 à 132 mots}. En dessous de 108 mots, vous avez probablement sacrifié une idée essentielle ; au-dessus de 132, vous n'avez pas assez réduit.
\end{exemplebox}

\begin{aretenir}[Compter les mots : les règles]
Chaque mot séparé par une espace compte pour un. Les mots composés à trait d'union (« porte-monnaie ») comptent pour \textbf{un} mot. Les contractions (« l'école », « d'accord ») comptent pour \textbf{un} mot. Entraînez-vous à compter ligne par ligne : comptez les mots d'une ligne moyenne, multipliez par le nombre de lignes, puis ajustez.
\end{aretenir}

\courssub{Entraînement gradué : de 400 mots à 100 mots}

Appliquons maintenant la méthode à un texte de presse camerounaise, en réduisant progressivement : 400 mots $\rightarrow$ 250 $\rightarrow$ 150 $\rightarrow$ 100 mots. Ce graduage montre que la réduction se fait par paliers : à chaque étape, on serre un peu plus la vis.

\begin{popfigure}
\begin{center}
\begin{tikzpicture}
\begin{axis}[
    xbar,
    width=\linewidth,
    height=6.5cm,
    xmin=0, xmax=450,
    symbolic y coords={Résumé final (100), Étape 3 (150), Étape 2 (250), Texte original (400)},
    ytick=data,
    xlabel={Nombre de mots},
    xlabel style={font=\small},
    y tick label style={font=\small},
    x tick label style={font=\small},
    bar width=14pt,
    nodes near coords,
    nodes near coords style={font=\small\bfseries},
    every axis plot/.append style={fill=popblue!70, draw=popblue},
    axis lines=left,
    enlarge y limits=0.2
]
\addplot coordinates {(400,{Texte original (400)}) (250,{Étape 2 (250)}) (150,{Étape 3 (150)}) (100,{Résumé final (100)})};
\end{axis}
\end{tikzpicture}
\end{center}
\legende{Réduction progressive d'un texte de presse : 400 mots $\rightarrow$ 100 mots (le quart), en trois paliers successifs de suppression, condensation et généralisation.}
\end{popfigure}

\begin{exoresolu}[Réduire un passage de presse camerounaise]
\textbf{Énoncé.} Réduisez le passage suivant (environ 100 mots) à environ 25 mots.

« Le marché central de Yaoundé, ce vaste ventre bruyant et coloré de la capitale, ce fourmillement incessant de vendeuses, de crieurs et d'acheteurs pressés, connaît depuis quelques années une véritable métamorphose. Les autorités municipales, soucieuses d'assainir le cadre de vie des commerçants et des clients, ont en effet entrepris de vastes travaux de réhabilitation : réfection des étals, modernisation des abords, amélioration de l'assainissement et réorganisation des différents secteurs de vente. Par exemple, le secteur des vivres frais, autrefois insalubre, a été entièrement reconstruit. »

\tcblower

\textbf{Solution détaillée.}

\textbf{Étape 1 — Lecture.} Thème : la réhabilitation du marché central de Yaoundé. Thèse : les autorités modernisent le marché.

\textbf{Étape 2 — Structure.} Idée 1 : le marché se transforme. Idée 2 : les autorités mènent des travaux de réhabilitation. Le reste est ornement ou exemple.

\textbf{Étape 3 — Suppression.} On élimine : les images (« vaste ventre bruyant et coloré », « fourmillement »), l'énumération (« vendeuses, crieurs, acheteurs »), la précision « soucieuses d'assainir... » (digression de justification) et l'exemple du secteur des vivres frais.

\textbf{Étape 4 — Condensation et généralisation.} L'énumération des travaux (« réfection des étals, modernisation des abords, amélioration de l'assainissement, réorganisation des secteurs ») se condense en « d'importants travaux de réhabilitation ».

\textbf{Étape 5 — Reformulation et comptage.}

« Le marché central de Yaoundé connaît une profonde transformation grâce aux importants travaux de réhabilitation entrepris par les autorités municipales. » (23 mots — dans la fourchette $25 \pm 10\,\%$, soit 23 à 27 mots.)
\end{exoresolu}

\begin{exercice}[À vous de jouer]
Réduisez à environ 30 mots le passage suivant (environ 120 mots) :

« Dans les villages de la région du Nord, les femmes jouent un rôle absolument capital, un rôle décisif, un rôle que l'on ne saurait surestimer, dans l'économie familiale. Elles cultivent le mil, le sorgho, le niébé, l'arachide et le maïs ; elles vendent leurs récoltes au marché hebdomadaire ; elles transforment artisanalement les produits agricoles en farines ou en huiles. Grâce à ces activités inlassables, elles assurent, jour après jour, la subsistance quotidienne de toute la famille, paient souvent la scolarité des enfants et font face aux dépenses de santé. »

Indiquez les techniques utilisées et comptez vos mots.
\end{exercice}

\begin{corrige}[Exercice]
\textbf{Éléments de correction.} Suppression des répétitions (« un rôle capital, décisif, que l'on ne saurait surestimer »), de la figure d'insistance (« inlassables », « jour après jour »). Généralisation de l'énumération des cultures (« mil, sorgho, niébé, arachide, maïs » $\rightarrow$ « diverses cultures vivrières »). Condensation des trois propositions sur les activités en une seule.

\textbf{Résumé possible (30 mots)} : « Dans le Nord, les femmes jouent un rôle essentiel dans l'économie familiale : leurs activités agricoles, commerciales et de transformation assurent la subsistance, la scolarisation des enfants et les soins de santé. » (30 mots)

Toute formulation équivalente, fidèle à la pensée de l'auteur et comprise entre 27 et 33 mots, est acceptable.
\end{corrige}

\begin{aretenir}[L'essentiel de la section]
Quatre techniques permettent de réduire un texte : la \cle{suppression} (éliminer le secondaire), la \cle{condensation} (fusionner et nominaliser), la \cle{généralisation} (remonter de l'hyponyme à l'hyperonyme) et la \cle{transformation} (changer la forme sans changer le sens). Ces techniques s'appliquent dans une méthode en cinq étapes : lire, dégager la structure, souligner l'essentiel, réduire, reformuler et compter. Le résumé final doit respecter la fraction imposée (le quart) avec une marge de $\pm 10\,\%$, être rédigé avec vos propres mots et rester parfaitement fidèle à la pensée de l'auteur.
\end{aretenir}

\coursec{Reformuler les idées essentielles d'un texte en vue d'un résumé}

\courssub{Registres de langue et communication}

La maîtrise des registres de langue est la condition \emph{sine qua non} d'une reformulation réussie. Au Probatoire comme en milieu professionnel, le registre \cle{soutenu} (ou \cle{formel/technique}) s'impose pour le résumé et le compte-rendu.

\begin{definition}[Registre de langue]
Niveau de langue adapté à une situation de communication donnée (émetteur, récepteur, objet, canal). Trois registres principaux : \textbf{soutenu} (écrit littéraire, administratif, technique, examen), \textbf{courant} (échange quotidien, email professionnel standard), \textbf{familier} (oral intime, proscrit à l'écrit d'examen).
\end{definition}

\begin{propriete}[Marqueurs du registre soutenu à l'écrit technique]
\begin{itemize}
    \item \textbf{Lexique :} Précis, désambiguïsé, vocabulaire spécialisé (ex: \emph{spoliation} vs \emph{vol}, \emph{optimiser} vs \emph{améliorer}, \emph{pérenniser} vs \emph{faire durer}).
    \item \textbf{Syntaxe :} Phrases complexes (subordination, relatives, participiales), nominalisation (verbe $\rightarrow$ nom : \emph{révolutionner} $\rightarrow$ \emph{la révolution}), passif/impersonnel (\emph{il convient de}, \emph{on observe que} $\rightarrow$ \emph{la nécessité de}).
    \item \textbf{Morphologie :} Peu d'élision, accords rigoureux, négation complète (\emph{ne... pas}, \emph{ne... point}).
    \item \textbf{Ponctuation :} Structurante (deux-points, point-virgule, parenthèses) pour hiérarchiser l'information.
\end{itemize}
\end{propriete}

\begin{attention}[Piège : le registre « courant » déguisé]
Évitez les tournures orales : \emph{« C'est pour ça que »} $\rightarrow$ \emph{« C'est pourquoi / Par conséquent »} ; \emph{« Il y a »} $\rightarrow$ \emph{« Il existe / On observe »} ; \emph{« Faire en sorte que »} $\rightarrow$ \emph{« Veiller à / Garantir »}.
\end{attention}

\begin{exemplebox}[Adaptation au registre soutenu]
\textbf{Courant :} « Les chefs de chantier n'ont pas assez d'argent pour acheter les drones. Du coup, les travaux prennent du retard. »
\textbf{Soutenu/Technique :} « L'insuffisance budgétaire des chefs de chantier entrave l'acquisition de drones, ce qui engendre des retards sur les chantiers. »
\end{exemplebox}

\courssub{La structure argumentative du texte}

Avant toute réduction, il faut \textbf{disséquer l'architecture logique} du texte. Un texte argumentatif vise à faire adhérer à une thèse.

\begin{definition}[Structure argumentative type]
\begin{itemize}
    \item \textbf{Thèse (Proposition centrale) :} L'idée que l'auteur défend (souvent en introduction ou conclusion).
    \item \textbf{Arguments (Preuves) :} Raisons logiques, faits, chiffres, exemples, citations d'autorité, analogies.
    \item \textbf{Connecteurs logiques :} Balises du raisonnement (cause, conséquence, opposition, concession, ajout).
    \item \textbf{Plan :} Dialectique (Thèse/Antithèse/Synthèse), Analytique (Causes/Conséquences/Solutions), Thématique (Aspects 1, 2, 3).
\end{itemize}
\end{definition}

\begin{methode}[Repérer la structure argumentative (3 lectures)]
\begin{enumerate}
    \item \textbf{Lecture 1 (Globale) :} Identifier le thème, la thèse, le ton, le type de texte.
    \item \textbf{Lecture 2 (Détaillée) :} Souligner les connecteurs, numéroter les paragraphes, noter l'idée-force de chacun en marge.
    \item \textbf{Lecture 3 (Structurale) :} Construire le plan détaillé (Thèse $\rightarrow$ Arg 1 $\rightarrow$ Arg 2 $\rightarrow$ ... $\rightarrow$ Conclusion/Proposition).
\end{enumerate}
\end{methode}

\begin{exemplebox}[Analyse structurelle sur « Le Lion et la Perle » (Extrait : Lakunle vs Baroka)]
\textbf{Thèse :} La tradition (Baroka) l'emporte sur la modernité superficielle (Lakunle) par la ruse et l'ancrage culturel.
\textbf{Arg 1 :} Lakunle incarne une modernité stérile (vocabulaire occidentalisé, rejet du prix de la mariée).
\textbf{Arg 2 :} Baroka utilise la tradition comme levier de pouvoir (ruse de l'impuissance, connaissance des coutumes).
\textbf{Arg 3 :} Sidi, prix de l'enjeu, choisit l'assurance foncière et symbolique de la tradition.
\textbf{Connecteurs clés :} \emph{Toutefois, Certes, En définitive, Paradoxalement}.
\end{exemplebox}

\courssub{Champ lexical, champ sémantique et relations de sens}

Ces outils lexicologiques sont le carburant de la \cle{généralisation} (technique de réduction) et de la \cle{reformulation} (synonymie, hyperonymie).

\begin{definition}[Distinctions fondamentales]
\begin{itemize}
    \item \textbf{Champ lexical :} Ensemble des mots (différents sens) liés à un même \textbf{thème/réalité concrète} (ex: champ lexical du BTP : \emph{dalle, ferraillage, coffrage, grue, béton, maçon}).
    \item \textbf{Champ sémantique :} Ensemble des sens (acceptions) d'un \textbf{même mot} (polysémie) (ex: \emph{volume} : sonorité / grandeur géométrique / tome d'ouvrage / quantité de données).
    \item \textbf{Paronymie :} Ressemblance formelle (son/graphie) entre mots de sens différents (ex: \emph{inférer / inférerer} -- attention \emph{inférer} = déduire ; \emph{prospectif / perspective}).
    \item \textbf{Hyperonymie / Hyponymie :} Relation \textbf{générique / spécifique}. L'hyperonyme (genre) englobe les hyponymes (espèces). \emph{Arbre} (hyperonyme) $\rightarrow$ \emph{Acacia, Balanites, Baobab} (hyponymes). \textbf{Outil clé de la généralisation.}
\end{itemize}
\end{definition}

\begin{propriete}[Application à la contraction]
\begin{itemize}
    \item \textbf{Généraliser une énumération :} Remplacer la liste d'hyponymes par l'hyperonyme (+ quantificateur si besoin). \emph{« Acacias, balanites, baobabs »} $\rightarrow$ \emph{« essences locales / arbres sahéliens »}.
    \item \textbf{Éviter la paronymie :} Ne pas confondre \emph{inciter / insiter}, \emph{préjudice / préjudiciable}, \emph{option / optation}. Vérifier le sens précis au dictionnaire.
    \item \textbf{Choisir le mot du champ lexical technique :} \emph{« Terre »} $\rightarrow$ \emph{« Foncier / Parcelle / Assise foncière »} selon le contexte (juridique, agricole, topographique).
\end{itemize}
\end{propriete}

\begin{exercice}[Entraînement lexical rapide]
Remplacer les énumérations par l'hyperonyme technique approprié :
\begin{enumerate}
    \item Béton, acier, bois, granulats $\rightarrow$ \dots
    \item Géomètre, architecte, ingénieur structure, économiste $\rightarrow$ \dots
    \item Forage, puits, adduction, pompage $\rightarrow$ \dots
\end{enumerate}
\end{exercice}

\begin{corrige}[Exercice lexical]
\begin{enumerate}
    \item Matériaux de construction / Intrants.
    \item Maîtrise d'œuvre / Acteurs de la conception.
    \item Ouvrages hydrauliques / Système d'approvisionnement en eau.
\end{enumerate}
\end{corrige}

\courssub{La phrase complexe au service de la reformulation}

La phrase complexe est le levier syntaxique n°1 de la \cle{condensation} (intégration) et de la \cle{reformulation}. Elle permet d'emboîter les idées (Cause $\rightarrow$ Conséquence, Moyen $\rightarrow$ But, Concession $\rightarrow$ Thèse).

\begin{propriete}[Transformations syntaxiques productives]
\begin{center}
\begin{tabularx}{\linewidth}{|l|X|X|}\hline
\textbf{Opération} & \textbf{Source (souvent phrases simples/juxtaposées)} & \textbf{Cible (Phrase complexe condensée)} \\ \hline
\textbf{Subordination} (Relative, Conjonctive) & « Le BTP pollue. Il doit verdir. » & « Le BTP, \textbf{qui} pollue, \textbf{doit} verdir. » / « \textbf{Parce qu'}il pollue, le BTP doit verdir. » \\ \hline
\textbf{Nominalisation} (Verbe/Adj $\rightarrow$ Nom) & « Les coûts augmentent. Cela freine l'adoption. » & « \textbf{L'augmentation des coûts freine} l'adoption. » \\ \hline
\textbf{Verbalisation} (Nom $\rightarrow$ Verbe) & « La mise en œuvre du projet est effective. » & « Le projet \textbf{se met en œuvre} / \textbf{démarre}. » \\ \hline
\textbf{Passif / Impersonnel} & « Les ingénieurs conçoivent le plan. » & « Le plan \textbf{est conçu par} les ingénieurs. » / « \textbf{La conception du plan} incombe aux ingénieurs. » \\ \hline
\textbf{Participiale / Gérondif} (Antériorité/Simultanéité/Moyen) & « Il a analysé le sol. Puis il a bâti. » & « \textbf{Ayant analysé} le sol, il a bâti. » / « \textbf{En analysant} le sol, il définit les fondations. » \\ \hline
\end{tabularx}
\end{center}
\end{propriete}

\begin{methode}[Entraînement : Transformer pour condenser]
Prenez deux phrases sources liées logiquement. Appliquez \textbf{une} transformation pour n'en faire qu'une.
\begin{enumerate}
    \item Identifier la relation logique (Cause, But, Opposition, Temps).
    \item Choisir l'outil syntaxique adapté (Subordonnant, Participial, Nominalisation).
    \item Réécrire. Vérifier que le sujet est clair (éviter les sujets flottants au gérondif).
\end{enumerate}
\end{methode}

\begin{attention}[Sujet du gérondif / Participiale]
\emph{« En utilisant des drones, le chantier est surveillé. \} $\rightarrow$ Faux (le chantier n'utilise pas).} \\
\emph{« En utilisant des drones, l'équipe surveille le chantier. \} $\rightarrow$ Correct.} \\
\emph{« Le chantier est surveillé par drones. \} $\rightarrow$ Solution nominale/passive préférable.}
\end{attention}

\courssub{Les techniques de réduction du texte}

La réduction (étape préalable à la reformulation) vise à passer du \textbf{texte source} (ex: 300 mots) à une \textbf{structure allégée} (ex: 100 mots) ne gardant que l'ossature argumentative.

\begin{propriete}[Les 3 techniques officielles de réduction (Programme MINESEC)]
\begin{enumerate}
    \item \textbf{Suppression (Élagage) :} Retirer ce qui n'est pas argumentatif : exemples illustratifs, détails descriptifs, répétitions, reformulations internes de l'auteur, connecteurs de simple liaison (\emph{aussi, or, mais} $\rightarrow$ gardés seulement s'ils structurent le plan), métadiscours (\emph{« nous allons voir », « notons que »}).
    \item \textbf{Généralisation (Montée en généralité) :} Remplacer une énumération d'hyponymes par l'hyperonyme ; remplacer une définition/explication par le terme technique ; remplacer une phrase illustrative par un adjectif/nom qualificatif.
    \item \textbf{Intégration (Fusion/Condensation) :} Fondre plusieurs phrases en une seule phrase complexe (subordination, nominalisation, participiale) pour lier Cause $\rightarrow$ Conséquence, Moyen $\rightarrow$ Résultat, Idée $\rightarrow$ Illustration.
\end{enumerate}
\end{propriete}

\begin{exemplebox}[Application des 3 techniques sur un paragraphe source (45 mots)]
\textbf{Source :} « Le coût initial des drones est très élevé. Il faut aussi former les pilotes. Cela demande du temps. Ces deux obstacles expliquent pourquoi les PME n'achètent pas de drones. C'est un frein majeur. »
\begin{itemize}
    \item \textit{Suppression :} « très », « aussi », « Cela demande du temps » (détail), « C'est un frein majeur » (répétition).
    \item \textit{Généralisation :} « Coût élevé + Formation » $\rightarrow$ « Coût global / Investissement initial ». « PME n'achètent pas » $\rightarrow$ « Adoption freinée ».
    \item \textit{Intégration :} Fusion Cause (Coût/Formation) $\rightarrow$ Conséquence (Non-adoption).
\end{itemize}
\textbf{Structure réduite (18 mots) :} « L'investissement initial (achat, formation) freine l'adoption des drones par les PME. »
\end{exemplebox}

\courssub{Reformuler les idées essentielles et rédiger le résumé}

Après avoir maîtrisé les techniques de réduction (suppression, généralisation, intégration), nous abordons l'étape finale et décisive : la \cle{reformulation}. Réduire un texte ne suffit pas ; il faut le réécrire avec ses propres mots pour produire un résumé autonome, fluide et fidèle. C'est ici que s'articulent vos compétences en langue (registres, syntaxe, lexique) et en analyse (structure argumentative).

\courssub{Principes de la reformulation : fidélité, clarté, cohérence, style personnel}

\begin{definition}[Reformulation]
Opération qui consiste à \textbf{exprimer les idées d'un texte source avec ses propres mots}, sans en altérer le sens, ni en ajouter, ni en retrancher l'essentiel. Elle vise à produire un \emph{nouveau texte} (le résumé) plus court, mais équivalent sur le plan informationnel.
\end{definition}

La reformulation repose sur quatre piliers indissociables :
\begin{enumerate}
    \item \textbf{Fidélité au sens :} Respect strict de la pensée de l'auteur. Pas d'interprétation, pas d'ajout d'idées personnelles.
    \item \textbf{Clarté :} Le résumé doit être compréhensible par quelqu'un qui n'a pas lu le texte original.
    \item \textbf{Cohérence :} Les idées s'enchaînent logiquement (connecteurs, progression thématique), ce n'est pas une liste de phrases disjointes.
    \item \textbf{Style personnel :} Vous écrivez \emph{votre} phrase (vocabulaire, syntaxe), pas celle de l'auteur. C'est la preuve de votre compréhension.
\end{enumerate}

\begin{attention}[Piège majeur : le « patchwork » ou le « copier-coller »]
Ne coupez pas des morceaux du texte original pour les coller bout à bout, même en changeant quelques mots. C'est du \emph{plagiat} ou de la \emph{citation}, pas un résumé. Au Probatoire, un résumé qui conserve trop de séquences textuelles identiques à l'original est sanctionné (note de langue faible, hors-sujet partiel).
\end{attention}

\courssub{Procédés de reformulation : outillage linguistique}

Pour passer du texte source à votre phrase, vous mobilisez trois leviers principaux, souvent combinés.

\begin{propriete}[Les trois leviers de la reformulation]
\begin{itemize}
    \item \textbf{1. Synonymie et paraphrase lexicale :} Remplacer un mot ou un groupe de mots par un équivalent (synonyme, hyperonyme, circonlocution). Utiliser le champ lexical technique.
    \item \textbf{2. Changement de construction syntaxique :} Modifier la structure de la phrase (voix active $\leftrightarrow$ passive, subordination $\leftrightarrow$ coordination, nominalisation $\leftrightarrow$ verbalisation, changement de catégorie grammaticale, gérondif/participiale).
    \item \textbf{3. Regroupement et hiérarchisation des idées :} Fusionner plusieurs phrases sources en une seule phrase cible ; subordonner le détail à l'essentiel (intégration).
\end{itemize}
\end{propriete}

\begin{exemplebox}[Mise en œuvre combinée des trois leviers]
\textbf{Texte source :} « L'urbanisation galopante de Douala engendre une insécurité foncière croissante. Les populations autochtones se voient dépossédées de leurs terres ancestrales par des promoteurs sans scrupules. »

\textbf{Analyse des procédés :}
\begin{itemize}
    \item \textit{Synonymie/Paraphrase :} « galopante » $\rightarrow$ « rapide » ; « engendre » $\rightarrow$ « provoque » ; « insécurité foncière » $\rightarrow$ « conflits de propriété » ; « autochtones » $\rightarrow$ « originaires » ; « dépossédées » $\rightarrow$ « spoliées ».
    \item \textit{Syntaxe :} Deux phrases $\rightarrow$ Une phrase complexe (subordination de conséquence). Nominalisation : « L'urbanisation... engendre » $\rightarrow$ « La croissance rapide... provoque ».
    \item \textit{Regroupement :} Cause (urbanisation) + Conséquence (insécurité/spoliation) fusionnées.
\end{itemize}

\textbf{Reformulation proposée :} « La croissance rapide de Douala provoque de plus en plus de conflits de propriété, spoliant les populations originaires de leurs terres ancestrales. »
\end{exemplebox}

\begin{methode}[Démarche pas à pas pour reformuler une idée]
\begin{enumerate}
    \item \textbf{Isoler l'idée :} Souligner le noyau informationnel (Sujet + Verbe + Complément essentiel) dans la phrase source.
    \item \textbf{Dégager le sens :} Reformuler mentalement l'idée en français courant, sans regarder le texte.
    \item \textbf{Choisir le levier :} Décider si on change le vocabulaire, la syntaxe, ou les deux.
    \item \textbf{Rédiger la phrase cible :} L'écrire sur brouillon.
    \item \textbf{Vérifier l'équivalence :} Relire source et cible : le sens est-il identique ? La nuance (certitude, doute, ironie) est-elle respectée ?
\end{enumerate}
\end{methode}

\courssub{Rédaction du résumé : de la structure à la phrase finale}

Le résumé n'est pas une suite de reformulations isolées ; c'est un \textbf{texte à part entière} qui respecte les normes de la rédaction administrative et technique (impersonnalité, objectivité, connecteurs logiques).

\begin{propriete}[Structure type du résumé (Probatoire / Bac Tech)]
\begin{itemize}
    \item \textbf{Introduction (1 phrase) :} Auteur (si connu), titre (si connu), nature du texte, thème général, thèse principale.
    \item \textbf{Développement (corps) :} Enchaînement des idées forces dans l'ordre logique du texte (ou ordre thématique si plus clair), reliées par des connecteurs adaptés (\emph{donc, par conséquent, en outre, cependant, en effet, tout d'abord, ensuite, enfin}).
    \item \textbf{Conclusion (souvent implicite ou 1 phrase) :} Ouverture, conséquence majeure, ou simple achèvement de la démonstration. Pas d'avis personnel.
\end{itemize}
\end{propriete}

\begin{exemplebox}[Introduction type]
\textit{« Dans son article "Ville et défis fonciers" (2023), l'urbaniste Jean Mbarga analyse les conséquences de l'expansion urbaine de Douala sur la sécurité des terres. Il démontre que la spéculation foncière, favorisée par un cadre juridique lacunaire, marginalise les populations autochtones. »}
\end{exemplebox}

\begin{aretenir}[Connecteurs logiques au service de la cohérence]
Le choix du connecteur reflète la relation logique entre les idées \textbf{du texte original} :
\begin{center}
\begin{tabularx}{\linewidth}{|l|X|}\hline
\textbf{Relation logique} & \textbf{Connecteurs privilégiés} \\ \hline
Cause / Conséquence & \emph{Par conséquent, ainsi, c'est pourquoi, d'où, du fait que} \\ \hline
Opposition / Concession & \emph{Cependant, néanmoins, pourtant, bien que, si... toutefois} \\ \hline
Ajout / Énumération & \emph{En outre, de plus, par ailleurs, d'une part... d'autre part} \\ \hline
Illustration / Explication & \emph{En effet, notamment, c'est-à-dire, par exemple} \\ \hline
Conclusion / Résumé & \emph{En définitive, au total, finalement, pour conclure} \\ \hline
\end{tabularx}
\end{center}
\end{aretenir}

\courssub{Vérification finale : le test des 4 questions}

Avant de rendre votre copie, appliquez systématiquement cette grille de contrôle. C'est votre assurance qualité.

\begin{methode}[Test des 4 questions : validation du résumé]
\begin{enumerate}
    \item \textbf{Fidèle ?} Chaque idée du résumé se trouve-t-elle dans le texte source ? N'ai-je rien inventé ? N'ai-je pas déformé une nuance (ex: « probable » devenu « certain ») ?
    \item \textbf{Complet sur l'essentiel ?} Les 3 à 5 idées forces (axes du plan argumentatif) sont-elles toutes présentes ? Ai-je éliminé les exemples, détails, répétitions ?
    \item \textbf{Bonne longueur ?} Respect de la consigne (ex: $\pm 10\%$ de la longueur demandée). \textbf{Compter les mots} (mots-outils inclus). Dépasser la marge = perte de points.
    \item \textbf{Bien écrit ?} Orthographe, grammaire, syntaxe, ponctuation, registres de langue (soutenu/technique), impersonnalité (pas de « je », « nous », « on » sauf citation). Fluidité de la lecture.
\end{enumerate}
\end{methode}

\begin{savaistu}[Au Probatoire / Bac Technique : barème type]
Le résumé est souvent noté sur \textbf{20 points} répartis ainsi (variable selon académies) :
\begin{itemize}
    \item \textbf{Contenu / Fidélité (8 à 10 pts) :} Idées justes, essentielles, non déformées.
    \item \textbf{Langue / Expression (6 à 8 pts) :} Correction grammaticale, richesse lexicale, syntaxe variée (phrases complexes), registres adaptés.
    \item \textbf{Respect des consignes / Forme (4 à 6 pts) :} Nombre de mots, présentation, introduction/conclusion, connecteurs.
\end{itemize}
\textbf{Astuce :} Visez la concision. Une phrase complexe bien construite vaut mieux que deux phrases simples. Utilisez la phrase complexe (subordination, relatives, participiales) apprise en section 4.
\end{savaistu}

\courssub{Le compte-rendu : objet professionnel et technique}

En milieu technique et professionnel (BTS, DUT, entreprise), le \cle{compte-rendu} est l'application directe de la compétence « résumer ». Il diffère du résumé scolaire par sa finalité \textbf{opérationnelle}.

\begin{definition}[Compte-rendu]
Document écrit, objectif et structuré, qui relate de manière synthétique le déroulement et les décisions d'une réunion, d'une visite de chantier, d'une expérience de laboratoire ou d'une lecture technique, à destination d'un destinataire précis (hiérarchie, client, équipe).
\end{definition}

\begin{center}
\begin{tabularx}{\linewidth}{|l|X|X|}\hline
\textbf{Critère} & \textbf{Résumé scolaire (examen)} & \textbf{Compte-rendu professionnel} \\ \hline
\textbf{Source} & Un texte unique, littéraire ou argumentatif & Réunion, visite, expérience, dossier technique \\ \hline
\textbf{Destinataire} & Correcteur anonyme & Hiérarchie, client, collègues (identifiés) \\ \hline
\textbf{Structure} & Introduction / Développement / Conclusion (fluide) & En-tête (objet, date, participants) / Ordre du jour / Décisions / Actions (Qui ? Quoi ? Quand ?) / Signature \\ \hline
\textbf{Style} & Littéraire, phrases complexes, connecteurs & Télégraphique, puces, verbes d'action, tableaux \\ \hline
\textbf{Subjectivité} & Interdite (neutralité absolue) & Interdite pour les faits ; \emph{avis techniques} identifiés comme tels \\ \hline
\end{tabularx}
\end{center}

\begin{experience}[Atelier : Du résumé au compte-rendu de chantier]
\textbf{Contexte :} Vous êtes technicien géomètre-topographe. Vous avez lu le rapport d'expertise foncière (texte source type « Le Lion et la Perle » : conflit coutumier / droit moderne).
\textbf{Consigne :} Rédiger le compte-rendu pour le chef de projet (150 mots max).
\textbf{Structure imposée :}
\begin{enumerate}
    \item \textbf{Objet :} Compte-rendu lecture rapport expertise lot 45 - Douala.
    \item \textbf{Constat principal :} Conflit limite entre famille N. et promoteur S. (1 phrase).
    \item \textbf{Analyse technique :} Bornage ancien imprécis + occupation anarchique (2 phrases).
    \item \textbf{Recommandation / Decision :} Nouveau bornage contradictoire + médiation chef de quartier (1 phrase).
    \item \textbf{Action :} Vous (Géomètre) : lever topographique J+5. Chef projet : convocation parties J+3.
\end{enumerate}
\textbf{Apport :} Passage du « style essai » au « style fiche technique ». Entraînement à la nominalisation et aux verbes d'action (\emph{prévoir, réaliser, valider, transmettre}).
\end{experience}

\courssub{Remédiation : correction collective et auto-évaluation}

La progression ne s'arrête pas à la production. L'analyse d'erreurs types (remédiation) ancre les acquis.

\begin{exoresolu}[Correction collective d'un résumé type]
\textbf{Sujet :} Contracter le texte ci-dessous à \textbf{60 mots $\pm$ 10\%} (soit 54 à 66 mots).

\textbf{Texte source (85 mots) :}
« L'intelligence artificielle (IA) transforme radicalement le secteur du BTP. Les algorithmes prédictifs optimisent la gestion des chantiers en anticipant les retards. Les drones couplés à la modélisation 3D (BIM) permettent un suivi précis de l'avancement. Cependant, l'investissement initial coûteux et la formation du personnel freinent l'adoption par les PME africaines. Les gouvernements doivent inciter la transition numérique par des subventions ciblées et l'adaptation des cursus de formation professionnelle. »

\textbf{Proposition d'élève (72 mots) - \textcolor{red}{HORS NORMES (trop long)}} :
« L'intelligence artificielle change le BTP. Les algorithmes aident à gérer les chantiers en prévoyant les retards. Les drones et la 3D permettent de suivre les travaux. Mais c'est cher au début et il faut former les ouvriers, ce qui bloque les petites entreprises en Afrique. Les États doivent donner de l'argent et changer les formations pour aider. »

\textbf{Analyse des erreurs (Grille) :}
\begin{itemize}
    \item \textbf{Longueur :} 72 mots $>$ 66 max. \textbf{\textcolor{red}{Échec critère 3.}}
    \item \textbf{Langue :} Registre familier (« change », « c'est cher », « donne de l'argent », « bloque »). Pas de subordination. \textbf{\textcolor{red}{Faible critère 2.}}
    \item \textbf{Fidélité :} Idées présentes mais « algorithmes prédictifs $\rightarrow$ algorithmes » (perte de précision), « PME africaines $\rightarrow$ petites entreprises en Afrique » (acceptable).
\end{itemize}

\textbf{Version corrigée (58 mots) - \textcolor{popgreen}{VALIDE}} :
« L'intelligence artificielle révolutionne le BTP : algorithmes prédictifs et drones couplés au BIM optimisent la planification et le suivi des chantiers. Toutefois, le coût initial élevé et la nécessité de former les équipes entravent son adoption par les PME africaines. Une transition numérique effective suppose des subventions ciblées des gouvernements et une refonte des cursus de formation professionnelle. »

\textbf{Gains :} Nominalisation (« révolutionne », « optimisent », « entravent »), subordination relative (« couplés au BIM »), vocabulaire technique précis (« refonte », « subventions ciblées »), longueur respectée.
\end{exoresolu}

\begin{aretenir}[Grille d'auto-évaluation (à coller dans le cahier)]
Avant de rendre \textbf{tout} résumé ou compte-rendu, cochez :
\begin{center}
\begin{tabular}{|c|p{6cm}|c|}\hline
\textbf{Critère} & \textbf{Question clé} & \textbf{OK ?} \\ \hline
\textbf{1. Fidélité} & Ai-je respecté la pensée de l'auteur sans l'interpréter ? & $\square$ \\ \hline
\textbf{2. Essentiel} & Les 3-4 arguments majeurs sont-ils là ? (Exemples exclus) & $\square$ \\ \hline
\textbf{3. Longueur} & Nombre de mots compté $\in$ [Consigne $-10\%$ ; Consigne $+10\%$] ? & $\square$ \\ \hline
\textbf{4. Langue} & Zéro faute d'accord/orthographe ? Phrases complexes ? Registre soutenu ? Impersonnel ? & $\square$ \\ \hline
\textbf{5. Cohérence} & Connecteurs logiques variés ? Enchaînement fluide ? Introduction claire ? & $\square$ \\ \hline
\end{tabular}
\end{center}
\end{aretenir}

\courssub{Préparation de l'exposé oral : présenter structure et résumé en 5 minutes}

L'oral (épreuve E1 ou contrôle continu) évalue votre capacité à \textbf{structurer votre pensée} et à \textbf{communiquer techniquement}.

\begin{methode}[Canevas de l'exposé oral de 5 minutes (1 minute = 130 mots env.)]
\begin{enumerate}
    \item \textbf{Accroche \& Identification (30 s) :} « Bonjour. Je vais présenter l'analyse du texte "Titre" de Auteur, publié dans Source/Année. Ce texte argumentatif traite de [Thème] et défend la thèse : [Thèse en 1 phrase]. »
    \item \textbf{Structure argumentative (1 min 30 s) :} « Le raisonnement s'articule en [3] étapes : \\ \textbf{1.} [Idée force 1 + argument clé + exemple/autorité]. \\ \textbf{2.} [Idée force 2 + argument clé + nuance/contre-argument levé]. \\ \textbf{3.} [Idée force 3 + argument clé + conséquence/proposition]. \\ \textit{Utilisez : "Dans un premier temps...", "Par ailleurs...", "Enfin..."} »
    \item \textbf{Résumé synthétique (1 min 30 s) :} Lisez/énoncez votre résumé rédigé (version orale, plus fluide). Regardez le jury, pas votre feuille.
    \item \textbf{Apport personnel / Contexte technique (1 min) :} « Ce texte éclaire mon futur métier de [Spécialité] car... / Il s'inscrit dans le débat actuel sur... / Il soulève la question technique de... »
    \item \textbf{Conclusion formelle (15 s) :} « Pour conclure, ce texte démontre que... Je vous remercie de votre attention. Je suis prêt à répondre à vos questions. »
\end{enumerate}
\end{methode}

\begin{attention}[Pièges de l'oral]
\begin{itemize}
    \item \textbf{Lire son résumé tête baissée} : Interdit. Des notes (mots-clés) sont tolérées, pas le texte intégral.
    \item \textbf{Dépasser le temps} : Chronométrez-vous. 5 min = coupure nette.
    \item \textbf{Confondre structure et résumé} : La structure = le squelette (plan). Le résumé = la chair (texte reformulé). Les deux sont demandés.
    \item \textbf{Jargon incompris} : Expliquez les termes techniques (BIM, algorithmique prédictive) en une phrase simple.
\end{itemize}
\end{attention}

\courssub{Synthèse visuelle : la démarche complète de contraction}

\begin{popfigure}
\begin{tikzpicture}[
    node distance=1.2cm and 1.5cm,
    every node/.style={font=\small},
    box/.style={rectangle, draw=popblue, thick, rounded corners, fill=popblueL, text width=3.2cm, align=center, minimum height=1.8cm},
    arrow/.style={-Stealth, thick, popdark},
    decision/.style={diamond, draw=poporange, thick, aspect=2, fill=poporangeL, text width=2.5cm, align=center, inner sep=2pt},
]
% Nodes
\node[box, align=center] (source){Texte Source\\ (Lecture active,\\ repérage thèse/arguments)};
\node[box, below=of source, align=center] (analyse){Analyse Structurelle\\ (Plan détaillé,\\ champs lexicaux,\\ connecteurs logiques)};
\node[box, below=of analyse, align=center] (reduction){Réduction\\ (Suppression détails,\\ Généralisation énumérations,\\ Intégration définitions)};
\node[decision, below=of reduction, align=center] (check1){Longueur\\ intermédiaire\\ respectée ?};
\node[box, below=of check1, align=center] (reform){Reformulation\\ (Synonymie, Syntaxe,\\ Regroupement,\\ Phrase complexe)};
\node[box, below=of reform, align=center] (draft){Rédaction Brouillon\\ (Intro, Développement,\\ Conclusion, Connecteurs)};
\node[decision, below=of draft, align=center] (check2){Test 4 Questions\\ (Fidèle, Essentiel,\\ Longueur, Langue) ?};
\node[box, below=of check2, fill=popgreenL, draw=popgreen, align=center] (final){Résumé Final\\ (Validé, Propre,\\ Nombre de mots exact)};

% Arrows
\draw[arrow] (source) -- (analyse);
\draw[arrow] (analyse) -- (reduction);
\draw[arrow] (reduction) -- (check1);
\draw[arrow] (check1) -- node[left, xshift=-2mm] {\small Non} ++(-2.5,0) |- (reduction);
\draw[arrow] (check1) -- node[right, xshift=2mm] {\small Oui} (reform);
\draw[arrow] (reform) -- (draft);
\draw[arrow] (draft) -- (check2);
\draw[arrow] (check2) -- node[left, xshift=-2mm] {\small Non} ++(-2.5,0) |- (reform);
\draw[arrow] (check2) -- node[right, xshift=2mm] {\small Oui} (final);
\end{tikzpicture}
\legende{Organigramme de la démarche complète de contraction : du texte source au résumé validé (boucles de rétroaction incluses).}
\end{popfigure}

\courssub{Grille d'évaluation type (sur 20 points)}

Cette grille vous permet de simuler la notation de l'examinateur.

\begin{popfigure}
\centering
\begin{tikzpicture}
\begin{axis}[
    ybar stacked,
    bar width=12pt,
    width=\linewidth,
    height=8cm,
    xlabel={Critères d'évaluation},
    ylabel={Points (sur 20)},
    symbolic x coords={Fidélité/Contenu, Langue/Expression, Respect Consignes, Cohérence/Structure},
    xtick=data,
    x tick label style={rotate=30, anchor=east, font=\small},
    ytick={0,2,4,6,8,10,12,14,16,18,20},
    ymin=0, ymax=20,
    legend style={at={(0.5,-0.25)}, anchor=north, legend columns=-1, font=\small},
    nodes near coords={\pgfmathprintnumber\pgfplotspointmeta},
    nodes near coords style={font=\tiny, text=popdark},
]
% Fidélité : 8 pts max
\addplot[fill=popblue!70] coordinates {(Fidélité/Contenu,8)};
% Langue : 7 pts max
\addplot[fill=popgreen!70] coordinates {(Langue/Expression,7)};
% Consignes : 3 pts max
\addplot[fill=poporange!70] coordinates {(Respect Consignes,3)};
% Cohérence : 2 pts max
\addplot[fill=poppurple!70] coordinates {(Cohérence/Structure,2)};

\legend{Fidélité (8), Langue (7), Consignes (3), Cohérence (2)}
\end{axis}
\end{tikzpicture}
\legende{Barème type de notation du résumé au Probatoire / Bac Technique (répartition indicative sur 20 points).}
\end{popfigure}

\begin{exercice}[Entraînement final : Contraction intégrée]
\textbf{Texte support (112 mots) :}
« Le développement durable s'impose comme un impératif pour les pays du Sahel face aux changements climatiques. La désertification progresse de 5 km par an au Nord-Cameroun, détruisant les terres arables et poussant les populations à l'exode rural. La Grande Muraille Verte, initiative panafricaine lancée en 2007, vise à restaurer 100 millions d'hectares de terres dégradées d'ici 2030 par la plantation d'essences locales (acacia, balanites). Cependant, le manque de financement pérenne, la gouvernance foncière instable et l'insécurité terroriste dans la zone sahélienne compromettent la survie des plants. Une approche intégrant gestion de l'eau, agroforesterie et appropriation communautaire est indispensable pour inverser la tendance. »

\textbf{Consignes :}
\begin{enumerate}
    \item Dégagez la structure argumentative (Thèse, 3 arguments, Conclusion/Proposition).
    \item Rédigez un résumé de \textbf{60 mots $\pm$ 10\%} (54 à 66 mots).
    \item Appliquez le Test des 4 questions. Comptez vos mots.
    \item Présentez l'oral (5 min) devant un pair : Structure + Résumé + Lien spécialité.
\end{enumerate}
\end{exercice}

\begin{corrige}[Exercice n°6 - Éléments de correction]
\textbf{1. Structure argumentative :}
\begin{itemize}
    \item \textbf{Thèse :} Développement durable = impératif Sahel / Nord-Cameroun (désertification + exode).
    \item \textbf{Arg 1 (Solution) :} Grande Muraille Verte (objectif 100M ha, essences locales).
    \item \textbf{Arg 2 (Freins) :} Financement, foncier, insécurité $\rightarrow$ mortalité plants.
    \item \textbf{Proposition (Conclusion) :} Approche intégrée (eau, agroforesterie, communauté).
\end{itemize}

\textbf{2. Résumé type (59 mots) :}
« Face à la désertification galopante au Nord-Cameroun, la Grande Muraille Verte ambitionne de restaurer 100 millions d'hectares via des essences locales. Toutefois, l'insuffisance des financements, l'instabilité foncière et l'insécurité sahélienne menacent la pérennité des plantations. Seule une stratégie intégrée, associant gestion hydraulique, agroforesterie et implication des communautés locales, permettra d'enrayer durablement la dégradation des sols. »

\textbf{3. Vérification :} Fidèle (oui), Essentiel (4 axes couverts), Longueur (59 $\in$ [54,66]), Langue (soutenu, complexe, impersonnel, connecteurs : toutefois, seule... permettra).
\end{corrige}

\begin{aretenir}[Check-list finale avant l'examen (Probatoire / Bac)]
\begin{itemize}
    \item \textbf{Je sais} analyser la structure argumentative (thèse, arguments, types d'arguments, connecteurs).
    \item \textbf{Je sais} réduire un texte (3 techniques : suppression, généralisation, intégration).
    \item \textbf{Je sais} reformuler (synonymie, syntaxe, regroupement) sans trahir le sens.
    \item \textbf{Je sais} rédiger un résumé structuré (Intro, Développement connecté, Conclusion) dans la longueur imposée.
    \item \textbf{Je sais} m'auto-évaluer (Grille 4 questions + Grille 20 pts).
    \item \textbf{Je sais} produire un compte-rendu technique (structure administrative).
    \item \textbf{Je sais} présenter l'analyse et le résumé à l'oral en 5 minutes (Canevas structuré).
\end{itemize}
\textbf{La contraction de texte n'est pas un exercice de style, c'est une compétence professionnelle : synthétiser l'information complexe pour décider et agir.}
\end{aretenir}

\newpage

\coursec{Activité d'intégration}

\coursec{Leçon 06 : Reformuler les idées essentielles d'un texte en vue d'un résumé}

\courssub{Objectifs de l'activité}
\noindent Cette activité d'intégration vise à mobiliser l'ensemble des savoirs et savoir-faire travaillés au cours de la leçon n°6 pour produire un résumé conforme aux exigences du Baccalauréat technique. À l'issue de cette séance, l'élève doit être capable de :
\begin{itemize}
    \item Identifier le \cle{registre de langue} d'un texte journalistique d'opinion et en expliquer les effets sur la réception.
    \item Dégager la \cle{structure argumentative} d'un éditorial (thèse, arguments, exemples, concession, conclusion) et la représenter sous forme d'arbre.
    \item Repérer les \cle{champs lexicaux} et \cle{champs sémantiques} dominants pour saisir la cohérence thématique.
    \item Appliquer les trois \cle{techniques de réduction} (suppression, condensation, généralisation) en respectant la \cle{phrase complexe} comme outil de condensation syntaxique.
    \item \cle{Reformuler} les idées essentielles en changeant la formulation (paronymie, hyperonymie/hyponymie, changement de voix, nominalisation) sans trahir le sens.
    \item Rédiger un \cle{résumé} de 125 mots (\textpm 10 \%) respectant la neutralité, la continuité et la cohérence.
    \item S'auto-évaluer à l'aide d'une grille critériée et présenter sa démarche à l'oral.
\end{itemize}

\courssub{Situation}
\noindent \textbf{Contexte :} Vous êtes élèves en classe de Première technique au Lycée Technique de Douala-Bassa. Dans le cadre du module « Communication et expression écrite », votre professeur vous propose de travailler sur un corpus de presse écrite camerounaise pour préparer l'épreuve de contraction de texte du Probatoire.

\noindent \textbf{Support :} L'éditorial ci-dessous, paru dans le quotidien \emph{Cameroon Tribune} (édition du 12 mars 2026, page 2), sous la plume du directeur de publication. Le texte compte \textbf{512 mots}.

\begin{center}
\begin{tabularx}{\linewidth}{|X|}\hline
\rowcolor{popblueL}
\textbf{\textsc{Éditorial : Riz importé vs riz local — Le courage de choisir la souveraineté}} \\ \hline
\emph{Par Jean-Baptiste Atemengue}

Depuis trois décennies, le Cameroun a fait le choix facile de l'importation massive de riz pour nourrir ses villes. En 2023, la facture a dépassé 350 milliards de FCFA, soit près de 1,2 million de tonnes débarquées au port de Douala, pendant que la plaine de Mbo, dans le Moungo, et les périmètres rizicoles de la Bénoué-Nord ne couvrent que 15 \% des besoins nationaux. Ce paradoxe n'est plus tenable. Le gouvernement annonce, dans le cadre de la Stratégie Nationale de Développement 2030 (SND30), une hausse progressive des droits de douane sur le riz importé, couplée à un plan d'urgence de 25 milliards FCFA pour la mécanisation de la plaine de Mbo. Faut-il y voir une provocation protectionniste ou un acte de lucidité patriotique ?

Les tenants du \emph{statu quo} brandissent l'argument du pouvoir d'achat. « Le riz thaïlandais ou vietnamien coûte 350 FCFA le kilo au détail, contre 550 FCFA pour le riz local décortiqué artisanalement », entend-on dans les marchés de Nkolndongo ou de Mfoundi. C'est vrai, mais c'est une vérité tronquée. Ce prix bas cache des coûts externes colossaux : dépendance aux cours mondiaux (la tonne a flambé à 650 \$ en 2024), fuite des devises, précarité des riziculteurs locaux contraints à l'exode rural. Interdire purement et simplement l'importation demain matin serait irresponsable : la production locale ne suit pas, les infrastructures de transformation manquent, et le consommateur urbain paierait la note. Cependant, maintenir les frontières grandes ouvertes sans contrepartie, c'est condamner la plaine de Mbo à rester un potentiel géologique au lieu de devenir un grenier économique.

L'exemple du Nigeria voisin est instructif. Depuis 2019, Abuja a fermé ses frontières terrestres au riz et injecté des milliards de nairas dans le programme \emph{Anchor Borrowers}. Résultat : la production a bondi de 4 à 9 millions de tonnes en quatre ans, créant des millions d'emplois ruraux. Le Cameroun a la terre, l'eau (le Moungo, le Wouri, la Sanaga), la main-d'œuvre jeune. Il lui manque la volonté politique constante et l'ingénierie financière adaptée. Le plan d'urgence annoncé doit donc aller au-delà de la simple distribution de tracteurs. Il faut des variétés améliorées (NERICA), des unités de décorticage industrielles, des routes rurales praticables en saison des pluies, et surtout un contrat de filière garantissant le prix plancher au producteur et le prix plafond au consommateur.

La souveraineté alimentaire ne se décrète pas, elle se construit. Elle exige du temps, de la rigueur dans la gestion des fonds publics — le scandale des 12 milliards FCFA détournés au Ministère de l'Agriculture en 2021 est encore dans les mémoires — et une communication transparente auprès des populations. Augmenter les droits de douane sans accompagner la production, c'est de la fiscalité déguisée ; accompagner la production sans protéger le marché naissant, c'est de la naïveté économique. Le juste milieu est une politique de \emph{substitution progressive} : taxer l'importé pour financer le local, tout en plafonnant le prix de vente du riz camerounais à 400 FCFA le kilo pendant la transition. C'est à ce prix que le riz de Mbo retrouvera la table des Camerounais, non par contrainte, mais par fierté et par goût. L'enjeu dépasse le grain de riz : c'est la dignité d'une nation qui se joue dans ses champs. \\ \hline
\end{tabularx}
\end{center}

\courssub{Consignes}
\noindent Travail en groupes de quatre élèves (20 min), puis travail individuel (30 min), enfin restitution orale (10 min).
\begin{enumerate}
    \item \textbf{Analyse du registre :} Quel est le registre de langue dominant de cet éditorial ? Citez trois indices lexicaux ou syntaxiques qui le caractérisent et expliquez l'effet recherché sur le lectorat.
    \item \textbf{Structure argumentative :} Représentez la structure argumentative du texte sous forme d'arbre (thèse, arguments principaux, arguments secondaires/exemples, concession, conclusion). Utilisez les connecteurs logiques du texte pour justifier votre arborescence.
    \item \textbf{Champs lexicaux et sémantiques :} Identifiez les deux champs lexicaux dominants. Pour chacun, donnez cinq occurrences et précisez la relation sémantique (synonymie, antonymie, hyperonymie/hyponymie) qui lie au moins deux de ces occurrences.
    \item \textbf{Techniques de réduction (brouillon) :} Sur une copie de brouillon, appliquez successivement :
    \begin{enumerate}
        \item La \cle{suppression} : barre les segments redondants, les exemples trop développés, les métaphores, les commentaires de l'auteur.
        \item La \cle{condensation} : regroupe les énumérations, remplace les propositions subordonnées par des groupes nominaux ou des participes, utilise l'hyperonymie.
        \item La \cle{généralisation} : remplace les données chiffrées précises par des ordres de grandeur, les noms propres par des termes génériques.
    \end{enumerate}
    \item \textbf{Rédaction du résumé (version finale) :} Rédigez individuellement le résumé en \textbf{125 mots (\textpm 10 \%, soit 113 à 138 mots)}. Respectez l'ordre des idées, la neutralité (pas de « je », pas d'opinion personnelle), la cohérence syntaxique (phrases complexes, connecteurs) et l'orthographe.
    \item \textbf{Auto-évaluation et exposé :} Complétez la grille d'auto-évaluation fournie (annexe). Deux volontaires par groupe présentent leur démarche (choix de suppression, difficultés de reformulation, gestion du nombre de mots) en 5 minutes devant la classe.
\end{enumerate}

\courssub{Ressources à mobiliser}
\begin{aretenir}[Rappel des méthodes clés]
\textbf{1. Registres de langue :} Soutenu (vocabulaire précis, syntaxe complexe, pas de familiarité) — Courant (standard) — Familier (ellipses, vocabulaire expressif). \textit{Indice : l'éditorial utilise le registre soutenu pour légitimer l'autorité de l'opinion.}

\textbf{2. Structure argumentative type :}
\begin{itemize}
    \item \cle{Thèse} : position défendue (souvent au début ou à la fin).
    \item \cle{Arguments} : raisons logiques (cause, conséquence, analogie, autorité).
    \item \cle{Exemples/Illustrations} : faits, chiffres, références (Nigeria, 2019).
    \item \cle{Concession} : admission partielle de l'adversaire (connecteurs : \emph{certes, il est vrai que, cependant}).
    \item \cle{Conclusion} : réaffirmation de la thèse + ouverture/préconisation.
\end{itemize}

\textbf{3. Techniques de réduction (Contraction) :}
\begin{itemize}
    \item \cle{Suppression} : adjectifs épithètes non essentiels, compléments de phrase redondants, exemples détaillés, incises, répétitions.
    \item \cle{Condensation} : \emph{Subordination $\to$ Participation / Groupe nominal} (ex : « \emph{parce que le prix est bas} » $\to$ « \emph{du fait du bas prix} » / « \emph{le bas prix} »). \emph{Énumération $\to$ Hyperonyme} (ex : « riz thaïlandais, vietnamien, pakistanais » $\to$ « riz asiatiques importés »).
    \item \cle{Généralisation} : Chiffres précis $\to$ « plus de 300 milliards », « plusieurs millions de tonnes ». Noms propres $\to$ « le pays voisin », « la zone de production ».
\end{itemize}

\textbf{4. Reformulation (Outils lexicaux/syntaxiques) :}
\begin{itemize}
    \item \cle{Paronymie} : changer la catégorie grammaticale (verbe $\to$ nom : « \emph{produire} » $\to$ « \emph{la production} »).
    \item \cle{Hyperonymie/Hyponymie} : « tracteurs, moissonneuses, décortiqueuses » $\to$ « matériel agricole ».
    \item \cle{Voix passive / impersonnelle} : « Le gouvernement annonce » $\to$ « Il est annoncé par le gouvernement » / « L'annonce est faite ».
    \item \cle{Connecteurs de liaison} : \emph{de plus, par conséquent, en revanche, ainsi}.
\end{itemize}

\textbf{5. Contraintes du Bac Tech :} Longueur imposée (\textpm 10 \%), fidélité au sens, neutralité, correction linguistique, présentation soignée (nombre de mots indiqué en fin de texte).
\end{aretenir}

\courssub{Démarche et corrigé commenté}
\begin{exoresolu}[Corrigé commenté de l'activité d'intégration]
\tcblower
\textbf{Étape 1 — Identification du registre de langue}
\noindent \textbf{Réponse :} Le registre est \cle{soutenu}.
\noindent \textbf{Justification (3 indices) :}
\begin{enumerate}
    \item \textbf{Vocabulaire précis et abstrait} : « \emph{paradoxe, tenants du statu quo, vérité tronquée, coûts externes, ingénierie financière, substitution progressive, fiscalité déguisée} ». Ces termes relèvent du lexique économique et politique pointu.
    \item \textbf{Syntaxe complexe} : phrases longues à subordonnées multiples (ex : « \emph{Ce prix bas cache des coûts externes colossaux : dépendance aux cours mondiaux..., fuite des devises..., précarité...} »), usage du passé simple (« \emph{brandissent, entend-on, condamner} »), structures impersonnelles (« \emph{Il faut, c'est à ce prix que} »).
    \item \textbf{Absence de marques de l'oralité} : pas de familiarité, pas d'ellipses conversationnelles, ponctuation rigoureuse.
\end{enumerate}
\noindent \textbf{Effet recherché :} Conférer une \cle{autorité intellectuelle} et une \cle{légitimité} à l'argumentaire pour convaincre les décideurs et l'opinion publique éclairée.

\textbf{Étape 2 — Structure argumentative (Arbre)}

\begin{popfigure}
\begin{tikzpicture}[
    mindmap, grow cyclic, every node/.style=concept, concept color=popblue,
    level 1/.append style={level distance=4cm, sibling angle=90},
    level 2/.append style={level distance=3.5cm, sibling angle=45},
    level 3/.append style={level distance=3cm, sibling angle=30}
]
\node {Thèse centrale : \textbf{Substitution progressive (taxer importé + financer local + plafonner prix)}}
    child[concept color=popgreen] { node {Argument 1 : \textbf{Échec du statu quo (dépendance, fuite devises)}}
        child { node {Exemple : Facture 350 Mds FCFA (2023)} }
        child { node {Exemple : 15\% besoins couverts localement} }
    }
    child[concept color=poporange] { node {Argument 2 : \textbf{Référence nigériane (preuve de faisabilité)}}
        child { node {Exemple : Frontières fermées 2019} }
        child { node {Résultat : Prod. 4 $\to$ 9 Mt en 4 ans} }
    }
    child[concept color=poppurple] { node {Concession : \textbf{Risque social (prix local 550 vs 350 FCFA)}}
        child { node {Réplique : Prix bas = vérité tronquée (coûts cachés)} }
    }
    child[concept color=poppink] { node {Préconisations (Conditions de réussite)}}
        child { node {Variétés améliorées (NERICA)} }
        child { node {Unités décorticage + Routes} }
        child { node {Contrat de filière (prix plancher/plafond)} }
        child { node {Gestion rigoureuse fonds (ref scandale 2021)} }
;
\end{tikzpicture}
\legende{Structure argumentative de l'éditorial sous forme de carte heuristique.}
\end{popfigure}

\textbf{Étape 3 — Champs lexicaux et relations sémantiques}

\begin{center}
\begin{tabularx}{\linewidth}{|l|X|X|}\hline
\rowcolor{popblueL}
\textbf{Champ lexical} & \textbf{Occurrences (5)} & \textbf{Relations sémantiques repérées} \\ \hline
\cellcolor{popgreenL}\textbf{Économie / Commerce} & importation, facture, devises, droits de douane, protectionnisme, fiscalité, marché, prix plancher, prix plafond, substitution & \textbf{Hyperonymie} : \emph{importation, exportation, douane} $\subset$ \textbf{commerce international}. \textbf{Antonymie} : \emph{importation} $\leftrightarrow$ \emph{production locale / substitution}. \textbf{Paronymie} : \emph{taxer} (verbe) $\to$ \emph{taxation} (nom) ; \emph{substituer} $\to$ \emph{substitution}. \\ \hline
\cellcolor{poporangeL}\textbf{Agriculture / Production} & plaine de Mbo, riziculteurs, mécanisation, variétés, décorticage, grenier, exode rural, périmètres rizicoles, Anchor Borrowers & \textbf{Hyperonymie} : \emph{tracteurs, moissonneuses, décortiqueuses} $\subset$ \textbf{matériel agricole / intrants}. \textbf{Synonymie (contexte)} : \emph{riziculteurs} $\approx$ \emph{producteurs / agriculteurs}. \textbf{Champ sémantique} : \emph{grenier} (métaphore $\to$ réserve alimentaire nationale). \\ \hline
\end{tabularx}
\end{center}

\textbf{Étape 4 — Application des techniques de réduction (Extrait commenté)}

\begin{center}
\begin{tabularx}{\linewidth}{|l|X|}\hline
\rowcolor{popblueL}
\textbf{Segment original (82 mots)} & \textbf{Après réduction (cible $\approx$ 30 mots)} \\ \hline
« Depuis trois décennies, le Cameroun a fait le choix facile de l'importation massive de riz pour nourrir ses villes. En 2023, la facture a dépassé 350 milliards de FCFA, soit près de 1,2 million de tonnes débarquées au port de Douala, pendant que la plaine de Mbo, dans le Moungo, et les périmètres rizicoles de la Bénoué-Nord ne couvrent que 15 \% des besoins nationaux. » & \textbf{Suppression} : « Depuis trois décennies », « facile », « pour nourrir ses villes », « soit près de 1,2 million de tonnes débarquées au port de Douala », « dans le Moungo », « et les périmètres rizicoles de la Bénoué-Nord ». \\
& \textbf{Condensation} : « le Cameroun a fait le choix... » $\to$ « Le choix camerounais d'importer massivement... » ; « la facture a dépassé... » $\to$ « une facture $>$ 300 Mds FCFA ». \\
& \textbf{Généralisation} : « 350 milliards » $\to$ « plus de 300 milliards » ; « 15 \% » $\to$ « une part minoritaire ». \\
& \textbf{Résultat} : \emph{« Le choix camerounais d'importer massivement le riz génère une facture de plus de 300 milliards FCFA, alors que la production locale ne couvre qu'une part minoritaire des besoins. »} (31 mots) \\ \hline
\end{tabularx}
\end{center}

\textbf{Étape 5 — Rédaction du résumé final (125 mots visés)}

\noindent \textbf{Proposition de résumé (124 mots) :}

\emph{Le Cameroun, dépendant à 85 \% du riz importé (facture $>$ 300 Mds FCFA), envisage une substitution progressive pour soutenir la plaine de Mbo. L'éditorialiste réfute l'argument du prix bas de l'importé (350 FCFA/kg), qui masque une dépendance aux cours mondiaux et la précarité des producteurs locaux. L'exemple nigérian prouve qu'une protection frontalière couplée à un financement ciblé (programme \emph{Anchor Borrowers}) a décuplé la production en quatre ans. Cependant, interdire l'importation sans préparer l'offre locale serait irresponsable. La réussite exige des variétés améliorées, des unités de transformation industrielles, des routes rurales et un contrat de filière garantissant prix plancher au producteur et prix plafond (400 FCFA) au consommateur. La souveraineté alimentaire se construit par une fiscalité protectrice financant l'accompagnement technique, dans la transparence budgétaire.}

\noindent \textbf{Comptage :} 124 mots. \textbf{Conforme} (113--138 mots).
\noindent \textbf{Vérification des critères :}
\begin{itemize}
    \item \cle{Fidélité} : Thèse, concession, exemple Nigeria, préconisations, conclusion respectées.
    \item \cle{Neutralité} : « L'éditorialiste réfute », « L'exemple... prouve », « La réussite exige ». Pas de « je pense que ».
    \item \cle{Reformulation} : « fait le choix facile » $\to$ « envisage une substitution » ; « tenants du statu quo » $\to$ « argument du prix bas » ; « dignité d'une nation » $\to$ « souveraineté alimentaire ».
    \item \cle{Syntaxe} : Phrases complexes (subordination, participes : « couplée à », « garantissant »), connecteurs logiques (cependant, cependant, par conséquent implicite).
\end{itemize}

\textbf{Étape 6 — Grille d'auto-évaluation (Extrait)}

\begin{center}
\begin{tabularx}{\linewidth}{|l|X|c|}\hline
\rowcolor{popblueL}
\textbf{Critère} & \textbf{Indicateurs} & \textbf{Auto-éval. /5} \\ \hline
Respect de la longueur (113--138 mots) & Nombre de mots exact indiqué en fin de texte & \\
Fidélité au sens (thèse, args, conc., concl.) & Aucune idée majeure omise, aucune idée ajoutée & \\
Neutralité / Objectivité & Absence de « je », verbes d'opinion à la 3e pers., pas de jugement personnel & \\
Qualité de la reformulation & Changement de formulation (synonymie, nominalisation, hyperonymie), pas de copier-coller & \\
Cohérence syntaxique & Phrases complexes correctes, connecteurs variés, enchaînement logique & \\
Correction linguistique & Orthographe, grammaire, ponctuation, accords & \\
\hline
\end{tabularx}
\end{center}
\end{exoresolu}

\courssub{Bilan de l'activité}
\noindent Cette activité d'intégration a permis de vérifier l'acquisition opérationnelle des compétences ciblées par la leçon n°6 dans une situation authentique d'examen (Probatoire).
\begin{itemize}
    \item \textbf{Articulation réception/production :} L'analyse fine du texte source (registre, structure, lexique) s'est révélée être le \cle{prérequis indispensable} à une contraction réussie. Les élèves qui ont sauté l'étape 2 (arbre argumentatif) ont produit des résumés désorganisés ou incomplets.
    \item \textbf{Maîtrise des techniques de réduction :} La distinction opératoire entre \cle{suppression} (gain quantitatif), \cle{condensation} (gain syntaxique via la phrase complexe) et \cle{généralisation} (gain sémantique via l'hyperonymie) a structuré le travail de brouillon. L'usage de la \cle{nominalisation} (paronymie) et de la \cle{subordination} (participes présents, gérondifs, relatives) a été le levier principal pour passer de 512 à 125 mots sans perte de sens.
    \item \textbf{Gestion de la contrainte de longueur :} Le comptage rigoureux et l'ajustement final (suppression d'adverbes, fusion de phrases) ont mobilisé la compétence « réviser son texte ».
    \item \textbf{Transfert oral :} L'exposé de 5 minutes a obligé les élèves à \cle{métacognitiver} leur démarche (expliquer \emph{pourquoi} tel segment a été coupé, \emph{comment} telle reformulation a été choisie), consolidant ainsi les apprentissages.
    \item \textbf{Ancrage contextuel :} Le choix d'un éditorial camerounais (Cameroon Tribune, plaine de Mbo, SND30, référence Nigeria) a suscité l'engagement et permis de travailler le \cle{champ lexical} de l'économie rurale et de la souveraineté nationale, conforme aux attendus de l'enseignement technique (lien avec les spécialités agro-alimentaires, commerce, gestion).
\end{itemize}
\noindent \textbf{Perspective :} Cette activité sert de jalon pour l'épreuve certificative. La prochaine séance (compte-rendu et remédiation) exploitera les copies pour corriger les écueils récurrents (copier-coller résiduel, dépassement de mots, rupture de cohérence).

\newpage

\coursec{Méthodes et exercices résolus}

\coursec{Reformuler les idées essentielles d'un texte en vue d'un résumé}

\courssub{Méthode 1 : Dégager la structure argumentative d'un texte}

\begin{methode}[Identifier la structure argumentative]
Pour résumer un texte argumentatif, il faut d'abord repérer son \cle{ossature logique}. Procéder en étapes :
\begin{enumerate}
\item \textbf{Lire le texte en entier} une première fois, sans rien noter, pour saisir le sujet et le point de vue de l'auteur.
\item \textbf{Repérer la thèse} : c'est l'idée principale que l'auteur défend. Elle se trouve souvent au début ou à la fin du texte. Formule-repère : \emph{« L'auteur affirme que... »}.
\item \textbf{Découper le texte en étapes de l'argumentation} grâce aux \cle{connecteurs logiques} : \emph{d'abord, ensuite, en effet, cependant, par conséquent, enfin}.
\item \textbf{Distinguer arguments et exemples} : l'\cle{argument} est une raison générale ; l'\cle{exemple} l'illustre concrètement. Dans le résumé, on garde les arguments, on supprime la plupart des exemples.
\item \textbf{Rédiger le schéma} sous la forme : Thèse $\rightarrow$ Argument 1 $\rightarrow$ Argument 2 $\rightarrow$ Argument 3 $\rightarrow$ Conclusion.
\end{enumerate}
\textbf{Réflexe} : souligner au brouillon les connecteurs logiques ; ils trahissent l'organisation du texte.
\end{methode}

\begin{exoresolu}[Structure argumentative d'un texte sur l'école]
\textbf{Énoncé.} Lisez le texte suivant, puis dégagez sa structure argumentative (thèse, arguments, exemples, conclusion).

\emph{« L'école reste le meilleur instrument de promotion sociale. D'abord, elle donne à chaque enfant, quelle que soit son origine, les mêmes chances d'apprendre : un fils de paysan peut devenir ingénieur. Ensuite, elle transmet des savoirs indispensables à la vie professionnelle, comme le calcul et la maîtrise de la langue. Certes, certains réussissent sans diplôme, mais ce sont des exceptions qui confirment la règle. Enfin, l'école forme des citoyens capables de réfléchir et de voter en conscience. C'est pourquoi tout État soucieux de son avenir doit investir dans l'éducation. »}

\tcblower
\textbf{Solution détaillée.}

\textbf{Étape 1 : repérage de la thèse.} La première phrase affirme une position générale : \emph{« L'école reste le meilleur instrument de promotion sociale »}. C'est la \textbf{thèse} défendue par l'auteur.

\textbf{Étape 2 : repérage des connecteurs.} On souligne : \emph{D'abord}, \emph{Ensuite}, \emph{Certes... mais}, \emph{Enfin}, \emph{C'est pourquoi}. Ils organisent le texte en trois arguments, une concession et une conclusion.

\textbf{Étape 3 : identification des arguments et des exemples.}
\begin{itemize}
\item \textbf{Argument 1} : l'école donne les mêmes chances à tous les enfants. \emph{Exemple} : le fils de paysan devenu ingénieur.
\item \textbf{Argument 2} : elle transmet des savoirs utiles à la vie professionnelle. \emph{Exemples} : le calcul, la maîtrise de la langue.
\item \textbf{Concession réfutée} : certains réussissent sans diplôme (\emph{certes}), mais ce sont des exceptions (\emph{mais}).
\item \textbf{Argument 3} : elle forme des citoyens éclairés.
\end{itemize}

\textbf{Étape 4 : conclusion.} Introduite par \emph{C'est pourquoi} : l'État doit investir dans l'éducation.

\textbf{Schéma final :} Thèse (l'école, instrument de promotion sociale) $\rightarrow$ Arg. 1 (égalité des chances) $\rightarrow$ Arg. 2 (savoirs professionnels) $\rightarrow$ concession réfutée (réussites sans diplôme = exceptions) $\rightarrow$ Arg. 3 (formation du citoyen) $\rightarrow$ Conclusion (investir dans l'éducation).

\textbf{Conclusion :} le texte suit un plan dialectique simple avec concession ; le résumé devra conserver la thèse, les trois arguments et la conclusion, et éliminer les exemples.
\end{exoresolu}

\courssub{Méthode 2 : Appliquer les techniques de réduction}

\begin{methode}[Réduire un texte sans le trahir]
La \cle{contraction de texte} exige de réduire le texte à environ \textbf{un quart} de sa longueur (au Probatoire/Bac technique, la consigne précise toujours le nombre de mots). Quatre techniques :
\begin{enumerate}
\item \textbf{La suppression} : éliminer les exemples, les répétitions, les détails descriptifs, les citations, les énumérations secondaires.
\item \textbf{La généralisation} : remplacer une énumération par un terme générique (\cle{hyperonyme}) : \emph{« manguiers, palmiers, bananiers »} $\rightarrow$ \emph{« arbres fruitiers »}.
\item \textbf{La condensation} : remplacer un groupe de mots par un seul mot : \emph{« d'une manière rapide »} $\rightarrow$ \emph{« rapidement »} ; \emph{« rendre meilleur »} $\rightarrow$ \emph{« améliorer »}.
\item \textbf{La fusion} : réunir deux phrases proches en une seule phrase complexe grâce à une subordonnée.
\end{enumerate}
\textbf{Formule de contrôle} : si le texte fait $N$ mots, le résumé doit faire environ $N/4$ mots, avec une tolérance de $\pm 10\,\%$. Compter les mots et l'indiquer à la fin.
\end{methode}

\begin{exoresolu}[Réduire un paragraphe]
\textbf{Énoncé.} Réduisez le paragraphe suivant (environ 55 mots) à environ 14 mots, en appliquant les techniques de réduction. Précisez les techniques utilisées.

\emph{« Dans la cour de l'école, il y avait des manguiers, des palmiers, des bananiers et des goyaviers qui donnaient une ombre fraîche et agréable. Les élèves, pendant la récréation, se précipitaient d'une manière rapide sous ces arbres pour manger, jouer, bavarder, rire et se reposer un peu avant la reprise des cours de l'après-midi. »}

\tcblower
\textbf{Solution détaillée.}

\textbf{Comptage préalable.} Le texte compte environ 55 mots ; le quart vaut environ 14 mots, on vise entre 12 et 15 mots (tolérance $\pm 10\,\%$).

\textbf{Application des techniques :}
\begin{itemize}
\item \textbf{Généralisation} : \emph{« manguiers, palmiers, bananiers et goyaviers »} $\rightarrow$ \emph{« des arbres fruitiers »} (hyperonyme).
\item \textbf{Suppression} : les détails descriptifs \emph{« qui donnaient une ombre fraîche et agréable »} et l'énumération \emph{« manger, jouer, bavarder, rire »} sont secondaires : on les supprime ou on les condense en \emph{« se détendre »}.
\item \textbf{Condensation} : \emph{« d'une manière rapide »} $\rightarrow$ \emph{« rapidement »} ; \emph{« pendant la récréation »} est conservé car il situe le moment.
\item \textbf{Suppression} : \emph{« avant la reprise des cours de l'après-midi »} est un détail supprimable.
\end{itemize}

\textbf{Résumé proposé (14 mots) :} \emph{« Pendant la récréation, les élèves se réfugient rapidement sous les arbres fruitiers pour se détendre. »}

\textbf{Vérification :} 14 mots $\approx$ 55/4 $\approx$ 14 mots (tolérance respectée) ; l'idée essentielle (les élèves profitent des arbres à la récréation) est conservée ; aucune idée nouvelle n'a été ajoutée.

\textbf{Conclusion :} la réduction respecte la consigne de longueur et la fidélité au texte.
\end{exoresolu}

\courssub{Méthode 3 : Reformuler les idées essentielles}

\begin{methode}[Reformuler sans copier]
La \cle{reformulation} consiste à exprimer les idées de l'auteur \textbf{avec ses propres mots}, sans changer le sens. Réflexes :
\begin{enumerate}
\item \textbf{Remplacer les mots par des synonymes} ou des termes du même \cle{champ lexical} : \emph{« demeure »} $\rightarrow$ \emph{« maison »}.
\item \textbf{Changer la construction de la phrase} : passer de l'actif au passif, transformer une subordonnée en proposition principale (ou l'inverse), nominaliser un verbe (\emph{« il réfléchit »} $\rightarrow$ \emph{« sa réflexion »}).
\item \textbf{Ne jamais recopier plus de trois mots consécutifs} du texte (sauf termes techniques irremplaçables).
\item \textbf{Conserver le point de vue de l'auteur} : ne pas écrire \emph{« je pense que »} ; le résumé est neutre et objectif.
\item \textbf{Vérifier} : relire en comparant avec le texte ; si le sens a changé, corriger.
\end{enumerate}
\textbf{Interdit} : la paraphrase mot à mot (changer un seul mot par phrase) et le collage de phrases du texte.
\end{methode}

\begin{exoresolu}[Reformuler deux phrases]
\textbf{Énoncé.} Reformulez les phrases suivantes sans en changer le sens, en utilisant au moins deux procédés différents par phrase.

1. \emph{« Parce qu'elle néglige l'éducation des filles, cette société se prive de la moitié de son intelligence. »}

2. \emph{« Le chef du village, qui était un homme sage et expérimenté, décida de consulter les anciens avant de répondre aux étrangers. »}

\tcblower
\textbf{Solution détaillée.}

\textbf{Phrase 1.} Procédés utilisés : changement de construction (la subordonnée de cause devient une principale reliée par \emph{donc}) et synonymie (\emph{se prive} $\rightarrow$ \emph{renonce} ; \emph{intelligence} $\rightarrow$ \emph{capacités intellectuelles}).

Reformulation : \emph{« Cette société néglige l'éducation des filles ; elle renonce donc à la moitié de ses capacités intellectuelles. »}

Vérification : la cause (négligence de l'éducation des filles) et la conséquence (perte de la moitié de l'intelligence collective) sont conservées ; aucun mot du texte n'est recopié en séquence longue.

\textbf{Phrase 2.} Procédés utilisés : suppression de la relative (l'épithète \emph{« sage et expérimenté »} est condensée en \emph{« avisé »}), synonymie (\emph{les anciens} $\rightarrow$ \emph{les sages du village} ; \emph{étrangers} $\rightarrow$ \emph{visiteurs}), inversion de l'ordre (l'action de consulter précède la réponse, exprimée par un gérondif ou une tournure temporelle).

Reformulation : \emph{« Avant de répondre aux visiteurs, le chef du village, homme avisé, prit l'avis des sages de la communauté. »}

Vérification : les trois idées (la sagesse du chef, la consultation des anciens, la réponse différée) sont présentes ; le sens est identique.

\textbf{Conclusion :} chaque reformulation combine synonymie et transformation syntaxique, ce qui garantit à la fois la fidélité et l'originalité de l'expression.
\end{exoresolu}

\courssub{Méthode 4 : Utiliser les registres de langue à bon escient}

\begin{methode}[Choisir le bon registre]
Le \cle{registre de langue} dépend de la situation de communication : le destinataire, le lieu, le sujet.
\begin{enumerate}
\item \textbf{Identifier la situation} : à qui parle-t-on (ou écrit-on) ? Dans quel cadre ?
\item \textbf{Choisir le registre adapté} :
\begin{itemize}
\item \cle{Registre soutenu} : discours officiels, textes littéraires, correspondance administrative (vocabulaire riche, syntaxe élaborée) ;
\item \cle{Registre courant} : conversations polies, lettres, exposés en classe (vocabulaire correct, phrases complètes) ;
\item \cle{Registre familier} : échanges entre proches (abréviations, argot, tournures relâchées).
\end{itemize}
\item \textbf{Pour transposer} d'un registre à l'autre : remplacer le vocabulaire familier par son équivalent courant ou soutenu (\emph{« boulot »} $\rightarrow$ \emph{« travail »} $\rightarrow$ \emph{« emploi »}), compléter les phrases, supprimer les interjections.
\end{enumerate}
\textbf{Réflexe Bac} : dans un résumé, une dissertation ou une lettre officielle, le registre exigé est toujours le registre \textbf{courant ou soutenu} : jamais de familier.
\end{methode}

\begin{exoresolu}[Transposer un dialogue en registre soutenu]
\textbf{Énoncé.} Le dialogue suivant est écrit en registre familier. Transposez-le en registre soutenu, comme dans un compte rendu officiel.

\emph{« — Hé patron, c'est pas cool, ton truc ! Le matos est naze et les gars sont crevés.\\
— T'inquiète, on va se débrouiller. Faut juste qu'on bosse plus tôt demain. »}

\tcblower
\textbf{Solution détaillée.}

\textbf{Étape 1 : repérage des marques du familier.} \emph{Hé} (interjection), \emph{patron} (adresse familière), \emph{cool} (argot), \emph{ton truc} (démonstratif vague + possessif familier), \emph{matos} (abréviation de \emph{matériel}), \emph{naze} (argot = en mauvais état), \emph{les gars} (= les ouvriers), \emph{crevés} (= fatigués), \emph{t'inquiète} (= ne vous inquiétez pas), \emph{se débrouiller} (= trouver une solution), \emph{bosser} (= travailler).

\textbf{Étape 2 : substitution par des équivalents soutenus.}
\begin{center}
\begin{tabularx}{\linewidth}{|l|X|}\hline
\rowcolor{popblueL} \textbf{Familier} & \textbf{Soutenu} \\ \hline
Hé patron & Monsieur le directeur \\ \hline
c'est pas cool, ton truc & cette situation est inacceptable \\ \hline
le matos est naze & le matériel est hors d'usage \\ \hline
les gars sont crevés & les ouvriers sont épuisés \\ \hline
t'inquiète, on va se débrouiller & soyez sans inquiétude, nous trouverons une solution \\ \hline
faut qu'on bosse plus tôt & il conviendra de commencer le travail plus tôt \\ \hline
\end{tabularx}
\end{center}

\textbf{Étape 3 : transposition rédigée.} \emph{« — Monsieur le directeur, cette situation est inacceptable : le matériel est hors d'usage et les ouvriers sont épuisés.\\
— Soyez sans inquiétude, nous trouverons une solution ; il conviendra simplement de commencer le travail plus tôt demain. »}

\textbf{Conclusion :} la transposition conserve le contenu (deux plaintes, une réponse rassurante, une proposition) tout en adoptant le vocabulaire et la syntaxe du registre soutenu.
\end{exoresolu}

\courssub{Méthode 5 : Exploiter champs lexicaux et relations de sens}

\begin{methode}[Champ lexical, paronymie, hyperonymie]
\begin{enumerate}
\item \textbf{Repérer un \cle{champ lexical}} : relever les mots qui se rapportent à une même notion (ex. champ de l'école : \emph{classe, cahier, maître, leçon}). Il sert à identifier le thème du texte, donc son idée essentielle.
\item \textbf{Distinguer champ lexical et \cle{champ sémantique}} : le champ sémantique rassemble les sens d'un seul mot selon le contexte (ex. \emph{voler} : champ sémantique du vol d'oiseau ou du délit).
\item \textbf{Utiliser l'\cle{hyperonymie}/\cle{hyponymie}} pour généraliser dans un résumé : l'hyperonyme englobe les hyponymes (\emph{fruit} $\supset$ \emph{mangue, banane}).
\item \textbf{Éviter les \cle{paronymes}} : mots de forme proche mais de sens différent (\emph{éminemment/immanquablement} ; \emph{inauguration/inaugural}). Vérifier le sens exact avant de reformuler.
\end{enumerate}
\textbf{Réflexe} : pour résumer une énumération, chercher l'hyperonyme ; pour reformuler, vérifier qu'on n'a pas glissé vers un paronyme.
\end{methode}

\begin{exoresolu}[Champ lexical et généralisation]
\textbf{Énoncé.} 1. Relevez le champ lexical dominant du passage suivant et déduisez-en le thème. 2. Remplacez l'énumération par un hyperonyme pour préparer un résumé. 3. Corrigez la faute de paronymie dans la phrase b.

a. \emph{« Le maître écrit au tableau, les élèves ouvrent leurs cahiers, recopient la leçon, puis répondent aux questions de l'interrogation. »}

b. \emph{« Le résumé est immanquablement plus court que le texte. » (le mot souligné est-il correct ?)}

\tcblower
\textbf{Solution détaillée.}

\textbf{1. Champ lexical.} Mots relevés : \emph{maître, tableau, élèves, cahiers, leçon, interrogation}. Tous renvoient à la notion d'\textbf{école / enseignement}. Le thème du passage est donc une scène de classe. Dans un résumé, ce thème doit apparaître dès la première phrase.

\textbf{2. Généralisation.} L'énumération \emph{« le maître écrit au tableau, les élèves ouvrent leurs cahiers, recopient la leçon »} peut être remplacée par l'hyperonyme d'ensemble : \emph{« les activités scolaires habituelles »} ou \emph{« le déroulement ordinaire d'une leçon »}. Gain : une douzaine de mots réduits à cinq, sans perte du sens essentiel.

\textbf{3. Paronymie.} \emph{Immanquablement} signifie \emph{« qui ne peut manquer d'arriver, inévitablement »}. La phrase veut dire que le résumé est \emph{nécessairement, par définition} plus court : le mot juste est \emph{nécessairement} ou \emph{par définition}. \emph{Éminemment} (au plus haut degré) ne convient pas au comparatif.

Correction : \emph{« Le résumé est par définition plus court que le texte »} (ou \emph{« nécessairement plus court »}).

\textbf{Conclusion :} le repérage du champ lexical confirme le thème (l'école), l'hyperonyme permet la réduction, et la vigilance sur les paronymes évite une faute de sens.
\end{exoresolu}

\courssub{Méthode 6 : Construire des phrases complexes pour rédiger le résumé}

\begin{methode}[La phrase complexe au service du résumé]
Une \cle{phrase complexe} contient au moins deux propositions. Elle permet de fusionner plusieurs idées du texte en peu de mots.
\begin{enumerate}
\item \textbf{Subordonnée relative} (pronom \emph{qui, que, où, dont}) : condense une information : \emph{« le chef, qui était sage, décida... »}.
\item \textbf{Subordonnée conjonctive} (conjonction \emph{parce que, bien que, lorsque, si}) : exprime la cause, la concession, le temps : elle reproduit les liens logiques du texte.
\item \textbf{Proposition participiale ou gérondif} : \emph{« négligeant l'école, cette société... »} : réduction élégante d'une subordonnée.
\item \textbf{Proposition infinitive} : \emph{« avant de répondre... »}.
\end{enumerate}
\textbf{Démarche de rédaction du résumé} : (1) rédiger une phrase d'introduction donnant le thème et la thèse ; (2) enchaîner les arguments dans des phrases complexes reliées par des connecteurs (\emph{d'abord, en effet, cependant, enfin}) ; (3) conclure par l'idée finale de l'auteur ; (4) compter les mots et vérifier la fidélité.
\end{methode}

\begin{exoresolu}[Rédiger un résumé complet]
\textbf{Énoncé.} Résumez le texte suivant (environ 95 mots) en environ 24 mots, en utilisant au moins une phrase complexe avec subordonnée.

\emph{« Dans de nombreux villages, l'eau potable manque cruellement. Les femmes parcourent chaque matin plusieurs kilomètres pour atteindre le puits le plus proche, ce qui les fatigue et les empêche de vaquer à d'autres occupations. Les enfants, eux, boivent souvent l'eau des marigots, ce qui provoque des maladies comme le choléra et la dysenterie. Pourtant, des solutions existent : le forage de puits modernes, l'installation de bornes-fontaines, la formation des populations à l'hygiène. Il appartient donc aux autorités, avec l'aide des organisations internationales, de financer ces projets afin que chaque village dispose enfin d'une eau saine. »}

\tcblower
\textbf{Solution détaillée.}

\textbf{Étape 1 : structure argumentative.}
\begin{itemize}
\item \textbf{Thèse} : l'eau potable manque dans les villages et il faut financer des projets pour y remédier.
\item \textbf{Argument 1} : les femmes s'épuisent à chercher l'eau (conséquence : fatigue, perte de temps).
\item \textbf{Argument 2} : les enfants boivent une eau insalubre (conséquence : maladies).
\item \textbf{Arguments-solutions} : forages, bornes-fontaines, formation à l'hygiène.
\item \textbf{Conclusion} : les autorités et les organisations internationales doivent financer.
\end{itemize}

\textbf{Étape 2 : réduction.} Suppression des exemples de maladies (\emph{choléra, dysenterie}) et du détail \emph{« chaque matin »} ; généralisation de \emph{« forage de puits modernes, installation de bornes-fontaines »} $\rightarrow$ \emph{« aménagements hydrauliques »}.

\textbf{Étape 3 : reformulation en phrases complexes.}

Résumé proposé : \emph{« Faute d'eau potable, les villageoises s'épuisent à la chercher et les enfants tombent malades ; des aménagements hydrauliques, financés par les autorités et les organisations internationales, régleraient ce problème. »}

\textbf{Étape 4 : vérifications.} Comptage : 26 mots. Le texte source en compte environ 95, le quart vaut 24 mots (tolérance $\pm 10\,\%$ : 22 à 26 mots). La longueur est conforme.

Subordonnées et procédés : proposition infinitive (\emph{à la chercher}), coordination de conséquences, généralisation (\emph{aménagements hydrauliques}), suppression des exemples. Le point de vue de l'auteur est conservé, aucune opinion personnelle n'est ajoutée.

\textbf{Conclusion :} le résumé respecte la longueur imposée, reproduit la thèse, les deux arguments et la conclusion, dans des phrases complexes correctes.
\end{exoresolu}

\begin{aretenir}[Synthèse de la démarche complète pour le Probatoire]
La contraction de texte suit un protocole rigoureux en quatre temps :
\begin{enumerate}
\item \textbf{Analyse} : lecture active, repérage de la thèse, découpage argumentatif (connecteurs), distinction arguments/exemples.
\item \textbf{Réduction} : suppression (exemples, détails), généralisation (hyperonymes), condensation (adverbes, noms), fusion (subordination).
\item \textbf{Reformulation} : synonymie, transposition syntaxique (nominalisation, passif, subordination), changement de point de vue (neutre), zéro copier-coller ($>3$ mots).
\item \textbf{Contrôle} : comptage exact des mots (indiqué en fin de copie), vérification de la fidélité (thèse, arguments, conclusion présents), correction de la langue (registre soutenu/courant, orthographe, paronymes).
\end{enumerate}
\end{aretenir}

\begin{savaistu}[Lien avec l'œuvre au programme : \emph{Le Lion et la Perle} de Wole Soyinka]
Les séquences de \emph{lecture méthodique n°4, 5, 6} et l'\emph{Exposé n°1} sur cette pièce exigent les mêmes compétences. Exemple d'application :
\begin{itemize}
\item \textbf{Structure argumentative} : dans la scène de la classe (Lakunle vs Sidi), repérer la thèse de Lakunle (« la modernité libère la femme »), ses arguments (éducation, fin de la dot), la concession (beauté de Sidi), la conclusion (mariage moderne).
\item \textbf{Réduction/Reformulation} : contracter le monologue de Baroka sur la ruse du lion (Acte 2) en remplaçant les images par des concepts abstraits (\emph{« ruse, patience, expérience »} $\rightarrow$ \emph{« sagacité politique »}).
\item \textbf{Registres} : opposer le registre soutenu/ironique de Lakunle (\emph{« exquisité, civilisation »}) au registre imagé/proverbial de Baroka et Sidi (\emph{« le vieux cuir, la perle »}).
\item \textbf{Champs lexicaux} : champ de la tradition (\emph{dot, fétiche, oracle}) vs champ de la modernité (\emph{école, machine, journal}).
\end{itemize}
Maîtriser ces outils techniques permet de réussir l'exposé et les questions de synthèse sur l'œuvre.
\end{savaistu}

\begin{attention}[Les 5 erreurs qui coûtent des points au Bac]
\begin{enumerate}
\item \textbf{Recopier des phrases entières du texte.} Le correcteur surligne les collages : plus de trois mots consécutifs repris sans guillemets est sanctionné. Reformulez toujours.
\item \textbf{Donner son avis personnel.} Le résumé est \emph{neutre} : écrire \emph{« je pense que l'auteur a raison »} fait perdre les points d'objectivité. On rapporte la pensée de l'auteur, pas la sienne.
\item \textbf{Négliger la consigne de longueur.} Un résumé de 80 mots quand 40 sont demandés est pénalisé, même s'il est bon. Comptez les mots et indiquez le total à la fin.
\item \textbf{Supprimer un argument essentiel en gardant un exemple.} C'est l'inverse qu'il faut faire : les exemples, citations et détails sautent ; la thèse, les arguments et la conclusion restent.
\item \textbf{Confondre les paronymes et mal construire la phrase complexe.} \emph{Immanquablement} pour \emph{nécessairement}, ou une relative sans antécédent clair (\emph{qui} flottant), créent des contresens. Relisez chaque phrase en vérifiant le sens exact des mots et le rapport des propositions.
\end{enumerate}
\end{attention}

\newpage

\coursec{Exercices}

\courssub{Je vérifie mes connaissances}

\begin{exercice}[QCM : Registres de langue et marqueurs]
Pour chacune des phrases ci-dessous, cochez le registre de langue correspondant et citez \textbf{deux} marqueurs linguistiques qui justifient votre choix.

\begin{enumerate}
    \item « Je vous saurais gré de bien vouloir m'indiquer la marche à suivre. »
    \item « T'inquiète, je gère, c'est bon. »
    \item « L'individu a été interpellé par les forces de l'ordre pour infraction au code de la route. »
\end{enumerate}

\textbf{Registres proposés :} Soutenu / Courant / Familier / Administratif (technique).
\end{exercice}

\begin{exercice}[Vrai ou Faux justifié : Structure argumentative]
Déclarez si les affirmations suivantes sont vraies ou fausses. Justifiez chaque réponse en vous appuyant sur les notions de thèse, d'arguments, de concession et de connecteurs logiques.

\begin{enumerate}
    \item La thèse d'un texte argumentatif se trouve toujours au premier paragraphe.
    \item Un argument d'autorité repose sur l'opinion personnelle de l'auteur.
    \item La concession (« Certes, ... mais ... ») affaiblit l'argumentation globale.
    \item Les connecteurs logiques (donc, cependant, en outre) organisent l'enchaînement des idées.
\end{enumerate}
\end{exercice}

\begin{exercice}[Paronymie : Choix et définition]
Complétez les phrases suivantes avec le paronyme correct (choisi entre parenthèses). Donnez ensuite la définition du mot \textbf{non utilisé} pour montrer la nuance de sens.

\begin{enumerate}
    \item L'\dots\ de l'arrivée du ministre a provoqué une grande agitation dans le protocole. (\emph{éminence} / \emph{imminence})
    \item Le témoin a dû \dots\ dire la vérité sous serment. (\emph{abjurer} / \emph{adjurer})
    \item L'\dots\ des plantes médicinales dans l'eau bouillante libère les principes actifs. (\emph{infusion} / \emph{confusion})
\end{enumerate}
\end{exercice}

\begin{exercice}[Réduction directe : Phrase complexe]
Voici une phrase extraite d'un rapport officiel (42 mots) :
\begin{quote}
« Il est impératif que l'ensemble des parties prenantes, y compris les partenaires techniques et financiers, puissent se réunir dans les plus brefs délais afin de procéder à l'évaluation minutieuse des résultats obtenus. »
\end{quote}
Appliquez \textbf{trois} techniques de réduction (suppression, condensation, généralisation) pour la ramener à \textbf{15 mots maximum}. Indiquez entre crochets la technique utilisée pour chaque modification.
\end{exercice}

\courssub{Je m'entraîne}

\begin{exercice}[Repérage des registres : Extrait « Le Lion et la perle » (adapté)]
Lisez l'extrait suivant, inspiré des dialogues de Wole Soyinka. Pour \textbf{chaque réplique}, identifiez le registre de langue et justifiez par \textbf{deux} marqueurs précis (vocabulaire, syntaxe, phonétique).

\begin{description}
    \item[LAKUNLE :] « Mon esprit s'élève au-dessus de ces coutumes archaïques. L'émancipation de la femme passe par l'instruction, non par la dot. »
    \item[SIDI :] « Toi, l'homme aux grands mots ! Tu parles comme un livre. Moi, je sais que la femme vaut son prix. »
    \item[BAROKA :] « Il convient d'examiner la requête avec la sagesse des ancêtres. Le progrès ne saurait effacer notre identité. »
\end{description}
\end{exercice}

\begin{exercice}[Champs lexicaux et hyperonymes : Agriculture au Cameroun]
Voici une liste de 20 mots extraits d'un texte sur l'agriculture camerounaise :
\begin{center}
cacao, café, bananier, plantain, maïs, mil, sorgho, igname, manioc, patate, cotonnier, palmier à huile, hévéa, engrais, pesticide, tracteur, charrue, irrigation, coopérative, rendement
\end{center}

\begin{enumerate}
    \item Classez ces mots en \textbf{trois champs lexicaux} cohérents (donnez un titre à chaque champ).
    \item Repérez \textbf{trois énumérations} implicites dans cette liste et remplacez chacune par son \textbf{hyperonyme} (ex : cacao, café, cotonnier $\rightarrow$ cultures d'exportation).
\end{enumerate}
\end{exercice}

\begin{exercice}[Transformation de phrase complexe : Discours officiel]
Transformez la phrase complexe suivante (46 mots) en une phrase simple de \textbf{16 mots maximum}, en conservant l'information essentielle. Justifiez chaque étape (suppression des incidants, nominalisation, pronominalisation, fusion).

\begin{quote}
« Le Ministre de l'Éducation Technique, considérant que la formation professionnelle doit répondre aux besoins réels du marché de l'emploi, a décidé de réviser les curricula en concertation avec les opérateurs économiques. »
\end{quote}
\end{exercice}

\begin{exercice}[Dégager et reformuler les idées essentielles]
Texte support (130 mots) : \emph{« L'introduction des langues nationales dans l'enseignement technique au Cameroun suscite un débat vif. Les partisans arguent qu'elle facilite l'acquisition des compétences techniques par une meilleure compréhension des concepts abstraits. Ils citent l'exemple de pays asiatiques ayant réussi leur industrialisation via leurs langues maternelles. Les opposants craignent une fragmentation linguistique nuisible à l'unité nationale et à l'insertion professionnelle internationale, l'anglais et le français restant les langues de la science et du commerce mondial. Une approche bilingue progressive, associant langue nationale et langue officielle, apparaît comme une piste de conciliation réaliste. »}

\begin{enumerate}
    \item Identifiez la \textbf{thèse} défendue (ou la thèse de conciliation).
    \item Dégagez la \textbf{structure argumentative} (arguments pour / contre / conciliation).
    \item Reformulez les \textbf{trois idées essentielles} en une phrase chacune, sans copier les phrases du texte.
\end{enumerate}
\end{exercice}

\courssub{Je me prépare au Bac}

\begin{exercice}[Sujet type Bac Tech : Contraction de texte argumentatif]
\textbf{Texte à contracter (environ 600 mots) :}

\emph{« La place des langues nationales dans l'enseignement technique au Cameroun : un levier pour l'insertion professionnelle ?}

\emph{Depuis l'indépendance, le système éducatif camerounais fonctionne officiellement en français et en anglais, héritage colonial qui structure encore l'administration et l'enseignement supérieur. Pourtant, plus de 250 langues nationales coexistent sur le territoire, véhicules de cultures et de savoirs endogènes. L'enseignement technique, voué à la formation de techniciens opérationnels, fait face à un paradoxe : les élèves maîtrisent mal les langues officielles d'enseignement, ce qui freine l'assimilation des concepts scientifiques et techniques complexes.}

\emph{Des expériences menées dans la région de l'Ouest, notamment en bamileke (ghomala') et en bassa, montrent que l'explication des phénomènes physiques (électricité, thermodynamique) ou des procédés industriels (soudure, maintenance) dans la langue maternelle accélère la compréhension. Les apprenants manipulent mieux le vocabulaire technique lorsqu'il est ancré dans leur référentiel culturel. L'UNESCO recommande d'ailleurs l'enseignement dans la langue maternelle au moins jusqu'au secondaire pour garantir la qualité des apprentissages.}

\emph{Toutefois, des obstacles majeurs persistent. La standardisation orthographique de nombreuses langues n'est pas achevée. Le manque de manuels techniques traduits et de formateurs bilingues (langue nationale / français ou anglais) est criant. Par ailleurs, le marché de l'emploi, tant national qu'international, exige la maîtrise des langues officielles. Une formation exclusivement monolingue en langue nationale enfermerait le technicien dans un marché local limité.}

\emph{La solution réside dans un bilinguisme additif : utiliser la langue nationale comme langue d'explication et d'ancrage cognitif au cycle d'observation et au début du cycle technique, tout en renforçant parallèlement l'apprentissage intensif du français ou de l'anglais comme langues de communication professionnelle et de documentation technique. Cette approche « passerelle » valorise l'identité culturelle sans sacrifier l'employabilité. Elle suppose une volonté politique forte : élaboration de terminologies techniques normalisées, formation des enseignants, édition de supports pédagogiques adaptés. C'est à ce prix que l'enseignement technique camerounais pourra former des techniciens compétents, fiers de leurs racines et ouverts au monde. »}

\begin{enumerate}
    \item \textbf{Analyse de la structure argumentative} (4 points) :\\
    a) Formulez la \textbf{thèse principale} défendue par l'auteur.\\
    b) Identifiez les \textbf{trois grandes parties} de l'argumentation (constat / arguments pour / obstacles / solution) et leur fonction.\\
    c) Relevez \textbf{deux connecteurs logiques} marquant l'opposition et la conséquence.
    \item \textbf{Dégagement des idées essentielles} (6 points) :\\
    Résumez le texte en \textbf{7 idées forces} numérotées, rédigées en phrases complètes, sans citation littérale, en respectant l'ordre logique du texte.
    \item \textbf{Rédaction du résumé} (10 points) :\\
    Rédigez le \textbf{résumé} du texte en \textbf{150 mots $\pm$ 10\%} (soit 135 à 165 mots).\\
    \textbf{Contraintes :} Registre soutenu, impersonnalité, cohérence, respect de la structure argumentative, zéro copier-coller, zéro avis personnel.
\end{enumerate}
\end{exercice}

\begin{exercice}[Compte-rendu professionnel : Réunion de coopérative]
\textbf{Procès-verbal de la réunion de la Coopérative « Espoir Paysan » de Bafoussam (extrait, 380 mots) :}

\emph{« Le 15 octobre 2024, à 10h, salle des fêtes de la mairie de Bafoussam II. Présents : 42 adhérents sur 60, le président M. Ngoufo, le secrétaire M. Talla, le trésorier Mme Kemayou. Absents excusés : 8. Ordre du jour : 1. Bilan de la campagne café 2023-2024. 2. Préparation campagne cacao 2024-2025. 3. Acquisition d'un camion. 4. Divers.}

\emph{Le président ouvre la séance. Bilan café : production en baisse de 15\% (120 tonnes contre 141 l'an passé) due aux pluies tardives et à la vétusté des pulpeuses. Recettes : 72 millions FCFA. Charges : 45 millions. Bénéfice net : 27 millions, soit 450 000 FCFA par part sociale. Approbation à l'unanimité.}

\emph{Campagne cacao : prévision 200 tonnes. Besoin urgent : 5 pulpeuses neuves (coût 1,2M FCFA/unité) et 2000 sacs de jute. Demande de subvention auprès du MINADER et du Fonds de développement rural. Le trésorier alerte sur la trésorerie : 8 millions disponibles, insuffisants pour l'avance intrants. Vote : recours à un crédit agricole CAMCCUL à 5\% sur 3 ans, adopté à 38 voix pour, 2 contre, 2 abstentions.}

\emph{Acquisition camion : le véhicule actuel (15 ans) tombe en panne fréquente. Devis Toyota Dyna 3,5T : 18 millions FCFA. Proposition : apport coopérative 10M + crédit 8M. Reporté à l'AG extraordinaire de décembre pour étude de rentabilité.}

\emph{Divers : formation 10 jeunes sur taille/ greffage (partenariat IRAD). Prochaine réunion : 15 janvier 2025. Clôture 13h30. »}

\begin{enumerate}
    \item Rédigez le \textbf{compte-rendu officiel} de cette réunion en \textbf{100 mots $\pm$ 10\%}.
    \item \textbf{Contraintes :} Style administratif (passif, infinitif, noms d'action), structure standard (objet, date, participants, décisions, votes, échéances), sélection de l'essentiel (chiffres clés, votes, reports), objectivité totale.
\end{enumerate}
\end{exercice}

\begin{exercice}[Remédiation : Correction d'un résumé d'élève]
Voici un résumé d'élève (165 mots) sur le texte « La place des langues nationales... » (sujet précédent). Il contient \textbf{cinq erreurs types} repérées ci-dessous.

\emph{« L'auteur pense que les langues nationales sont très importantes pour l'enseignement technique au Cameroun. Il dit que "les élèves maîtrisent mal les langues officielles" et que ça freine l'apprentissage. Il donne l'exemple du bamileke et du bassa où ça marche mieux. Moi je pense qu'il a raison car chez moi on parle bassa à la maison. Mais il y a des problèmes : pas de livres, pas de profs formés. L'anglais et le français sont quand même nécessaires pour le travail. L'auteur propose le bilinguisme additif comme solution miracle. Il faut que l'État fasse des dictionnaires techniques et forme les maîtres. C'est comme ça qu'on aura de bons techniciens. »}

\textbf{Erreurs à identifier et corriger :}
\begin{enumerate}
    \item Copie de phrases / passages littéraux (repérez la citation).
    \item Avis personnel / subjectivité (repérez la marque de la 1ère personne).
    \item Dépassement de la longueur prescrite (calculez le nombre de mots).
    \item Registre familier / oral (repérez au moins deux marqueurs).
    \item Contresens par paronymie / approximation lexicale (repérez « solution miracle » vs « piste de conciliation réaliste »).
\end{enumerate}

\textbf{Travail demandé :}
\begin{enumerate}
    \item Pour chaque erreur, citez le passage fautif, nommez le type d'erreur et expliquez pourquoi c'est pénalisant au Bac.
    \item Rédigez la \textbf{version corrigée} du résumé (150 mots $\pm$ 10\%), en appliquant les techniques de reformulation, de réduction et le registre soutenu.
\end{enumerate}
\end{exercice}

\newpage

\coursec{Corrigés des exercices}

\coursec{Leçon 06 : Reformuler les idées essentielles d'un texte en vue d'un résumé}
\courssub{Exercices corrigés et remédiation}

\begin{corrige}[Exercice 1 — QCM : Registres de langue et marqueurs]
\textbf{1. « Je vous saurais gré de bien vouloir m'indiquer la marche à suivre. »} $\rightarrow$ Registre \cle{soutenu}. Marqueurs : la tournure de politesse élaborée et conditionnelle « je vous saurais gré » ; l'infinitif de politesse « de bien vouloir » ; le vocabulaire recherché (« la marche à suivre »). Syntaxe longue et soignée.

\textbf{2. « T'inquiète, je gère, c'est bon. »} $\rightarrow$ Registre \cle{familier}. Marqueurs : l'élision orale « T'inquiète » (pour « ne t'inquiète pas ») avec suppression de la négation ; le verbe argotique « gère » ; la phrase télégraphique « c'est bon ».

\textbf{3. « L'individu a été interpellé par les forces de l'ordre pour infraction au code de la route. »} $\rightarrow$ Registre \cle{administratif} (technique). Marqueurs : la voix passive « a été interpellé » ; le vocabulaire juridique et institutionnel (« interpellé », « forces de l'ordre », « infraction au code de la route ») ; l'impersonnalité (pas de nom propre, « l'individu »).
\end{corrige}

\begin{attention}[Point de vigilance]
Ne pas confondre \emph{soutenu} et \emph{administratif} : le soutenu vise l'élégance de l'expression, l'administratif vise la précision juridique ou technique, souvent au passif. Au Bac, citez toujours \textbf{deux marqueurs précis} (vocabulaire + syntaxe), jamais une impression vague.
\end{attention}

\begin{corrige}[Exercice 2 — Vrai ou Faux justifié : Structure argumentative]
\textbf{1. FAUX.} La thèse peut être placée au début (plan déductif), à la fin (plan inductif, après l'examen des faits) ou au milieu. Ce qui compte est de la repérer grâce aux marqueurs (« nous soutenons que », « il faut donc », « en définitive ») et à sa fonction : c'est l'idée que l'auteur cherche à faire admettre.

\textbf{2. FAUX.} L'argument d'autorité repose sur la parole d'une personne ou d'une institution reconnue comme compétente (expert, organisme comme l'UNESCO, texte officiel), et non sur l'opinion personnelle de l'auteur. L'opinion personnelle relève de l'argument d'opinion ou du témoignage.

\textbf{3. FAUX.} La concession (« Certes... mais... ») \emph{renforce} l'argumentation : en reconnaissant un argument adverse avant de le réfuter, l'auteur montre son honnêteté intellectuelle et prévient les objections. C'est un procédé de persuasion efficace.

\textbf{4. VRAI.} Les connecteurs logiques organisent l'enchaînement des idées : addition (\emph{en outre, de plus}), opposition (\emph{cependant, toutefois, mais}), conséquence (\emph{donc, par conséquent, c'est pourquoi}), cause (\emph{car, puisque}). Ils rendent la démonstration lisible et cohérente.
\end{corrige}

\begin{attention}[Point de vigilance]
Au Probatoire, la question sur la structure argumentative exige de \textbf{nommer le type de plan} (déductif, inductif, dialectique) et de \textbf{citer les connecteurs} précis qui articulent les parties. Une réponse vague (« l'auteur explique puis conclut ») ne rapporte pas de points.
\end{attention}

\begin{corrige}[Exercice 3 — Paronymie : Choix et définition]
\textbf{1.} L'\textbf{imminence} de l'arrivée du ministre a provoqué une grande agitation dans le protocole.\\
Mot non utilisé : \emph{éminence} = qualité de ce qui est élevé, supérieur ; titre honorifique (Son Éminence) ; au sens propre, petite élévation de terrain. \emph{Imminence} = caractère de ce qui va arriver très bientôt.

\textbf{2.} Le témoin a dû \textbf{abjurer} dire la vérité sous serment.\\
Correction de sens : le verbe attendu est \textbf{adjurer} (= prier instamment, supplier avec insistance : « il l'adjura de dire la vérité »). La phrase correcte est donc : « Le témoin a dû être \textbf{adjuré} de dire la vérité sous serment » ou « On a dû \textbf{adjurer} le témoin de dire la vérité ».\\
Mot non utilisé : \emph{abjurer} = renier solennellement (une religion, une opinion), désavouer publiquement.

\textbf{3.} L'\textbf{infusion} des plantes médicinales dans l'eau bouillante libère les principes actifs.\\
Mot non utilisé : \emph{confusion} = trouble, mélange désordonné, erreur consistant à prendre une chose pour une autre. \emph{Infusion} = action de verser un liquide bouillant sur une substance pour en extraire les principes.
\end{corrige}

\begin{attention}[Point de vigilance]
La paronymie est un piège classique au Bac : deux mots proches par la forme mais différents par le sens. Vérifiez toujours le sens dans le contexte avant d'écrire. Entraînez-vous sur les paires fréquentes : \emph{inférer/infirmer}, \emph{préjudice/préjugé}, \emph{addition/édiction}, \emph{événement/avènement}.
\end{attention}

\begin{corrige}[Exercice 4 — Réduction directe : Phrase complexe]
Phrase de départ (42 mots). Réduction proposée (14 mots) :

\begin{center}
\textbf{« Les partenaires doivent se réunir rapidement pour évaluer les résultats obtenus. »}
\end{center}

Détail des opérations :
\begin{itemize}
    \item « Il est impératif que ... puissent » $\rightarrow$ « doivent » : [\emph{condensation} de la tournure impersonnelle en modalité simple].
    \item « l'ensemble des parties prenantes, y compris les partenaires techniques et financiers » $\rightarrow$ « les partenaires » : [\emph{généralisation} par hyperonyme + \emph{suppression} de l'incise « y compris... »].
    \item « dans les plus brefs délais » $\rightarrow$ « rapidement » : [\emph{condensation} de la locution en un adverbe].
    \item « afin de procéder à l'évaluation minutieuse » $\rightarrow$ « pour évaluer » : [\emph{condensation} par infinitif direct + \emph{suppression} de l'adjectif « minutieuse », détail secondaire].
\end{itemize}

Vérification du comptage : Les(1) partenaires(2) doivent(3) se(4) réunir(5) rapidement(6) pour(7) évaluer(8) les(9) résultats(10) obtenus(11). Soit \textbf{11 mots} $\leq$ 15. L'information essentielle (qui / quoi / quand / pourquoi) est conservée.
\end{corrige}

\begin{attention}[Point de vigilance]
Dans la contraction, chaque mot doit \textbf{porter de l'information}. Supprimez les adjectifs qualificatifs non techniques, les incises, les périphrases verbales. Comptez les mots (mots-outils inclus) et annotez le total en marge.
\end{attention}

\begin{corrige}[Exercice 5 — Repérage des registres : « Le Lion et la perle »]
\textbf{LAKUNLE} $\rightarrow$ Registre \cle{soutenu} (et didactique). Marqueurs : vocabulaire abstrait et recherché (« émancipation », « instruction », « coutumes archaïques ») ; syntaxe complète et construite avec opposition équilibrée (« par l'instruction, non par la dot »). Ce registre traduit son prétentieux vernis occidental.

\textbf{SIDI} $\rightarrow$ Registre \cle{courant} avec une pointe familière. Marqueurs : apostrophe directe et moqueuse « Toi, l'homme aux grands mots ! » ; phrases courtes, ton parlé (« Tu parles comme un livre », comparaison populaire) ; vocabulaire concret (« la femme vaut son prix »).

\textbf{BAROKA} $\rightarrow$ Registre \cle{soutenu} (solennel, traditionnel). Marqueurs : tournure impersonnelle de type administratif « Il convient d'examiner » ; vocabulaire de la sagesse et de l'autorité (« la sagesse des ancêtres », « notre identité ») ; conditionnel de prudence « ne saurait effacer ».
\end{corrige}

\begin{attention}[Point de vigilance]
Dans une pièce de théâtre, le registre de langue est un \textbf{indice de caractérisation} : il révèle le statut social, l'éducation et la stratégie du personnage. Le justifier par la situation d'énonciation (qui parle à qui, dans quel but) valorise la copie.
\end{attention}

\begin{corrige}[Exercice 6 — Champs lexicaux et hyperonymes : Agriculture au Cameroun]
\textbf{1. Classement en trois champs lexicaux :}

\begin{center}
\begin{tabularx}{\linewidth}{|l|X|}\hline
\rowcolor{popblueL} \textbf{Champ lexical} & \textbf{Mots} \\ \hline
Les cultures vivrières & bananier, plantain, maïs, mil, sorgho, igname, manioc, patate \\ \hline
Les cultures industrielles et d'exportation & cacao, café, cotonnier, palmier à huile, hévéa \\ \hline
Les techniques et l'organisation agricoles & engrais, pesticide, tracteur, charrue, irrigation, coopérative, rendement \\ \hline
\end{tabularx}
\end{center}

\textbf{2. Trois énumérations remplacées par leur hyperonyme :}
\begin{itemize}
    \item cacao, café, cotonnier $\rightarrow$ \cle{cultures d'exportation} (ou cultures de rente) ;
    \item maïs, mil, sorgho $\rightarrow$ \cle{céréales} ;
    \item tracteur, charrue $\rightarrow$ \cle{matériel agricole} (ou outillage agricole) ;
    \item (autre possibilité) igname, manioc, patate $\rightarrow$ \cle{tubercules}.
\end{itemize}

L'hyperonyme est le mot au sens plus général qui englobe les hyponymes : remplacer une énumération par son hyperonyme est une \textbf{technique de réduction} essentielle du résumé.
\end{corrige}

\begin{attention}[Point de vigilance]
L'hyperonymie/hyponymie structure la pensée technique. Au Bac, on vous demandera souvent de \textbf{remplacer une liste par son hyperonyme} pour gagner des mots. Ex. : « le maïs, le mil, le sorgho, le riz » $\rightarrow$ « les céréales » (gain de 3 mots).
\end{attention}

\begin{corrige}[Exercice 7 — Transformation de phrase complexe : Discours officiel]
Phrase de départ (46 mots). Transformation proposée (15 mots) :

\begin{center}
\textbf{« Le Ministre révisera les curricula avec les opérateurs économiques pour adapter la formation à l'emploi. »}
\end{center}

Comptage : Le(1) Ministre(2) révisera(3) les(4) curricula(5) avec(6) les(7) opérateurs(8) économiques(9) pour(10) adapter(11) la(12) formation(13) à(14) l'emploi(15). Soit \textbf{15 mots} $\leq$ 16.

Étapes justifiées :
\begin{enumerate}
    \item \textbf{Suppression de l'incident} : « de l'Éducation Technique » (précision secondaire, le contexte suffit).
    \item \textbf{Nominalisation} : « considérant que la formation professionnelle doit répondre aux besoins réels du marché de l'emploi » $\rightarrow$ « pour adapter la formation à l'emploi » (la subordonnée de cause devient un complément de but infinitif).
    \item \textbf{Fusion} : « a décidé de réviser » $\rightarrow$ « révisera » (le verbe de décision est absorbé par le futur simple).
    \item \textbf{Condensation} : « en concertation avec » $\rightarrow$ « avec ».
\end{enumerate}
L'information essentielle (qui : le Ministre ; quoi : révision des curricula ; comment : avec les opérateurs ; pourquoi : adapter la formation à l'emploi) est intégralement conservée.
\end{corrige}

\begin{attention}[Point de vigilance]
La \textbf{nominalisation} (verbe $\rightarrow$ nom, ou subordonnée $\rightarrow$ infinitif/nom) est la technique n°1 de condensation en français administratif et technique. Maîtrisez les suffixes : \emph{-tion, -ment, -age, -ance, -eur}.
\end{attention}

\begin{corrige}[Exercice 8 — Dégager et reformuler les idées essentielles]
\textbf{1. Thèse défendue :} c'est une thèse de \cle{conciliation} : l'introduction des langues nationales dans l'enseignement technique est souhaitable à condition d'adopter une approche bilingue progressive associant langue nationale et langue officielle.

\textbf{2. Structure argumentative :}

\begin{center}
\begin{tabularx}{\linewidth}{|l|X|}\hline
\rowcolor{popblueL} \textbf{Étape} & \textbf{Contenu} \\ \hline
Arguments POUR & Meilleure compréhension des concepts abstraits ; exemple des pays asiatiques industrialisés via leurs langues maternelles (argument d'exemple). \\ \hline
Arguments CONTRE & Risque de fragmentation linguistique (unité nationale menacée) ; l'anglais et le français restent indispensables à l'insertion internationale. \\ \hline
Conciliation & Bilinguisme progressif : associer langue nationale et langue officielle. \\ \hline
\end{tabularx}
\end{center}

\textbf{3. Reformulation des trois idées essentielles (sans copier le texte) :}
\begin{enumerate}
    \item Enseigner les techniques dans les langues nationales rendrait les notions difficiles plus accessibles aux apprenants, comme l'illustre la réussite industrielle de certains pays d'Asie.
    \item Cette réforme comporte toutefois des dangers : elle pourrait diviser le pays linguistiquement et priver les techniciens des langues officielles exigées par le monde professionnel.
    \item La voie la plus réaliste consiste donc à combiner progressivement la langue nationale et la langue officielle dans un enseignement bilingue.
\end{enumerate}
\end{corrige}

\begin{attention}[Point de vigilance]
Reformuler, c'est dire la même idée avec d'autres mots (synonymes, changement de tournure, nominalisation), jamais recopier. Toute phrase reprise mot pour mot est pénalisée comme un copier-coller. Vérifiez : \emph{sujet changé ? verbe changé ? structure changée ?}
\end{attention}

\begin{corrige}[Exercice 9 — Sujet type Bac Tech : Contraction de texte argumentatif]
\textbf{1. Analyse de la structure argumentative (4 points)}

a) \textbf{Thèse principale (1 pt)} : l'enseignement technique camerounais doit adopter un bilinguisme additif — la langue nationale comme langue d'explication et d'ancrage, le français ou l'anglais comme langues professionnelles — pour former des techniciens compétents, enracinés et employables.

b) \textbf{Les grandes parties (2 pts)} :
\begin{itemize}
    \item \emph{Constat} (paragraphes 1-2) : les élèves maîtrisent mal les langues officielles, ce qui freine l'apprentissage technique ; fonction : poser le problème.
    \item \emph{Arguments pour} (paragraphe 3) : les expériences régionales et la recommandation de l'UNESCO (argument d'autorité) ; fonction : étayer la thèse.
    \item \emph{Obstacles} (paragraphe 4) : standardisation inachevée, manque de manuels et de formateurs, exigences du marché de l'emploi ; fonction : concession, examen honnête des difficultés.
    \item \emph{Solution} (paragraphe 5) : le bilinguisme additif et ses conditions politiques ; fonction : conclure par une proposition réaliste.
\end{itemize}

c) \textbf{Connecteurs (1 pt)} : opposition : « \emph{Pourtant} », « \emph{Toutefois} » ; conséquence : « \emph{C'est à ce prix que} », « \emph{ce qui} » (freine).

\textbf{2. Les 7 idées forces (6 points, environ 0,85 pt par idée)}
\begin{enumerate}
    \item Le système éducatif camerounais repose sur le français et l'anglais, alors que plus de 250 langues nationales portent les cultures locales.
    \item La mauvaise maîtrise des langues officielles empêche les élèves d'assimiler les notions scientifiques et techniques.
    \item Des expériences régionales prouvent qu'expliquer les techniques dans la langue maternelle facilite et accélère la compréhension.
    \item L'UNESCO recommande l'enseignement en langue maternelle au moins jusqu'au secondaire.
    \item Des obstacles subsistent : orthographes non standardisées, pénurie de manuels traduits et de formateurs bilingues.
    \item Le marché de l'emploi exige les langues officielles ; un enseignement uniquement en langue nationale limiterait l'insertion.
    \item La solution est un bilinguisme additif, exigeant terminologies normalisées, formation des enseignants et supports adaptés.
\end{enumerate}

\textbf{3. Résumé rédigé (10 points) — proposition (environ 150 mots) :}

\begin{quote}
Au Cameroun, l'enseignement officiel en français et en anglais coexiste avec plus de 250 langues nationales. Or la faible maîtrise des langues officielles freine l'assimilation des concepts techniques. Des expériences menées dans l'Ouest démontrent que l'explication des savoirs dans la langue maternelle améliore la compréhension, ce que confirment les recommandations de l'UNESCO. Cependant, des obstacles demeurent : standardisation inachevée, manque de manuels et de formateurs bilingues, exigence des langues officielles sur le marché de l'emploi. La solution réside dans un bilinguisme additif : la langue nationale sert d'ancrage cognitif tandis que le français ou l'anglais restent les langues professionnelles. Cette approche requiert une volonté politique : terminologies normalisées, formation des enseignants et supports adaptés, afin de former des techniciens compétents, enracinés et ouverts au monde.
\end{quote}

Comptage indicatif : 149 mots $\rightarrow$ dans la fourchette 135-165.

\textbf{Barème indicatif (rédaction, 10 pts)} : respect de la longueur (1 pt) ; fidélité aux idées essentielles (3 pts) ; reformulation sans copier-coller (2 pts) ; registre soutenu et impersonnalité (2 pts) ; cohérence et connecteurs (1 pt) ; correction linguistique (1 pt).
\end{corrige}

\begin{attention}[Points de vigilance Bac]
\begin{itemize}
    \item Toute phrase copiée du texte = point perdu en reformulation.
    \item Aucun « je », aucun avis personnel : le résumé est \textbf{impersonnel}.
    \item Annoncer le nombre de mots en fin de copie et le respecter à $\pm$ 10\%.
    \item Conserver l'ordre logique du texte (constat $\rightarrow$ pour $\rightarrow$ obstacles $\rightarrow$ solution).
\end{itemize}
\end{attention}

\begin{corrige}[Exercice 10 — Compte-rendu professionnel : Réunion de coopérative]
\textbf{Proposition de compte-rendu (environ 100 mots) :}

\begin{quote}
\textbf{Compte-rendu de la réunion de la Coopérative « Espoir Paysan »} — Bafoussam II, 15 octobre 2024, en présence de 42 adhérents sur 60. Le bilan de la campagne café 2023-2024 (production : 120 tonnes, en baisse de 15\% ; bénéfice net : 27 millions FCFA) a été approuvé à l'unanimité. Pour la campagne cacao 2024-2025 (prévision : 200 tonnes), le recours à un crédit agricole CAMCCUL de 5\% sur 3 ans a été adopté (38 voix pour, 2 contre, 2 abstentions). L'acquisition d'un camion (18 millions FCFA) est reportée à l'assemblée générale extraordinaire de décembre. Prochaine réunion : 15 janvier 2025.
\end{quote}

Comptage indicatif : 98 mots $\rightarrow$ dans la fourchette 90-110.

\textbf{Techniques appliquées :}
\begin{itemize}
    \item \emph{Style administratif} : voix passive et noms d'action (« a été approuvé », « a été adopté », « l'acquisition... est reportée ») ; impersonnalité totale.
    \item \emph{Sélection de l'essentiel} : chiffres clés (42/60 ; 120 t ; 27 millions ; 200 t ; 18 millions), votes et échéances conservés ; détails secondaires supprimés (noms des absents, devis détaillé, divers).
    \item \emph{Structure standard} : objet, date, participants, décisions, votes, échéances.
\end{itemize}
\end{corrige}

\begin{attention}[Point de vigilance]
Un compte-rendu n'est pas un procès-verbal intégral : il ne rapporte que les \textbf{décisions, chiffres et échéances}. Toute appréciation personnelle ou tout détail anecdotique est hors sujet. Respectez le format : en-tête (objet, date, lieu, présents), corps (décisions), pied (prochaine échéance, signature).
\end{attention}

\begin{corrige}[Exercice 11 — Remédiation : Correction d'un résumé d'élève]
\textbf{1. Identification des cinq erreurs :}

\begin{center}
\begin{tabularx}{\linewidth}{|l|X|X|}\hline
\rowcolor{popblueL} \textbf{N°} & \textbf{Passage fautif} & \textbf{Type d'erreur et gravité} \\ \hline
1 & « les élèves maîtrisent mal les langues officielles » (entre guillemets) & \emph{Copie littérale} : citation du texte, interdite dans un résumé ; pénalise la reformulation. \\ \hline
2 & « Moi je pense qu'il a raison car chez moi on parle bassa » & \emph{Avis personnel / subjectivité} : le résumé doit être impersonnel ; faute grave de méthode. \\ \hline
3 & 165 mots au lieu de 150 $\pm$ 10\% (max 165 : limite atteinte, mais le texte fourni en contient environ 165 avec des redondances inutiles) & \emph{Dépassement / non-concision} : le respect du nombre de mots est une consigne contractuelle. \\ \hline
4 & « ça freine », « ça marche mieux », « pas de profs », « C'est comme ça qu'on aura » & \emph{Registre familier/oral} : « ça », « profs », tournures parlées ; le résumé exige le registre soutenu. \\ \hline
5 & « solution miracle » & \emph{Contresens} : le texte dit « piste de conciliation réaliste » ; déformer la pensée de l'auteur est l'erreur la plus pénalisée. \\ \hline
\end{tabularx}
\end{center}

\textbf{2. Version corrigée du résumé (environ 150 mots) :}

\begin{quote}
L'enseignement technique au Cameroun, dispensé en français et en anglais, se heurte à la faible maîtrise de ces langues par les élèves, ce qui nuit à l'assimilation des notions scientifiques. Des expériences conduites dans l'Ouest montrent que l'explication des savoirs techniques en langue maternelle améliore sensiblement la compréhension, conformément aux recommandations de l'UNESCO. Néanmoins, des difficultés persistent : standardisation orthographique inachevée, insuffisance de manuels traduits et de formateurs bilingues, exigence des langues officielles sur le marché de l'emploi. La voie préconisée est un bilinguisme additif : la langue nationale assure l'ancrage cognitif, tandis que le français ou l'anglais demeurent les langues professionnelles. Cette réforme suppose terminologies normalisées, formation des enseignants et supports adaptés, pour former des techniciens compétents et ouverts au monde.
\end{quote}

Comptage indicatif : 138 mots $\rightarrow$ dans la fourchette 135-165.

\textbf{Corrections appliquées :} citation supprimée et reformulée ; avis personnel éliminé ; « ça » $\rightarrow$ « ce qui », « profs » $\rightarrow$ « enseignants/formateurs » ; « solution miracle » $\rightarrow$ « voie préconisée » (fidélité au sens) ; connecteurs soutenus (« Néanmoins », « tandis que ») ; impersonnalité totale.
\end{corrige}

\begin{attention}[Point de vigilance]
Au Bac, la grille de correction sanctionne en priorité : le \textbf{contresens}, le \textbf{copier-coller} et l'\textbf{avis personnel}. Avant de rendre la copie, relisez en posant trois questions : ai-je copié ? ai-je dit « je » ? ai-je respecté le sens de l'auteur et le nombre de mots ?
\end{attention}

\newpage

\coursec{Fiche bilan}

\courssub{Carte mentale de la leçon}
\begin{popfigure}
\begin{tikzpicture}[
    mindmap,
    grow cyclic,
    every node/.style={concept, execute at begin node=\hskip0pt},
    concept color=popblue,
    level 1/.append style={level distance=4.5cm, sibling angle=360/6},
    level 2/.append style={level distance=3cm, sibling angle=360/8},
    text=white,
    font=\small\bfseries
]
\node[concept, scale=1.2, align=center]{Résumé \&\\ Reformulation}
    child[concept color=poporange] { node[concept, align=center]{Registres\\ de langue} }
    child[concept color=popgreen] { node[concept, align=center]{Structure\\ argumentative} }
    child[concept color=poppurple] { node[concept, align=center]{Champ lexical\\ \& sémantique} }
    child[concept color=poppink] { node[concept, align=center]{Phrase\\ complexe} }
    child[concept color=popgold] { node[concept, align=center]{Techniques\\ de réduction} }
    child[concept color=popblue] { node[concept, align=center]{Rédaction\\ du résumé} };
\end{tikzpicture}
\legende{Carte mentale : les 6 piliers de la contraction de texte en Première Technique.}
\end{popfigure}

\courssub{Bilan synthétique}

\begin{bilan}

\courssub*{1. Registres de langue et communication}
\begin{itemize}
    \item \cle{Registre soutenu} : vocabulaire précis, syntaxe complexe, pas d'argot (ex. : \emph{« Il est impératif de... »}).
    \item \cle{Registre courant} : langue standard, communication quotidienne (ex. : \emph{« Il faut... »}).
    \item \cle{Registre familier} : vocabulaire expressif, syntaxe relâchée, argot (ex. : \emph{« Faut que... »}).
    \item \textbf{Adaptation} : le résumé s'écrit en \textbf{registre soutenu} (objectivité, concision).
\end{itemize}

\courssub*{2. La structure argumentative du texte}
\begin{center}
\begin{tabularx}{\linewidth}{|l|X|}\hline
\rowcolor{popblueL}\textbf{Étape} & \textbf{Description} \\ \hline
\cle{Thèse} & Idée principale défendue par l'auteur. \\ \hline
\cle{Arguments} & Raisons logiques, faits, exemples, citations. \\ \hline
\cle{Connecteurs} & \emph{En effet, par conséquent, toutefois, de plus, ainsi}. \\ \hline
\cle{Plan} & Dialectique (thèse/antithèse/synthèse), analytique (causes/conséquences), thématique. \\ \hline
\end{tabularx}
\end{center}

\courssub*{3. Champ lexical, champ sémantique et relations de sens}
\begin{definition}[Champ lexical]
Ensemble des mots (formes) appartenant à une même catégorie grammaticale et liés à un même thème (ex. : \emph{mer, vague, sable, marée} = champ lexical de la mer).
\end{definition}
\begin{definition}[Champ sémantique]
Ensemble des \emph{sens} (concepts) liés à une notion, toutes catégories confondues (ex. : \emph{naviguer, bleu, sel, naufrage}).
\end{definition}
\begin{itemize}
    \item \cle{Hyperonymie} : relation genre/espèce. L'hyperonyme (genre) \rightarrow \emph{animal}. L'hyponyme (espèce) \rightarrow \emph{chien, chat}.
    \item \cle{Paronymie} : mots proches par la forme, différents par le sens (ex. : \emph{inférer / inférer} -- attention aux confusions Bac).
\end{itemize}
\textbf{Utilité pour le résumé} : remplacer des énumérations par l'hyperonyme (réduction).

\courssub*{4. La phrase complexe au service de la reformulation}
\begin{propriete}[Types de subordination]
\begin{itemize}
    \item \cle{Relative} : antecedent + pronom relatif (\emph{qui, que, dont, où}) $\rightarrow$ intègre une définition.
    \item \cle{Conjonctive (complétive)} : \emph{que, si} $\rightarrow$ rapporte un propos ou un fait.
    \item \cle{Circonstancielle} : \emph{parce que, bien que, quand, si, afin que} $\rightarrow$ exprime cause, concession, temps, condition, but.
\end{itemize}
\end{propriete}
\begin{methode}[Transformer pour condenser]
1. Repérer les phrases simples juxtaposées ou coordonnées.\\
2. Subordonner l'accessoire au principal (cause $\rightarrow$ \emph{parce que} ; concession $\rightarrow$ \emph{bien que}).\\
3. Nominaliser (verbe $\rightarrow$ nom : \emph{il argue} $\rightarrow$ \emph{son argumentation}).\\
4. Utiliser les participes présents/gérondifs pour l'antériorité/simultanéité (\emph{En argumentant, il convainc}).
\end{methode}

\courssub*{5. Les techniques de réduction (contraction)}
\begin{center}
\begin{tabularx}{\linewidth}{|l|X|l|}\hline
\rowcolor{popblueL}\textbf{Technique} & \textbf{Explication} & \textbf{Exemple} \\ \hline
\cle{Suppression} & Enlever exemples, détails, répétitions, adjectifs accessoires. & \emph{« Le grand et beau lion »} $\rightarrow$ \emph{« Le lion »}. \\ \hline
\cle{Généralisation} & Remplacer liste/hyponymes par hyperonyme. & \emph{« Chiens, chats, oiseaux »} $\rightarrow$ \emph{« Animaux »}. \\ \hline
\cle{Nominalisation} & Transformer verbe/adjectif en nom. & \emph{« Il a décidé »} $\rightarrow$ \emph{« Sa décision »}. \\ \hline
\cle{Subordination} & Fondre plusieurs phrases en une phrase complexe. & \emph{« Il pleut. Je sors. »} $\rightarrow$ \emph{« Bien qu'il pleuve, je sors. »} \\ \hline
\cle{Pronominalisation} & Remplacer groupes nominaux répétés par pronoms. & \emph{« Sidi... Sidi... »} $\rightarrow$ \emph{« Il... »}. \\ \hline
\cle{Synthèse} & Regrouper plusieurs arguments en une idée force. & 3 arguments $\rightarrow$ 1 phrase résumée. \\ \hline
\end{tabularx}
\end{center}

\courssub*{6. Reformuler les idées essentielles et rédiger le résumé}
\begin{methode}[Procédé en 5 étapes (Méthode Bac)]
1. \textbf{Lecture active} : identifier thèse, arguments, structure, ton.\\
2. \textbf{Découpage} : segmenter en macro-propositions (1 paragraphe = 1 idée force).\\
3. \textbf{Prise de notes} : mots-clés, hyperonymes, connecteurs logiques (pas de phrases complètes).\\
4. \textbf{Rédaction du brouillon} : reformuler \emph{avec ses mots} (registre soutenu), respecter l'ordre logique, appliquer techniques de réduction. \textbf{Ne pas commenter, ne pas citer}.\\
5. \textbf{Contrôle final} : nombre de mots ($\pm 10\%$), cohérence, syntaxe (phrases complexes), orthographe, fidélité au sens original.
\end{methode}

\begin{attention}[Pièges fréquents au Probatoire]
\begin{enumerate}
\item Dépasser la longueur imposée (souvent 1/4 ou 1/3 du texte original).
\item Copier-coller des phrases du texte (pas de reformulation = note basse).
\item Ajouter son opinion, des exemples personnels ou des connaissances extérieures.
\item Oublier un argument majeur (perte de sens).
\item Mal gérer les connecteurs logiques (incohérence).
\item Confondre paronymes (\emph{inférer / inférer, attention / intention}).
\end{enumerate}
\end{attention}

\end{bilan}

\begin{aretenir}[Checklist avant le Bac Probatoire]
$\square$ Je distingue les 3 registres de langue et j'écris en registre soutenu.\\
$\square$ J'identifie la thèse, les arguments et les connecteurs d'un texte argumentatif.\\
$\square$ J'utilise l'hyperonymie pour généraliser (réduire les listes).\\
$\square$ Je maîtrise la paronymie (pièges lexicaux courants).\\
$\square$ Je construis des phrases complexes (relatives, conjonctives, circonstancielles).\\
$\square$ J'applique les 5 techniques de réduction : suppression, généralisation, nominalisation, subordination, pronominalisation.\\
$\square$ Je respecte la procédure en 5 étapes (lecture $\rightarrow$ notes $\rightarrow$ brouillon $\rightarrow$ rédaction $\rightarrow$ contrôle).\\
$\square$ Je respecte la contrainte de longueur (nombre de mots imposé $\pm 10\%$).\\
$\square$ Je ne cite pas, je ne commente pas, je ne donne pas mon avis.\\
$\square$ Je vérifie la cohérence logique (connecteurs) et la correction syntaxique/orthographique.\\
$\square$ Je gère mon temps : lecture (15\%), notes (20\%), rédaction (50\%), relecture (15\%).\\
\end{aretenir}

\findecours

\end{document}
